% Les Entreprises — Philosophie 1re Tech — Cours signé OnBuch+
% Programme officiel MINESEC. Compiler avec : tectonic N12-entreprises.tex
\newcommand{\DOCMATIERE}{Philosophie}
\newcommand{\DOCNIVEAU}{1re Tech}
\newcommand{\DOCMODULE}{Philosophie – classe de Première technique}
\newcommand{\DOCLECON}{N12}
\newcommand{\DOCTITRE}{Les Entreprises}
% OnBuch+ — préambule « cours signés » ENRICHI (compilé avec tectonic / XeLaTeX)
% Système visuel « Pop » d'OnBuch+ : fond crème (jamais blanc), bordures encre
% épaisses, ombres « dures » décalées, titres Archivo Black en capitales.
% Version MATHÉMATIQUES : graphes (pgfplots/tikz), boîtes pédagogiques
% (définition, propriété, méthode, attention, le savais-tu, exercice résolu,
% exercice, TP, bilan), page de garde et signature OnBuch+ ancrée en bas.
%
% Macros attendues dans main.tex AVANT \input{preamble} :
%   \DOCMATIERE  \DOCNIVEAU  \DOCMODULE  \DOCLECON (numéro)  \DOCTITRE
\documentclass[11pt]{article}

\usepackage{fontspec}
\usepackage[french]{babel}
\usepackage[a4paper,margin=2cm,top=3.4cm,bottom=2.8cm,headheight=1.4cm,footskip=1.3cm]{geometry}
\usepackage{amsmath,amssymb,amsfonts}
\usepackage{mathtools} % \xmapsto, \coloneqq... souvent utilisés par les modèles
\usepackage{cancel} % simplifications barrées (bilans de forces)
% \abs{z} (module, valeur absolue) et \norm{u} (norme), utilisés par les modèles
\providecommand{\abs}{}\let\abs\relax\DeclarePairedDelimiter{\abs}{\lvert}{\rvert}
\providecommand{\norm}{}\let\norm\relax\DeclarePairedDelimiter{\norm}{\lVert}{\rVert}
\usepackage[table]{xcolor} % [table] : fournit aussi \rowcolors (alternance de lignes)
\usepackage{pagecolor}
\usepackage{tikz}
\usetikzlibrary{arrows.meta,positioning,calc,shapes.geometric,shapes.misc,shapes.arrows,decorations.pathmorphing,decorations.pathreplacing,decorations.markings,patterns,shadows,mindmap,trees,fit,backgrounds,matrix,intersections,angles,quotes}
\usepackage{pgfplots}
\usepackage[version=4]{mhchem} % équations nucléaires \ce{^{235}_{92}U}
\usepackage[european,siunitx]{circuitikz} % schémas électriques
\pgfplotsset{compat=1.18}
\usepgfplotslibrary{fillbetween,groupplots} % aires entre courbes (calcul intégral)
\usepackage{enumitem}
\usepackage{fancyhdr}
% [most] charge toutes les bibliothèques tcolorbox (listings, minted, raster,
% documentation...) et gonfle inutilement la mémoire TeX sur un document long
% (~300 boîtes) : on ne charge que ce qui est réellement utilisé.
\usepackage[skins,breakable]{tcolorbox}
\usepackage{graphicx}
\usepackage{colortbl}
\usepackage{booktabs}
\usepackage{tabularx}
\usepackage{multirow}
\usepackage{diagbox} % case d'en-tête diagonale des tableaux à double entrée
% Colonnes extensibles centrée / gauche / droite (C, L, R), souvent utilisées par les modèles
\newcolumntype{C}{>{\centering\arraybackslash}X}
\newcolumntype{L}{>{\raggedright\arraybackslash}X}
\newcolumntype{R}{>{\raggedleft\arraybackslash}X}
\usepackage{array}
\usepackage{multicol}
\usepackage{siunitx}
% parse-numbers=false : accepte aussi \SI{2,5 \times 10^{-3}}{...}
\sisetup{locale=FR,output-decimal-marker={,},group-minimum-digits=4,parse-numbers=false}
\DeclareSIUnit\ohm{\ensuremath{\symup{\Omega}}} % U+2126 absent de la police
\DeclareSIUnit\division{div} % réglages d'oscilloscope (ms/div, V/div)
\DeclareSIUnit\tr{tr} % tours (vitesse de rotation, ex. \SI{1500}{\tr\per\minute})
\DeclareSIUnit\molar{M} % concentration molaire (ex. \SI{3}{\molar}, \SI{1}{\micro\molar})
\DeclareSIUnit\an{an} % année (le modèle confond parfois avec \year, un compteur TeX) — ex. \SI{5}{\kilo\meter\per\an}
\DeclareSIUnit\FCFA{FCFA} % franc CFA (ex. \SI{500}{\FCFA\per\kilo\gram})
\DeclareSIUnit\degreeBrix{\textdegree Brix} % teneur en sucre d'un sirop/jus (réfractométrie)
\DeclareSIUnit\month{mois} % le modèle utilise parfois \month, un compteur TeX, comme unité de durée
\DeclareSIUnit\microliter{\micro\liter} % le modèle écrit parfois \microliter en un seul token
\DeclareSIUnit\million{millions} % dénombrement (ex. \SI{15}{\million\per\mL} spermatozoïdes)
\usepackage{qrcode}
% \degree : commande fréquemment utilisée par les modèles pour un angle ou une
% température en dehors de \SI{}{} (non fournie par les paquets chargés).
\providecommand{\degree}{^{\circ}}
% \qed (fin de démonstration), utilisé par les modèles sans amsthm
\providecommand{\qed}{\hfill$\square$}
% \barre : double barre des valeurs interdites dans un tableau de variations (tkz-tab)
\providecommand{\barre}{\ensuremath{\|}}
% environnement proof (amsthm non chargé) utilisé par les modèles
\ifdefined\proof\else\newenvironment{proof}[1][Démonstration]{\par\noindent\textit{#1.}\ }{\qed\par}\fi
\usepackage{lastpage}
\usepackage{hyperref}
\hypersetup{hidelinks}

% ============================================================
% Palette « Pop » OnBuch+ — mêmes hex que la classe P (plus_ui.dart)
% ============================================================
\definecolor{poporange}{HTML}{F59321}
\definecolor{popdark}{HTML}{E07A0C}
\definecolor{popcream}{HTML}{FFF6E8}
\definecolor{popink}{HTML}{1C1714}
\definecolor{poppurple}{HTML}{7A5AE0}
\definecolor{popgreen}{HTML}{2E9E6A}
\definecolor{poppink}{HTML}{E0457B}
\definecolor{popblue}{HTML}{2F7DE1}
\definecolor{popgold}{HTML}{C98A12}
\definecolor{popmuted}{HTML}{8A7F76}
\definecolor{popline}{HTML}{EADFD2}
\definecolor{popgray}{HTML}{8A7F76}
\colorlet{popgrey}{popgray}
% Alias tolérants (le modèle nomme parfois les couleurs différemment)
\colorlet{popwhite}{white}
\colorlet{popblack}{popink}
\colorlet{popred}{poppink}
\colorlet{popyellow}{popgold}
% Teintes claires (fonds de tableaux/encadrés) — le modèle nomme parfois
% les couleurs avec un suffixe L (ex. popblueL) sans les avoir définies.
\colorlet{poporangeL}{poporange!20}
\colorlet{popdarkL}{popdark!20}
\colorlet{poppurpleL}{poppurple!20}
\colorlet{popgreenL}{popgreen!20}
\colorlet{poppinkL}{poppink!20}
\colorlet{popblueL}{popblue!20}
\colorlet{popgoldL}{popgold!20}
\colorlet{popredL}{popred!20}
\colorlet{popgrayL}{popgray!20}
% Teintes claires pour fonds de tableaux / zones
\colorlet{popblueL}{popblue!10!white}
\colorlet{poporangeL}{poporange!14!white}
\colorlet{popgreenL}{popgreen!12!white}
\colorlet{poppurpleL}{poppurple!12!white}
\colorlet{poppinkL}{poppink!12!white}
\colorlet{popgoldL}{popgold!14!white}

\pagecolor{popcream}

% ============================================================
% Typographie
% ============================================================
\defaultfontfeatures{Ligatures=TeX}
\setmainfont{PlusJakartaSans-Medium.ttf}[Path=../../../fonts/,BoldFont=PlusJakartaSans-Bold.ttf]
% Maths en sans-serif (Fira Math) avec chiffres et lettres droites de Plus
% Jakarta, pour que les formules se fondent dans le texte.
\usepackage[math-style=ISO,bold-style=ISO]{unicode-math}
\setmathfont{FiraMath-Regular.otf}[Path=../../../fonts/]
\setmathfont{PlusJakartaSans-Medium.ttf}[Path=../../../fonts/,range={up/num,up/{Latin,latin},\mathup}]
\usepackage{icomma} % virgule décimale sans espace en mode math
% Caractères mathématiques Unicode absents des polices de texte (Plus Jakarta,
% Archivo Black) : rendus par leur équivalent mathématique, en texte comme en
% formule — sinon ils s'affichent en carré vide (ex. « Arithmétique dans ℤ »).
\usepackage{newunicodechar}
\newunicodechar{ℝ}{\ensuremath{\mathbb{R}}}
\newunicodechar{ℤ}{\ensuremath{\mathbb{Z}}}
\newunicodechar{ℂ}{\ensuremath{\mathbb{C}}}
\newunicodechar{ℕ}{\ensuremath{\mathbb{N}}}
\newunicodechar{ℚ}{\ensuremath{\mathbb{Q}}}
\newunicodechar{𝔻}{\ensuremath{\mathbb{D}}}
\newunicodechar{α}{\ensuremath{\alpha}}
\newunicodechar{β}{\ensuremath{\beta}}
\newunicodechar{γ}{\ensuremath{\gamma}}
\newunicodechar{δ}{\ensuremath{\delta}}
\newunicodechar{ε}{\ensuremath{\varepsilon}}
\newunicodechar{θ}{\ensuremath{\theta}}
\newunicodechar{λ}{\ensuremath{\lambda}}
\newunicodechar{μ}{\ensuremath{\mu}}
\newunicodechar{σ}{\ensuremath{\sigma}}
\newunicodechar{τ}{\ensuremath{\tau}}
\newunicodechar{φ}{\ensuremath{\varphi}}
\newunicodechar{ψ}{\ensuremath{\psi}}
\newunicodechar{ω}{\ensuremath{\omega}}
\newunicodechar{Σ}{\ensuremath{\Sigma}}
% majuscules produites par \MakeUppercase dans les titres (α-aminés -> Α)
\newunicodechar{Α}{\ensuremath{\alpha}}
\newunicodechar{Β}{\ensuremath{\beta}}
\newunicodechar{Γ}{\ensuremath{\Gamma}}
\newunicodechar{Δ}{\ensuremath{\Delta}}
\newunicodechar{Ε}{\ensuremath{\varepsilon}}
\newunicodechar{Θ}{\ensuremath{\Theta}}
\newunicodechar{Λ}{\ensuremath{\Lambda}}
\newunicodechar{Μ}{\ensuremath{\mu}}
\newunicodechar{Τ}{\ensuremath{\tau}}
\newunicodechar{Φ}{\ensuremath{\Phi}}
\newunicodechar{Ψ}{\ensuremath{\Psi}}
\newunicodechar{Ω}{\ensuremath{\Omega}}
\newunicodechar{↦}{\ensuremath{\mapsto}}
\newunicodechar{∈}{\ensuremath{\in}}
\newunicodechar{∉}{\ensuremath{\notin}}
\newunicodechar{∧}{\ensuremath{\wedge}}
\newunicodechar{⇔}{\ensuremath{\Leftrightarrow}}
\newunicodechar{⟺}{\ensuremath{\Longleftrightarrow}}
\newunicodechar{⇒}{\ensuremath{\Rightarrow}}
\newunicodechar{⟹}{\ensuremath{\Longrightarrow}}
\newunicodechar{∘}{\ensuremath{\circ}}
\newunicodechar{∪}{\ensuremath{\cup}}
\newunicodechar{∩}{\ensuremath{\cap}}
\newunicodechar{≡}{\ensuremath{\equiv}}
\newunicodechar{⊂}{\ensuremath{\subset}}
\newunicodechar{⊥}{\ensuremath{\perp}}
\newunicodechar{∃}{\ensuremath{\exists}}
\newunicodechar{∀}{\ensuremath{\forall}}
\newunicodechar{∛}{\ensuremath{\sqrt[3]{\,}}}
\newunicodechar{∜}{\ensuremath{\sqrt[4]{\,}}}
\newunicodechar{⃗}{\ensuremath{{}^{\to}}}
\newunicodechar{ⁿ}{\ifmmode^{n}\else\textsuperscript{\ensuremath{n}}\fi}
\newunicodechar{ˣ}{\ifmmode^{x}\else\textsuperscript{\ensuremath{x}}\fi}
\newunicodechar{ᵇ}{\ifmmode^{b}\else\textsuperscript{\ensuremath{b}}\fi}
\newunicodechar{ʳ}{\ifmmode^{r}\else\textsuperscript{\ensuremath{r}}\fi}
\newunicodechar{ᵘ}{\ifmmode^{u}\else\textsuperscript{\ensuremath{u}}\fi}
\newunicodechar{ʸ}{\ifmmode^{y}\else\textsuperscript{\ensuremath{y}}\fi}
\newunicodechar{ᵃ}{\ifmmode^{a}\else\textsuperscript{\ensuremath{a}}\fi}
\newunicodechar{ᵉ}{\ifmmode^{e}\else\textsuperscript{\ensuremath{e}}\fi}
\newunicodechar{ᵏ}{\ifmmode^{k}\else\textsuperscript{\ensuremath{k}}\fi}
\newunicodechar{ᵖ}{\ifmmode^{p}\else\textsuperscript{\ensuremath{p}}\fi}
\newunicodechar{ᴺ}{\ifmmode^{N}\else\textsuperscript{\ensuremath{N}}\fi}
\newunicodechar{ᶜ}{\ifmmode^{c}\else\textsuperscript{\ensuremath{c}}\fi}
\newunicodechar{ʰ}{\ifmmode^{h}\else\textsuperscript{\ensuremath{h}}\fi}
\newunicodechar{ᵠ}{\ifmmode^{q}\else\textsuperscript{\ensuremath{q}}\fi}
\newunicodechar{ₙ}{\ifmmode_{n}\else\textsubscript{\ensuremath{n}}\fi}
\newunicodechar{ₐ}{\ifmmode_{a}\else\textsubscript{\ensuremath{a}}\fi}
\newunicodechar{ᵢ}{\ifmmode_{i}\else\textsubscript{\ensuremath{i}}\fi}
\newunicodechar{ᵣ}{\ifmmode_{r}\else\textsubscript{\ensuremath{r}}\fi}
\newunicodechar{ₖ}{\ifmmode_{k}\else\textsubscript{\ensuremath{k}}\fi}
\newunicodechar{ᵦ}{\ifmmode_{\beta}\else\textsubscript{\ensuremath{\beta}}\fi}
\newunicodechar{ₓ}{\ifmmode_{x}\else\textsubscript{\ensuremath{x}}\fi}
\newfontfamily\popdisplay{ArchivoBlack-Regular.ttf}[Path=../../../fonts/]
\newfontfamily\poplogo{SpaceGrotesk-Bold.ttf}[Path=../../../fonts/]
\newfontfamily\popsemi{PlusJakartaSans-SemiBold.ttf}[Path=../../../fonts/]
\newfontfamily\popxbold{PlusJakartaSans-ExtraBold.ttf}[Path=../../../fonts/]
\newfontfamily\popxboldls{PlusJakartaSans-ExtraBold.ttf}[Path=../../../fonts/,LetterSpace=3.0]

\setlist[enumerate]{leftmargin=1.7em,itemsep=5pt,topsep=4pt}
\setlist[itemize]{leftmargin=1.7em,itemsep=4pt,topsep=4pt,label={\color{poporange}\textbullet}}
\setlength{\parindent}{0pt}
\setlength{\parskip}{5pt}
\renewcommand{\arraystretch}{1.25}

% pgfplots : style Pop par défaut
\pgfplotsset{
  every axis/.append style={axis line style={popink,line width=1pt},tick style={popink},
    grid style={popline,line width=0.5pt},label style={font=\small},tick label style={font=\footnotesize}},
  popcurve/.style={poporange,line width=1.8pt,no markers,smooth},
}
% Même style disponible dans un \draw TikZ simple (hors pgfplots)
\tikzset{popcurve/.style={poporange,line width=1.8pt}}
\tikzset{popgrid/.style={popline,very thin}}
\colorlet{popcurve}{poporange} % le nom du style est parfois utilisé comme couleur

% ============================================================
% Pill / squiggle
% ============================================================
\newcommand{\poppill}[3]{%
  \tikz[baseline=-0.55ex]\node[draw=popink,line width=1.2pt,rounded corners=2.6pt,%
    fill=#2,inner xsep=6.5pt,inner ysep=3pt]{%
    \popxboldls\fontsize{7}{7}\selectfont\color{#3}\MakeUppercase{#1}};%
}
\newcommand{\popsquiggle}[1]{%
  \tikz[baseline=-0.4ex]{
    \draw[#1,line width=2.6pt,line cap=round]
      plot [smooth] coordinates {(0,0.09) (0.34,0.01) (0.68,0.21) (1.02,0.01) (1.36,0.10)};
  }%
}
% Mot-clé surligné (marqueur orange sous le texte)
\newcommand{\cle}[1]{{\popxbold\color{popdark}#1}}

% ============================================================
% En-tête / pied
% ============================================================
\pagestyle{fancy}
\fancyhf{}
\renewcommand{\headrulewidth}{0pt}
\renewcommand{\footrulewidth}{0pt}
\lhead{\raisebox{-0.15ex}{\poplogo\fontsize{13}{13}\selectfont\color{popink}OnBuch\color{poporange}+}}
\rhead{\poppill{\DOCMATIERE\ \textbullet\ \DOCNIVEAU\ \textbullet\ Leçon \DOCLECON}{white}{popink}}
\lfoot{\popxbold\fontsize{7.6}{7.6}\selectfont\color{popmuted}\MakeUppercase{Cours signé OnBuch+}\ \textbullet\ {\popsemi\DOCTITRE}}
\rfoot{\tikz[baseline=-0.6ex]\node[rounded corners=8.5pt,fill=popink,minimum height=17pt,minimum width=17pt,inner xsep=5pt,inner sep=0]{\popxbold\fontsize{8}{8}\selectfont\color{popcream}\thepage\,/\,\pageref*{LastPage}};}

% ============================================================
% Titres
% ============================================================
\newcommand{\chaptitre}[2]{%
  \begin{tcolorbox}[enhanced,colback=poporange,colframe=popink,boxrule=2.6pt,arc=11pt,
      drop shadow=popink,left=15pt,right=15pt,top=12pt,bottom=11pt,before skip=3pt,after skip=18pt]
    \poppill{#2}{popink}{popcream}\par\vspace{9pt}
    {\raggedright\hyphenpenalty=10000\exhyphenpenalty=10000\popdisplay\fontsize{22}{24}\selectfont\color{popcream}\MakeUppercase{#1}\par}
  \end{tcolorbox}
}
% Section numérotée : pastille orange avec le numéro + titre Archivo
\newcounter{popsec}
\newcommand{\coursec}[1]{%
  \stepcounter{popsec}\setcounter{popsub}{0}%
  \par\vspace{16pt}\noindent
  \begin{minipage}{\linewidth}
  \tikz[baseline=-0.7ex]\node[draw=popink,line width=1.6pt,fill=poporange,rounded corners=4pt,minimum size=22pt,inner sep=0]{\popdisplay\fontsize{12}{12}\selectfont\color{popcream}\thepopsec};%
  \hspace{8pt}{\popdisplay\fontsize{15}{17}\selectfont\color{popink}\MakeUppercase{#1}}\par
  \vspace{4pt}\noindent{\color{poporange}\rule{36pt}{3.6pt}}
  \end{minipage}\par\nopagebreak\vspace{8pt}
}
\newcounter{popsub}
\newcommand{\courssub}[1]{%
  \stepcounter{popsub}%
  \par\vspace{9pt}\noindent{\popxbold\fontsize{11.5}{13}\selectfont\color{popink}{\color{poporange}\thepopsec.\thepopsub}\hspace{6pt}#1}\par\nopagebreak\vspace{4pt}
}

% ============================================================
% Boîtes pédagogiques — même recette : fond blanc, bordure épaisse colorée,
% ombre dure encre, badge plein. Titre optionnel [#1].
% ============================================================
\tcbset{popbase/.style={breakable,enhanced,colback=white,boxrule=2.3pt,arc=9pt,
  shadow={3pt}{-3pt}{0pt}{popink},left=14pt,right=14pt,top=11pt,bottom=11pt,before skip=11pt,after skip=15pt,
  fontupper=\popsemi\fontsize{10.6}{14.5}\selectfont\color{popink},
  fontlower=\popsemi\fontsize{10.6}{14.5}\selectfont\color{popink}}}
% Coche dessinée (le glyphe ✓ n'existe pas dans les polices du document)
\newcommand{\popcheck}[1]{\tikz[baseline=-0.5ex]\draw[#1,line width=1.6pt,line cap=round,line join=round](-0.13,0.0)--(-0.04,-0.09)--(0.13,0.1);}
\providecommand{\coche}{\popcheck{popgreen}}
\providecommand{\resultat}[1]{\par\smallskip\noindent\fcolorbox{popgreen}{popgreenL}{\parbox{\dimexpr\linewidth-2\fboxsep-2\fboxrule\relax}{\textbf{Résultat.} #1}}\par\smallskip}
\newcommand{\popboxhead}[3]{\poppill{#1}{#2}{white}\ifx\relax#3\relax\else\hspace{6pt}{\popxbold\fontsize{10.6}{12}\selectfont\color{popink}#3}\fi\par\vspace{7pt}}

\newtcolorbox{aretenir}[1][]{popbase,colframe=popgreen,before upper={\popboxhead{À retenir}{popgreen}{#1}}}
\newtcolorbox{exemplebox}[1][]{popbase,colframe=poporange,before upper={\popboxhead{Exemple}{poporange}{#1}}}
\newtcolorbox{definition}[1][]{popbase,colframe=popblue,before upper={\popboxhead{Définition}{popblue}{#1}}}
\newtcolorbox{propriete}[1][]{popbase,colframe=poppurple,before upper={\popboxhead{Propriété}{poppurple}{#1}}}
\newtcolorbox{methode}[1][]{popbase,colframe=popgold,before upper={\popboxhead{Méthode}{popgold}{#1}}}
\newtcolorbox{attention}[1][]{popbase,colframe=poppink,before upper={\popboxhead{Attention}{poppink}{#1}}}
\newtcolorbox{experience}[1][]{popbase,colframe=popblue,colback=popblueL,before upper={\popboxhead{Expérience}{popblue}{#1}}}
% « Le savais-tu ? » : carte violette pleine (contexte, culture, Cameroun)
\newtcolorbox{savaistu}[1][]{popbase,colback=poppurple,colframe=popink,
  fontupper=\popsemi\fontsize{10.4}{14.2}\selectfont\color{white},
  before upper={\poppill{Le savais-tu\,?}{white}{poppurple}\ifx\relax#1\relax\else\hspace{6pt}{\popxbold\color{white}#1}\fi\par\vspace{7pt}}}
% Exercice résolu : énoncé en haut, \tcblower puis la solution sur fond crème
\newcounter{popexo}\newcounter{popexr}
\newtcolorbox{exoresolu}[1][]{popbase,colframe=poporange,colbacklower=popcream,
  segmentation style={popink,line width=1.2pt,dashed},
  before upper={\stepcounter{popexr}\popboxhead{Exercice résolu \thepopexr}{poporange}{#1}},
  before lower={\poppill{Solution}{popink}{popcream}\par\vspace{6pt}}}
% Exercice (non résolu en place) — la correction est dans la partie Corrigés
\newtcolorbox{exercice}[1][]{popbase,colframe=popink,
  before upper={\stepcounter{popexo}\popboxhead{Exercice \thepopexo}{popink}{#1}}}
% Corrigé
\newtcolorbox{corrige}[1][]{popbase,colframe=popgreen,colback=popgreenL,
  before upper={\popboxhead{Corrigé}{popgreen}{#1}}}
% Objectifs / prérequis / bilan
\newtcolorbox{objectifs}{popbase,colframe=popink,colback=poporangeL,
  before upper={\popboxhead{Objectifs de la leçon}{popink}{}}}
\newtcolorbox{prerequis}{popbase,colframe=popink,colback=popblueL,
  before upper={\popboxhead{Prérequis}{popblue}{}}}
\newtcolorbox{bilan}{popbase,colframe=popink,colback=white,boxrule=2.8pt,
  before upper={\popboxhead{Fiche bilan}{poporange}{L'essentiel à connaître pour l'examen}}}
% Figure encadrée avec légende
\newtcolorbox{popfigure}[1][]{enhanced,breakable=false,colback=white,colframe=popline,boxrule=1.4pt,arc=8pt,
  left=10pt,right=10pt,top=10pt,bottom=8pt,before skip=10pt,after skip=12pt,halign=center,
  lower separated=false,halign lower=center,
  fontlower=\popsemi\fontsize{9}{11.5}\selectfont\color{popmuted},
  #1}
\newcommand{\legende}[1]{\par\vspace{6pt}{\popsemi\fontsize{9}{11.5}\selectfont\color{popmuted}#1}}

% ============================================================
% Signature OnBuch+
% ============================================================
\newcommand{\poplogoinline}[1]{{\poplogo\fontsize{#1}{#1}\selectfont\color{popink}OnBuch\color{poporange}+}}
\newcommand{\signatureonbuch}{%
  \begin{tcolorbox}[enhanced,colback=popink,colframe=popink,boxrule=2.2pt,arc=10pt,
      drop shadow=poporange,left=14pt,right=14pt,top=10pt,bottom=10pt,before skip=0pt,after skip=0pt]
    \tikz[baseline=-0.7ex]\node[circle,fill=popgreen,draw=popcream,line width=1.3pt,minimum size=22pt,inner sep=0]{\popcheck{white}};%
    \hspace{9pt}\begin{minipage}[c]{0.62\linewidth}
      {\popxbold\fontsize{12}{13}\selectfont\color{popcream}Signé\ \textbullet\ {\poplogo OnBuch\color{poporange}+}}\par\vspace{2pt}
      {\raggedright\popsemi\fontsize{8}{10.5}\selectfont\color{popline}\MakeUppercase{Équipe pédagogique OnBuch+}\\\MakeUppercase{Conforme au programme officiel MINESEC}\par}
    \end{minipage}\hfill
    {\poplogo\fontsize{22}{22}\selectfont\color{popcream}OnBuch\color{poporange}+}
  \end{tcolorbox}%
}

% ============================================================
% Page de garde
% #1 titre · #2 matière · #3 niveau · #4 module · #5 n° leçon · #6 durée ·
% #7 description / objectifs (texte ou itemize)
% ============================================================
\newcommand{\pagegarde}[7]{%
  \thispagestyle{empty}
  \newgeometry{a4paper,margin=1.7cm,top=1.6cm,bottom=1.4cm}%
  \begin{tikzpicture}[remember picture,overlay]
    \fill[poporange!18] ([xshift=-3.2cm,yshift=-3.6cm]current page.north east) circle (4.6cm);
    \fill[poppurple!14] ([xshift=2.4cm,yshift=7.6cm]current page.south west) circle (3.8cm);
  \end{tikzpicture}%
  {\poplogo\fontsize{30}{30}\selectfont\color{popink}OnBuch\color{poporange}+}\hfill
  \raisebox{0.6ex}{\poppill{Cours signé \textbullet\ Premium}{popink}{popcream}}\par
  \vspace{1pt}\popsquiggle{poppurple}\par
  \vspace{8pt}
  \begin{tcolorbox}[enhanced,colback=poporange,colframe=popink,boxrule=2.8pt,arc=13pt,
      drop shadow=popink,left=18pt,right=18pt,top=12pt,bottom=13pt,after skip=10pt]
    \poppill{#2 \textbullet\ #3}{popink}{popcream}\hspace{5pt}\poppill{#4}{white}{popink}\par\vspace{8pt}
    {\popdisplay\fontsize{14}{14}\selectfont\color{popink}LEÇON #5}\par\vspace{4pt}
    {\raggedright\hyphenpenalty=10000\exhyphenpenalty=10000\popdisplay\fontsize{25}{27}\selectfont\color{popcream}\MakeUppercase{#1}\par}
    \vspace{8pt}
    {\popxbold\fontsize{9.5}{9.5}\selectfont\color{popink}\MakeUppercase{Durée conseillée : #6}}
  \end{tcolorbox}
  \begin{tcolorbox}[enhanced,colback=white,colframe=popink,boxrule=2.2pt,arc=10pt,
      drop shadow=popink,left=16pt,right=16pt,top=10pt,bottom=10pt,after skip=8pt]
    \poppill{Dans cette leçon}{poppurple}{white}\par\vspace{5pt}
    \popsemi\fontsize{9.3}{12.6}\selectfont\color{popink}#7
  \end{tcolorbox}
  \noindent\begin{minipage}[t]{0.487\linewidth}
    \begin{tcolorbox}[enhanced,colback=popgreen,colframe=popink,boxrule=2.2pt,arc=10pt,
        drop shadow=popink,left=11pt,right=11pt,top=10pt,bottom=10pt,height=3.1cm,valign=center]
      \tcbox[colback=white,colframe=popink,boxrule=1.3pt,arc=5pt,boxsep=0pt,nobeforeafter,
        left=3pt,right=3pt,top=3pt,bottom=3pt,valign=center]{\qrcode[height=1.9cm]{https://play.google.com/store/apps/details?id=com.luvvix.onbuch&hl=fr}}%
      \hspace{8pt}\begin{minipage}[c]{0.5\linewidth}
        {\popdisplay\fontsize{11}{12.5}\selectfont\color{white}\MakeUppercase{Télécharge OnBuch}}\par\vspace{4pt}
        {\raggedright\popsemi\fontsize{8.4}{11}\selectfont\color{white}Gratuit \textbullet\ Google Play\\Tuteur IA, annales, concours, résultats\par}
      \end{minipage}
    \end{tcolorbox}
  \end{minipage}\hfill
  \begin{minipage}[t]{0.487\linewidth}
    \begin{tcolorbox}[enhanced,colback=poppurple,colframe=popink,boxrule=2.2pt,arc=10pt,
        drop shadow=popink,left=11pt,right=11pt,top=10pt,bottom=10pt,height=3.1cm,valign=center]
      \tikz[baseline=-0.7ex]\node[circle,fill=white,draw=popink,line width=1.3pt,minimum size=1.6cm,inner sep=0]{\popdisplay\fontsize{24}{24}\selectfont\color{poppurple}+};%
      \hspace{8pt}\begin{minipage}[c]{0.6\linewidth}
        {\popdisplay\fontsize{11}{12.5}\selectfont\color{white}\MakeUppercase{Passe à OnBuch+}}\par\vspace{4pt}
        {\raggedright\popsemi\fontsize{8.4}{11}\selectfont\color{white}Tous les cours signés, Tuteur IA illimité, fiches PDF — onglet OnBuch+ de l'app\par}
      \end{minipage}
    \end{tcolorbox}
  \end{minipage}
  \vspace{8pt}
  \popcontenu
  \vfill
  \signatureonbuch
  \restoregeometry
  \newpage
}

% Rangée « Ce cours contient » (page de garde)
\newcommand{\poptile}[3]{% #1 couleur #2 gros texte #3 légende
  \begin{tcolorbox}[enhanced,colback=#1,colframe=popink,boxrule=2pt,arc=9pt,shadow={2.5pt}{-2.5pt}{0pt}{popink},
      left=6pt,right=6pt,top=9pt,bottom=9pt,halign=center,valign=center,height=1.75cm,width=\linewidth,nobeforeafter]
    {\popdisplay\fontsize{12.5}{13}\selectfont\color{white}\MakeUppercase{#2}}\par\vspace{3pt}
    {\centering\popsemi\fontsize{7.8}{9.5}\selectfont\color{white}#3\par}
  \end{tcolorbox}}
\newcommand{\popcontenu}{%
  \par\noindent{\popxboldls\fontsize{7.5}{7.5}\selectfont\color{popmuted}CE COURS SIGNÉ CONTIENT}\par\vspace{7pt}
  \noindent\begin{minipage}[t]{0.235\linewidth}\poptile{popblue}{Cours}{complet, illustré, pas à pas}\end{minipage}\hfill
  \begin{minipage}[t]{0.235\linewidth}\poptile{poporange}{Méthodes}{et exercices résolus}\end{minipage}\hfill
  \begin{minipage}[t]{0.235\linewidth}\poptile{poppink}{Exercices}{type examen + corrigés}\end{minipage}\hfill
  \begin{minipage}[t]{0.235\linewidth}\poptile{popgreen}{Fiche bilan}{l'essentiel à retenir}\end{minipage}\par}

% Sommaire
\newtcolorbox{sommaire}{enhanced,colback=white,colframe=popink,boxrule=2.2pt,arc=10pt,
  drop shadow=popink,left=18pt,right=18pt,top=13pt,bottom=13pt,before skip=6pt,after skip=20pt,
  before upper={\poppill{Sommaire}{poppurple}{white}\par\vspace{9pt}\popsemi\fontsize{10.6}{14.5}\selectfont\color{popink}}}

% Signature de fin de cours — poussée tout en bas de la dernière page
\newcommand{\findecours}{%
  \par\vfill\nopagebreak
  \begin{center}\popsquiggle{poporange}\end{center}\vspace{4pt}
  \signatureonbuch
}
\AtBeginDocument{\def\llbracket{\mathopen{[\mkern-3.5mu[}}\def\rrbracket{\mathclose{]\mkern-3.5mu]}}} % intervalles d'entiers (glyphe ⟦ absent de la police)
% Glyphes absents de Fira Math : redéfinis à partir de glyphes disponibles
\AtBeginDocument{%
  \def\setminus{\mathbin{\backslash}}\let\smallsetminus\setminus
  \def\vdots{\mathord{\vbox{\baselineskip3.2pt\lineskiplimit0pt\kern4pt\hbox{.}\hbox{.}\hbox{.}}}}%
  \def\ddots{\mathinner{\mkern1mu\raise6.5pt\hbox{.}\mkern2mu\raise3.6pt\hbox{.}\mkern2mu\raise0.7pt\hbox{.}\mkern1mu}}%
  \def\triangle{\mathord{\scalebox{0.9}{$\symup{\Delta}$}}}%
  \def\nabla{\mathord{\raisebox{\depth}{\scalebox{1}[-1]{$\symup{\Delta}$}}}}%
  \def\checkmark{\popcheck{popdark}}%
  \def\star{\mathbin{\ast}}\def\bigstar{\mathord{\ast}}%
  \def\diamond{\mathbin{\scalebox{0.6}{$\Diamond$}}}%
  \def\bigcup{\mathop{\mathchoice{\raisebox{-0.3em}{\scalebox{1.7}{$\cup$}}}{\scalebox{1.25}{$\cup$}}{\cup}{\cup}}\displaylimits}%
  \def\bigcap{\mathop{\mathchoice{\raisebox{-0.3em}{\scalebox{1.7}{$\cap$}}}{\scalebox{1.25}{$\cap$}}{\cap}{\cap}}\displaylimits}%
  \def\lesseqqgtr{\mathrel{\lessgtr}}\def\bigoplus{\mathop{\oplus}\displaylimits}%
}
\providecommand{\ssi}{si et seulement si} % abréviation fréquente
\AtBeginDocument{\providecommand{\wideparen}{\overparen}} % arc de cercle

%% FIGFIX-BEGIN v4 — correctifs globaux des figures (docs/onbuch-generation/tools/figfix)
\usepackage{adjustbox}\usepackage{etoolbox}
% toute figure TikZ/pgfplots plus large que la ligne est réduite à la largeur disponible
\BeforeBeginEnvironment{tikzpicture}{\begin{adjustbox}{max width=\linewidth}}
\AfterEndEnvironment{tikzpicture}{\end{adjustbox}}
% légendes pgfplots lisibles par-dessus les courbes
\pgfplotsset{every axis/.append style={legend style={fill=white,fill opacity=0.92,text opacity=1,draw=black!15,inner sep=3pt}}}
% étiquettes d'arêtes (syntaxe edge["..."]) avec fond blanc
\tikzset{every edge quotes/.append style={fill=white,inner sep=1.2pt}}
%% FIGFIX-END
\begin{document}

\pagegarde{Les Entreprises}{Philosophie}{1re Tech}{Philosophie – classe de Première technique}{N12}{2 séances (fiche de progression harmonisée)}{Cette leçon analyse la notion d'entreprise comme unité de production et organisation sociale, en distinguant ses formes juridiques et ses modes de fonctionnement interne. Elle permet à l'élève de mobiliser des concepts philosophiques (finalité, responsabilité, organisation) pour analyser des situations économiques concrètes du Cameroun.
\vspace{4pt}
\begin{multicols}{2}\small
\begin{itemize}
\item À la fin de cette leçon, je sais définir l'entreprise en distinguant ses dimensions technique, juridique et humaine.
\item À la fin de cette leçon, je sais classifier les entreprises selon des critères pertinents (taille, secteur, forme juridique, capital) en m'appuyant sur le droit camerounais (OHADA).
\item À la fin de cette leçon, j'identifie les facteurs de production (travail, capital, organisation) et explique leur combinaison dans le processus de création de valeur.
\item À la fin de cette leçon, j'analyse la structure organisationnelle d'une entreprise (organigramme, services, flux d'information) à partir de documents réels.
\item À la fin de cette leçon, je comprends les enjeux de la gouvernance (associés, dirigeants, conseil d'administration) et de la responsabilité sociale (RSE).
\item À la fin de cette leçon, je sais interpréter des documents chiffrés (compte de résultat simplifié, effectifs) et graphiques (répartition sectorielle) pour poser un diagnostic.
\end{itemize}\end{multicols}}

\begin{sommaire}
\begin{multicols}{2}
\begin{enumerate}
\item 1. Qu'est-ce qu'une entreprise ? Définition et nature
\item 2. Typologie et classification des entreprises au Cameroun
\item 3. Le fonctionnement interne : Facteurs de production et Organisation
\item 4. Gouvernance, Stratégie et Responsabilité Sociétale (RSE)
\item 5. TP Intégré : Diagnostic simplifié d'une PME camerounaise (Cas 'Bois \& Design')
\item Activité d'intégration
\item Méthodes et exercices résolus
\item Exercices
\item Corrigés des exercices
\item Fiche bilan
\end{enumerate}
\end{multicols}
\end{sommaire}

\begin{objectifs}
\begin{itemize}
\item À la fin de cette leçon, je sais définir l'entreprise en distinguant ses dimensions technique, juridique et humaine.
\item À la fin de cette leçon, je sais classifier les entreprises selon des critères pertinents (taille, secteur, forme juridique, capital) en m'appuyant sur le droit camerounais (OHADA).
\item À la fin de cette leçon, j'identifie les facteurs de production (travail, capital, organisation) et explique leur combinaison dans le processus de création de valeur.
\item À la fin de cette leçon, j'analyse la structure organisationnelle d'une entreprise (organigramme, services, flux d'information) à partir de documents réels.
\item À la fin de cette leçon, je comprends les enjeux de la gouvernance (associés, dirigeants, conseil d'administration) et de la responsabilité sociale (RSE).
\item À la fin de cette leçon, je sais interpréter des documents chiffrés (compte de résultat simplifié, effectifs) et graphiques (répartition sectorielle) pour poser un diagnostic.
\item À la fin de cette leçon, je sais rédiger une argumentation structurée justifiant le choix d'une forme juridique ou d'une stratégie organisationnelle face à un cas pratique.
\item À la fin de cette leçon, je sais construire un croquis organisationnel ou un diagramme de flux à partir d'une description textuelle d'une PME camerounaise.
\end{itemize}
\end{objectifs}

\begin{prerequis}
\begin{itemize}
\item Notion d'activité économique et d'agents économiques (ménages, entreprises, État) vue en 4ème / 3ème.
\item Distinction entre bien et service, besoin et désir.
\item Notions de base de logique argumentative (thèse, argument, exemple, conclusion).
\item Repères géographiques et économiques du Cameroun (régions, principales villes, secteurs porteurs).
\item Vocabulaire de base : capital, travail, salaire, profit, marché.
\end{itemize}
\end{prerequis}

\newpage

\coursec{Qu'est-ce qu'une entreprise ? Définition et nature}

\courssub{Situation d'accroche et problématique}

À Douala, Mme Ngo organise la transformation artisanale de plantains en chips dans sa cour et vend ses produits aux collègues de son mari ; à Yaoundé, la société anonyme « Cameroun Telecom » emploie des centaines d'ingénieurs et gère un budget de plusieurs milliards de FCFA. Toutes deux produisent des biens ou des services, mais leur organisation, leur responsabilité juridique et leur pouvoir de décision diffèrent radicalement. Mme Ngo confond son argent personnel avec celui de son activité ; la SA, elle, possède un patrimoine propre, des actionnaires, un conseil d'administration.

Comment définir l'entreprise au-delà de cette diversité ? Quels sont les rouages internes qui permettent à toute entreprise, petite ou grande, de fonctionner et de durer ? Telle est la problématique de cette première section : nous chercherons une définition rigoureuse de l'entreprise, puis nous en dégagerons la triple nature (technique, juridique, humaine), avant de distinguer soigneusement entreprise, société et établissement, et de situer la figure de l'entrepreneur dans le cadre légal camerounais.

\courssub{Intuition : de la cour de Mme Ngo à la définition}

Observons d'abord ce que les deux structures de notre accroche ont en commun. Mme Ngo achète des plantains (matière première), utilise une friteuse et du gaz (capital technique), travaille avec sa fille (travail), et vend des chips sur un marché. Cameroun Telecom achète des équipements réseau, emploie des salariés, et vend des services de communication. Dans les deux cas, on trouve : une \cle{combinaison de facteurs} (travail + capital), une \cle{production} destinée à être \cle{vendue}, et une personne qui \cle{décide} (quoi produire, à quel prix, comment).

C'est à partir de cette observation que les économistes construisent la définition générale : l'entreprise n'est pas d'abord un bâtiment ni une taille, c'est une \emph{unité organisée de production marchande dotée d'une autonomie de décision}.

\begin{definition}[L'entreprise]
L'\cle{entreprise} est une unité économique et sociale, dotée d'une \cle{autonomie de décision}, qui combine des \cle{facteurs de production} (travail et capital) pour créer de la \cle{valeur} destinée au \cle{marché}, dans une finalité généralement lucrative.
\end{definition}

Décomposons cette définition :
\begin{itemize}
\item \textbf{Unité économique} : elle produit des biens (chips, ciment, meubles) ou des services (transport, téléphonie, coiffure) destinés à être vendus à un prix couvrant au moins les coûts.
\item \textbf{Unité sociale} : elle réunit des hommes et des femmes qui coopèrent selon des règles et une hiérarchie.
\item \textbf{Autonomie de décision} : le dirigeant choisit librement ses produits, ses prix, ses fournisseurs. C'est ce qui distingue l'entreprise d'une simple administration exécutant un budget décidé ailleurs.
\item \textbf{Finalité généralement lucrative} : elle cherche en principe un \cle{profit}, c'est-à-dire un excédent des recettes sur les coûts. Nous verrons que cette finalité comporte d'importantes nuances.
\end{itemize}

\begin{methode}[Comment reconnaître une entreprise ?]
Pour qualifier une structure d'\emph{entreprise}, vérifier trois critères \textbf{cumulatifs} :
\begin{enumerate}
\item \textbf{Autonomie de gestion et de décision} : la structure décide elle-même de sa production, de ses prix et de ses achats.
\item \textbf{Combinaison de facteurs de production} : elle réunit du travail et du capital (machines, locaux, fonds) pour produire.
\item \textbf{Marchandise du produit ou du service} : la production est vendue sur un marché, à un prix, et non distribuée gratuitement.
\end{enumerate}
Si l'un des trois critères manque, on n'est pas en présence d'une entreprise au sens strict (exemple : une école publique gratuite ne vend pas son service).
\end{methode}

\begin{exoresolu}[Mme Ngo est-elle chef d'entreprise ?]
Énoncé : Mme Ngo transforme des plantains en chips dans sa cour et les vend aux collègues de son mari. Une voisine affirme : « Ce n'est pas une entreprise, c'est juste une activité de famille. » Qui a raison ? Justifier à l'aide des trois critères.
\tcblower
Solution détaillée : appliquons la méthode des trois critères cumulatifs.
\begin{enumerate}
\item \emph{Autonomie de décision} : Mme Ngo choisit elle-même ses quantités, ses prix de vente et ses fournisseurs de plantains. Critère rempli.
\item \emph{Combinaison de facteurs} : elle réunit du travail (le sien et celui de sa fille) et du capital (friteuse, gaz, ustensiles). Critère rempli.
\item \emph{Marchandise} : les chips sont vendues contre de l'argent, à un prix. Critère rempli.
\end{enumerate}
Conclusion : l'activité de Mme Ngo est bien une \textbf{entreprise}, au sens économique du terme, même si elle est de très petite taille et informelle. La taille n'entre pas dans la définition. La voisine a donc tort : elle confond « entreprise » et « grande société », confusion que nous dissiperons plus loin.
\end{exoresolu}

\courssub{La triple dimension de l'entreprise}

L'entreprise est une réalité complexe qu'on ne peut enfermer dans un seul regard. Les sciences économiques et sociales en distinguent trois dimensions complémentaires.

\textbf{1. La dimension technique.}\par
L'entreprise est d'abord un \cle{appareil de production} : elle maîtrise des \emph{procédés} (recettes, machines, savoir-faire, organisation du travail) qui transforment des \emph{intrants} (matières premières, énergie, informations) en \emph{extrants} (biens et services). Le menuisier de Bafoussam transforme le bois en meubles ; l'opérateur télécom transforme des signaux en communication. C'est la division du travail et la combinaison technique des facteurs qui créent la \cle{productivité}.

\textbf{2. La dimension juridique.}\par
L'entreprise peut être dotée d'une \cle{personnalité morale} : le droit la traite alors comme une « personne » fictive, distincte des personnes physiques qui la composent. Elle possède un \cle{patrimoine propre} (biens, comptes bancaires, dettes), peut signer des contrats, ester en justice. Cette dimension est encadrée au Cameroun par le droit \cle{OHADA} (Organisation pour l'Harmonisation en Afrique du Droit des Affaires), dont l'Acte uniforme relatif au droit des sociétés commerciales et du GIE, révisé en 2014, définit les formes sociales (SARL, SA, GIE...), et par le Code du travail camerounais pour les relations avec les salariés.

\textbf{3. La dimension humaine.}\par
L'entreprise est un \cle{collectif de travail} : des individus y coopèrent, mais aussi y négocient et parfois s'y affrontent. Elle est traversée par des \emph{relations de pouvoir} (hiérarchie, autorité du dirigeant) et des \emph{relations de coopération} (esprit d'équipe, syndicats, délégués du personnel). Un dirigeant qui ignore cette dimension — motivation, conflits, communication — condamne son entreprise, même si la technique et le droit sont parfaits.

\begin{popfigure}
\begin{center}
\begin{tikzpicture}[scale=0.9]
\begin{scope}[blend group=soft light]
\fill[popblue!40] (90:1.6) circle (2.4);
\fill[poporange!40] (210:1.6) circle (2.4);
\fill[popgreen!40] (330:1.6) circle (2.4);
\end{scope}
\draw[popblue,thick] (90:1.6) circle (2.4);
\draw[poporange,thick] (210:1.6) circle (2.4);
\draw[popgreen,thick] (330:1.6) circle (2.4);
\node[popblue,font=\bfseries] at (90:3.6) {Unité de production};
\node[popblue,font=\small] at (90:2.9) {(technique)};
\node[poporange,font=\bfseries,align=center] at (210:3.4) {Personne\\morale};
\node[poporange,font=\small] at (210:2.7) {(juridique)};
\node[popgreen,font=\bfseries,align=center] at (330:3.4) {Collectif\\humain};
\node[popgreen,font=\small] at (330:2.7) {(social)};
\node[font=\bfseries\large,popdark] at (0,0) {ENTREPRISE};
\end{tikzpicture}
\end{center}
\legende{Diagramme de Venn : l'entreprise se situe à l'intersection de trois ensembles — l'unité technique de production, la personne morale juridique et le collectif humain de travail.}
\end{popfigure}

\begin{popfigure}
\begin{center}
\begin{tikzpicture}[mindmap, grow cyclic, every node/.style={concept}, concept color=popblue!50, text=popdark,
level 1/.append style={level distance=3.4cm, sibling angle=120},
level 2/.append style={level distance=2.2cm, sibling angle=45}]
\node[font=\bfseries] {L'ENTREPRISE}
child[concept color=poporange!50] { node {Technique}
  child { node[font=\small] {Procédés, machines} }
  child { node[font=\small] {Division du travail} }
  child { node[font=\small] {Productivité} }
}
child[concept color=popgreen!50] { node {Juridique}
  child { node[font=\small] {Personnalité morale} }
  child { node[font=\small] {Patrimoine propre} }
  child { node[font=\small] {Droit OHADA} }
}
child[concept color=poppurple!50] { node {Humaine}
  child { node[font=\small] {Collectif de travail} }
  child { node[font=\small] {Pouvoir, hiérarchie} }
  child { node[font=\small] {Coopération, conflits} }
};
\end{tikzpicture}
\end{center}
\legende{Carte heuristique (mindmap) : les trois dimensions de l'entreprise et leurs sous-branches.}
\end{popfigure}

\begin{aretenir}[L'essentiel de la triple dimension]
Une entreprise est à la fois un \textbf{outil} (dimension technique : elle produit), une \textbf{personne fictive} (dimension juridique : elle a des droits et des obligations) et un \textbf{groupe humain} (dimension sociale : elle organise la coopération et le pouvoir). Oublier l'une de ces dimensions, c'est ne pas comprendre ce qu'est une entreprise.
\end{aretenir}

\courssub{Entreprise, société, établissement : une distinction fondamentale}

C'est ici que se situe le piège classique du Probatoire et du Baccalauréat technique. Trois mots du langage courant désignent des réalités différentes :

\begin{definition}[Trois notions à ne pas confondre]
\begin{itemize}
\item L'\cle{entreprise} est une \textbf{unité économique} : une réalité réelle, un ensemble de moyens humains et matériels organisés pour produire.
\item La \cle{société} est une \textbf{personne morale} : une réalité juridique, créée par un contrat entre associés, dotée d'un patrimoine propre.
\item L'\cle{établissement} est une \textbf{unité géographique} : un lieu de production localisé (usine, atelier, magasin, agence).
\end{itemize}
\end{definition}

\begin{propriete}[Distinction entreprise / société / établissement]
Une même personne morale (société) peut exploiter \textbf{plusieurs entreprises} : c'est le cas d'une \emph{holding} qui détient une entreprise de transport et une entreprise agroalimentaire. Inversement, un même établissement peut abriter \textbf{plusieurs entreprises} : dans une zone industrielle de Douala, un atelier de sous-traitance peut travailler pour plusieurs donneurs d'ordre sous le même toit. Enfin, une entreprise peut comporter plusieurs établissements (une banque et ses agences).
\end{propriete}

\begin{exemplebox}[Illustration camerounaise]
La société « SABC » (Société Anonyme des Brasseries du Cameroun) est \emph{une} personne morale, mais elle exploite \emph{plusieurs} établissements : des brasseries à Douala, Yaoundé, Bafoussam, Garoua. Chaque brasserie est un établissement ; l'ensemble forme l'entreprise ; la SA est la forme juridique. De même, une banque comme Afriland First Bank (une société) possède des dizaines d'agences (des établissements) réparties sur le territoire national.
\end{exemplebox}

\begin{attention}[Piège fréquent : entreprise individuelle et société unipersonnelle]
Ne pas confondre :
\begin{itemize}
\item L'\textbf{entreprise individuelle} (celle de Mme Ngo) : il n'y a \emph{aucune distinction} entre le patrimoine personnel et le patrimoine professionnel. En cas de dettes, le créancier peut saisir les biens personnels de l'entrepreneur.
\item La \textbf{société unipersonnelle} (ex. SARL unipersonnelle) : même avec un seul associé, il existe une \emph{personnalité morale distincte} et la responsabilité est \emph{limitée} à l'apport. Le patrimoine personnel est protégé.
\end{itemize}
« Un seul homme » ne signifie donc pas automatiquement « entreprise individuelle » : tout dépend de la forme juridique choisie.
\end{attention}

\courssub{La finalité : lucrative ou non ?}

La définition de l'entreprise mentionne une finalité « généralement lucrative ». Cette prudence de langage est essentielle pour l'examen.

\textbf{La finalité lucrative.}\par
La plupart des entreprises recherchent le \cle{profit} : la différence entre les recettes (chiffre d'affaires) et les coûts (salaires, matières, loyers, impôts). Le profit rémunère le risque de l'entrepreneur, permet d'investir et de récompenser les associés (dividendes). Sans profit durable, l'entreprise disparaît : c'est la logique de l'économie de marché.

\textbf{Les finalités non lucratives.}\par
Certaines organisations produisent des biens ou des services sans rechercher le profit :
\begin{itemize}
\item les \textbf{associations} et \textbf{ONG} (loi camerounaise de 1990 sur la liberté d'association) : elles poursuivent un but social, culturel ou humanitaire ; un éventuel excédent est réinvesti dans l'objet social, jamais distribué ;
\item les \textbf{établissements publics à caractère industriel et commercial} (EPIC) : ils vendent des services (eau, électricité) mais appartiennent à l'État et poursuivent une mission de service public ;
\item l'\textbf{économie sociale et solidaire} (coopératives, mutuelles) : elles dégagent des excédents mais privilégient l'intérêt collectif des membres.
\end{itemize}

\begin{attention}[Nuance exigible au Bac]
Ne jamais écrire : « toute entreprise cherche le profit ». La formulation correcte est : « la \emph{plupart} des entreprises ont une finalité lucrative, mais certaines organisations productives (associations, coopératives, EPIC) poursuivent des finalités non lucratives ou d'intérêt général ». Le critère décisif reste la \emph{production marchande autonome}, pas le seul profit.
\end{attention}

\courssub{L'entrepreneur : porteur de projet, preneur de risque, innovateur}

Au cœur de toute entreprise se trouve une personne (ou une équipe) : l'\cle{entrepreneur}. Trois traits le définissent.

\textbf{1. Le porteur de projet.}\par
L'entrepreneur repère un besoin insatisfait et imagine une réponse : Mme Ngo voit que ses voisins aiment les chips de plantain ; un jeune de Buea crée une application de livraison. Il rassemble ensuite les moyens (capitaux, compétences) pour transformer l'idée en activité.

\textbf{2. Le preneur de risque.}\par
Contrairement au salarié, qui reçoit un salaire fixe, l'entrepreneur engage ses biens sans garantie de gain : si les chips ne se vendent pas, la perte est pour lui. L'économiste Richard Cantillon, dès le XVIII\textsuperscript{e} siècle, définissait déjà l'entrepreneur par cette \emph{incertitude} assumée : il achète à prix certain (matières, salaires) pour revendre à prix incertain.

\textbf{3. L'innovateur.}\par
Joseph \cle{Schumpeter} (1883-1950) a fait de l'entrepreneur le moteur du capitalisme : par l'\emph{innovation} (produit nouveau, procédé nouveau, marché nouveau, organisation nouvelle), il détruit les anciennes combinaisons économiques — c'est la célèbre « \cle{destruction créatrice} ». Le téléphone mobile a ainsi « détruit » les cabines téléphoniques au Cameroun, tout en créant des milliers d'emplois (revendeurs de crédit, réparateurs, opérateurs Mobile Money).

\begin{savaistu}[Le GIE : entreprendre ensemble sans perdre son autonomie]
Le \cle{Groupement d'Intérêt Économique} (GIE), prévu par le droit OHADA, permet à des entrepreneurs indépendants de \emph{mutualiser des moyens} sans fusionner ni perdre leur autonomie juridique. Exemple : des producteurs de café du Nord-Ouest créent un GIE pour acheter ensemble les engrais (moins cher), partager un camion et exporter en commun. Chacun reste propriétaire de sa propre exploitation ; le GIE n'est qu'un outil commun. C'est une forme juridique précieuse pour les petits producteurs camerounais confrontés aux coûts élevés de l'exportation.
\end{savaistu}

\courssub{Le poids réel des entreprises au Cameroun}

\begin{exemplebox}[Un tissu économique de très petites unités]
Selon l'INS (Recensement général des entreprises, 2016, confirmé par les enquêtes récentes), environ \textbf{98\,\%} des entreprises camerounaises sont des \cle{très petites entreprises} (TPE) et \cle{PME} de moins de 50 employés, souvent informelles, comme l'atelier de Mme Ngo. Pourtant, les \textbf{grandes entreprises} (GE), très peu nombreuses, concentrent environ \textbf{70\,\%} de la valeur ajoutée produite dans le secteur formel. Leçon à retenir : la \emph{nombre} ne coïncide pas avec le \emph{poids économique}. Cette dualité (un océan de micro-unités, quelques grands pôles) est une caractéristique majeure de l'économie camerounaise.
\end{exemplebox}

\begin{center}
\begin{tabularx}{\linewidth}{|l|X|X|}
\hline
\rowcolor{popblueL} \textbf{Critère} & \textbf{TPE / PME ($\approx$ 98\,\% des unités)} & \textbf{Grandes entreprises} \\
\hline
Effectif & moins de 50 salariés & 100 salariés et plus \\
\hline
Part de la valeur ajoutée & environ 30\,\% & environ 70\,\% \\
\hline
Exemples & atelier de chips, menuiserie, salon de coiffure & SABC, ENEO, grandes banques \\
\hline
Forme juridique fréquente & entreprise individuelle, informel & SA, SARL \\
\hline
\end{tabularx}
\end{center}

\courssub{Synthèse et ouverture}

\begin{aretenir}[Ce qu'il faut retenir de la section]
\begin{enumerate}
\item L'entreprise est une unité économique et sociale, autonome dans ses décisions, qui combine travail et capital pour produire des biens ou services destinés au marché.
\item Elle possède une triple dimension : technique (procédés de production), juridique (personnalité morale, patrimoine propre, droit OHADA), humaine (collectif de travail, pouvoir et coopération).
\item Entreprise (unité économique), société (personne morale) et établissement (unité géographique) sont trois notions distinctes qu'une même réalité peut combiner de multiples façons.
\item La finalité lucrative domine, mais n'est pas universelle (associations, coopératives, EPIC).
\item L'entrepreneur — porteur de projet, preneur de risque, innovateur selon Schumpeter — est la figure centrale de la création d'entreprise.
\end{enumerate}
\end{aretenir}

Nous savons désormais \emph{ce qu'est} une entreprise. Mais toutes ne se ressemblent pas : de la cour de Mme Ngo à la SA aux milliards de FCFA, les tailles, les secteurs et les formes juridiques varient considérablement. La section suivante proposera une \emph{typologie} des entreprises au Cameroun.

\begin{exercice}[Pour s'entraîner]
1. Le lycée technique de ta ville possède un atelier où les élèves fabriquent des tabourets vendus aux parents d'élèves. Cet atelier est-il une entreprise ? Argumente à l'aide des trois critères cumulatifs.
2. « Une société ne peut exploiter qu'une seule entreprise. » Corrige cette affirmation.
3. Explique pourquoi Schumpeter qualifie l'entrepreneur d'agent de « destruction créatrice », en t'appuyant sur un exemple camerounais.
\end{exercice}

\begin{corrige}[Exercice]
1. \emph{Autonomie} : l'atelier dépend des décisions de l'administration scolaire (programmes, horaires), donc l'autonomie de décision est limitée ; la vente des tabourets relève d'une activité pédagogique accessoire. Le critère d'autonomie fait défaut : ce n'est pas une entreprise au sens strict, mais une activité productive pédagogique. (Accepter la discussion nuancée.)
2. Faux : une même personne morale peut exploiter plusieurs entreprises (cas d'une holding) et une entreprise peut comporter plusieurs établissements. Le lien entre forme juridique et unité économique n'est pas biunivoque.
3. Pour Schumpeter, l'innovation de l'entrepreneur rend obsolètes les productions anciennes tout en créant de nouvelles activités : l'arrivée du téléphone mobile et du Mobile Money au Cameroun a fait disparaître les cabines téléphoniques et une partie du courrier traditionnel, mais a créé des milliers de points de vente de crédit, de kiosques de transfert d'argent et d'emplois de réparation. La destruction des anciennes combinaisons est le prix de la création de richesses nouvelles.
\end{corrige}

\coursec{Typologie et classification des entreprises au Cameroun}

Dans la section précédente, nous avons défini l'entreprise comme une unité économique, juridiquement autonome, qui combine des facteurs de production pour offrir des biens et des services sur un marché en vue de réaliser un profit. Mais dire « entreprise », c'est encore parler de manière abstraite. En réalité, le paysage économique camerounais est d'une diversité étonnante : il n'y a rien de commun, en apparence, entre la vendeuse de beignets de quartier à Mokolo, la menuiserie « Bois \& Design » de Douala, la SONARA de Limbé et la filiale camerounaise d'Orange. Comment dès lors ordonner cette diversité ? Comment un économiste, un juriste ou un décideur public peut-il « classer » les entreprises ? C'est toute la question de la \cle{typologie} : il n'existe pas une classification unique, mais plusieurs \cle{critères de classification} que l'on croise selon l'objectif poursuivi. Nous étudierons successivement le critère de la taille, le critère sectoriel, le critère juridique, le critère de la propriété et de l'origine du capital, avant d'examiner le cas particulier et massif du secteur informel.

\courssub{Le critère de la taille : TPE, PME et Grandes Entreprises}

\textbf{L'intuition.} La première manière de distinguer les entreprises est la plus spontanée : la taille. Mais comment mesurer la taille d'une entreprise ? Par son nombre de salariés ? Par son chiffre d'affaires ? Par la valeur de ses machines ? Le droit camerounais a tranché en retenant deux indicateurs combinés : l'\cle{effectif} (nombre de salariés permanents) et le \cle{chiffre d'affaires} annuel (CA).

\textbf{L'encadrement légal.} Le Décret n°2010/0882/PM fixe les seuils officiels applicables au Cameroun. On distingue trois catégories :

\begin{definition}[Classes de taille des entreprises au Cameroun (Décret n°2010/0882/PM)]
\begin{itemize}
\item \cle{Très Petite Entreprise (TPE)} : moins de 10 salariés permanents et un chiffre d'affaires annuel inférieur à 50 millions de FCFA.
\item \cle{Petite et Moyenne Entreprise (PME)} : entre 10 et 199 salariés permanents, et un chiffre d'affaires compris entre 50 millions et 10 milliards de FCFA.
\item \cle{Grande Entreprise (GE)} : 200 salariés permanents et plus, \textbf{ou} un chiffre d'affaires supérieur à 10 milliards de FCFA.
\end{itemize}
\end{definition}

\begin{aretenir}[Le cœur du dispositif — Exigible au Probatoire]
Une PME camerounaise occupe entre 10 et 199 salariés permanents et réalise un chiffre d'affaires annuel entre 50 millions et 10 milliards de FCFA. Les deux seuils (effectif \textbf{et} CA) doivent être connus par cœur pour l'examen. Attention : pour la GE, une \textbf{seule} des deux conditions suffit (effectif $\ge$ 200 \textbf{ou} CA $>$ 10 Mds).
\end{aretenir}

\textbf{Pourquoi cette classification est-elle décisive ?} Parce qu'elle révèle un paradoxe majeur de l'économie camerounaise. Les statistiques de l'INS (Recensement Général des Entreprises) montrent que l'écrasante majorité des entreprises enregistrées sont des TPE : elles représentent environ 90 \% du nombre d'unités, mais une part modeste de la valeur ajoutée produite. À l'inverse, quelques centaines de grandes entreprises concentrent l'essentiel de la richesse créée, des exportations et des recettes fiscales. Entre les deux, les PME structurées sont relativement peu nombreuses : les économistes parlent du \emph{« missing middle »} (le « maillon manquant »). Ce ventre creux du tissu productif est un problème : ce sont normalement les PME qui industrialisent un pays, créent des emplois stables et qualifiés, et servent de sous-traitants aux grands groupes. D'où les politiques publiques d'appui aux PME (guichets uniques, BC-PME, fonds de garantie).

\begin{popfigure}
\begin{center}
\begin{tikzpicture}
\begin{axis}[
ybar, bar width=0.85cm,
width=0.95\linewidth, height=6.5cm,
symbolic x coords={TPE,PME,GE},
xtick=data,
ymin=0, ymax=100,
ylabel={Part (en \%)},
legend style={at={(0.5,-0.22)},anchor=north,legend columns=2,font=\small},
nodes near coords, every node near coord/.append style={font=\small},
]
\addplot[fill=popblue] coordinates {(TPE,90) (PME,8) (GE,2)};
\addlegendentry{Part dans le nombre d'entreprises}
\addplot[fill=poporange] coordinates {(TPE,25) (PME,30) (GE,45)};
\addlegendentry{Part dans la valeur ajoutée}
\end{axis}
\end{tikzpicture}
\end{center}
\legende{Répartition approximative des entreprises et de la valeur ajoutée par classe de taille au Cameroun (ordres de grandeur INS, RGE). Les TPE dominent en nombre (90 \%), les GE dominent en richesse produite (45 \%) : c'est le « missing middle ».}
\end{popfigure}

\begin{attention}[Taille et forme juridique : deux choses différentes — Piège classique]
Une « Grande Entreprise » n'est pas nécessairement une « Société Anonyme ». Une SARL peut être une GE si elle dépasse les seuils d'effectif ou de chiffre d'affaires ; inversement, une SA peut être une PME si elle reste sous les seuils. Ne confondez \textbf{jamais} le critère de taille (économique, Décret 2010) et le critère juridique (forme sociale, Acte uniforme OHADA).
\end{attention}

\courssub{Le critère sectoriel : les trois secteurs d'activité}

\textbf{L'intuition.} Une deuxième manière de classer consiste à demander : que produit l'entreprise ? D'où tire-t-elle sa matière première ? L'économiste Colin Clark a proposé une classification en trois \cle{secteurs}, reprise dans les nomenclatures statistiques modernes (type NACE, adaptée au Cameroun dans la nomenclature nationale des activités).

\begin{definition}[Les trois secteurs d'activité]
\begin{itemize}
\item \cle{Secteur primaire} : activités qui exploitent directement la nature — agriculture, élevage, pêche, sylviculture, mines. Exemples : la CDC (bananeraies, palmeraies), la SONARA (raffinage de pétrole brut, à la frontière primaire/secondaire), les Plantations du Haut Penja (PHP, bananes/ananas).
\item \cle{Secteur secondaire} : activités de transformation des matières premières — industrie manufacturière, agroalimentaire, BTP, énergie. Exemples : cimenteries (CIMENCAM, Dangote Cement), brasseries (SABC), chocolateries, entreprises de construction.
\item \cle{Secteur tertiaire} : activités de services — commerce, transport, banques, assurances, télécommunications, éducation, santé, NTIC. Exemples : Jumia Cameroun (commerce électronique), les fintechs de mobile money, les agences de voyage, les hôpitaux privés.
\end{itemize}
\end{definition}

\begin{exemplebox}[Deux entreprises, deux logiques sectorielles]
Comparons la PHP (secteur primaire) et une start-up fintech « Djangui » (secteur tertiaire). La PHP exige une \textbf{forte intensité capitalistique} (terrains, usines, flotte logistique) et une main-d'œuvre nombreuse mais souvent saisonnière (récoltes). La fintech, elle, demande peu de capital matériel : son actif critique est le \textbf{capital humain} (développeurs, data analysts) et sa force est la \emph{scalabilité} : une fois l'application construite, servir 10 000 ou 1 million d'utilisateurs coûte presque le même prix. Le secteur détermine donc la structure des coûts, le type de main-d'œuvre et la vitesse de croissance possible.
\end{exemplebox}

\begin{savaistu}[La tertiarisation de l'économie camerounaise]
Comme toutes les économies qui se développent, le Cameroun connaît une \emph{tertiarisation} : le secteur tertiaire représente aujourd'hui environ 50 \% du PIB, devant le primaire et le secondaire. Mais attention : une part importante de ce tertiaire est constituée de petits services de survie (commerce de rue, transport informel), ce qui n'a pas la même signification économique que la tertiarisation des pays industrialisés, fondée sur des services à haute valeur ajoutée (finance, conseil, R\&D).
\end{savaistu}

\courssub{Le critère juridique : les formes sociales de l'OHADA}

\textbf{L'intuition.} Créer une entreprise, c'est aussi choisir un « vêtement juridique » : qui est responsable des dettes ? Sur quels biens ? Comment prend-on les décisions ? Le droit des sociétés applicable au Cameroun est celui de l'\cle{OHADA} (Organisation pour l'Harmonisation en Afrique du Droit des Affaires), dont l'Acte uniforme sur le droit des sociétés commerciales et du GIE (révisé en 2014) régit les formes sociales dans 17 États africains.

\begin{propriete}[Principe de liberté contractuelle (OHADA)]
Les associés peuvent organiser librement leur société par les statuts, sous réserve des dispositions d'ordre public impératives : capitaux minimaux légaux, organes obligatoires (assemblée générale, dirigeant), règles de publicité (immatriculation au RCCM). La liberté est donc large mais encadrée pour protéger les tiers et les associés minoritaires.
\end{propriete}

\textbf{Le panorama des formes juridiques.}

\begin{center}
\begin{tabularx}{\linewidth}{|l|X|X|}
\hline
\rowcolor{popblueL} \textbf{Forme} & \textbf{Caractéristiques essentielles} & \textbf{Profil type} \\
\hline
Entreprise individuelle & Une seule personne, pas de capital minimum, patrimoine confondu avec celui de l'exploitant (responsabilité illimitée, sauf option EIRL) & Commerçant, artisan, libéral \\
\hline
SNC (Société en Nom Collectif) & Associés répondent indéfiniment et solidairement des dettes ; parts non librement cessibles & Petite structure familiale fondée sur la confiance \\
\hline
SARL & 1 associé ou plus ; capital minimum 100 000 FCFA (réforme 2014) ; responsabilité limitée aux apports ; gérant nommé & PME familiale, « Bois \& Design » \\
\hline
SA (Société Anonyme) & Capital minimum 10 millions de FCFA ; actionnaires ; conseil d'administration ou administrateur général ; commissaire aux comptes obligatoire & Grandes entreprises, banques, filiales de groupes \\
\hline
SAS & Pas de capital minimum légal ; organisation statutaire très libre ; président obligatoire & Startups, filiales de groupes, joint-ventures \\
\hline
SCA (Société en Commandite par Actions) & Associe commandités (responsabilité illimitée, gèrent) et commanditaires (apporteurs, responsabilité limitée) & Structures hybrides famille/investisseurs \\
\hline
GIE & Groupement sans but lucratif propre : met en commun des moyens entre membres & Consortia d'entreprises, achats groupés \\
\hline
Coopérative & Gestion démocratique (1 personne = 1 voix), excédents répartis selon l'usage & Coopératives agricoles, d'épargne-crédit \\
\hline
\end{tabularx}
\end{center}

\begin{aretenir}[Les deux formes à maîtriser pour l'examen]
\cle{SARL} : capital minimum de 100 000 FCFA (depuis 2014), responsabilité limitée aux apports, adaptée aux PME. \cle{SA} : capital minimum de 10 millions de FCFA, actions librement cessibles, adaptée aux grands projets nécessitant des capitaux importants.
\end{aretenir}

\begin{savaistu}[La SAS, la forme « startup » créée par la réforme 2014]
La réforme OHADA de 2014 a introduit la \cle{SAS} (Société par Actions Simplifiée) : pas de capital minimum, organes librement organisés par les statuts, seul un président est obligatoire. Cette souplesse la rend très prisée des startups et des filiales de groupes internationaux, qui peuvent y prévoir librement pactes d'actionnaires et clauses de sortie.
\end{savaistu}

\textbf{Comment choisir ?} Le choix dépend de quatre questions : combien sommes-nous ? Quel capital pouvons-nous réunir ? Quel risque acceptons-nous de prendre sur notre patrimoine personnel ? Comment transmettra-t-on l'entreprise ? L'arbre de décision ci-dessous synthétise la démarche.

\begin{popfigure}
\begin{center}
\begin{tikzpicture}[
font=\small,
dec/.style={rectangle,draw=popblue,fill=popblueL,rounded corners,align=center,inner sep=3pt},
res/.style={rectangle,draw=poporange,fill=poporangeL,rounded corners,align=center,inner sep=3pt},
edge from parent/.style={draw=popmuted,-latex},
level distance=1.6cm,
level 1/.style={sibling distance=5.5cm},
level 2/.style={sibling distance=3.0cm},
]
\node[dec]{Seul ou plusieurs associés ?}
  child { node[dec]{Seul} 
    child { node[res, align=center]{Entreprise\\individuelle} edge from parent node[left,font=\scriptsize]{risque accepté} }
    child { node[res, align=center]{SARL unipersonnelle\\ou SASU} edge from parent node[right,font=\scriptsize]{patrimoine protégé} }
    edge from parent node[above left,font=\scriptsize]{1 personne} }
  child { node[dec]{Capital disponible ?}
    child { node[res, align=center]{SARL\\(min 100 000 F)} edge from parent node[left,font=\scriptsize]{faible/moyen} }
    child { node[dec]{Gouvernance ?}
      child { node[res, align=center]{SA\\(min 10 M F)} edge from parent node[left,font=\scriptsize]{conseil d'admin.} }
      child { node[res, align=center]{SAS\\(statuts libres)} edge from parent node[right,font=\scriptsize]{souplesse} }
      edge from parent node[right,font=\scriptsize]{élevé} }
    edge from parent node[above right,font=\scriptsize]{plusieurs} };
\end{tikzpicture}
\end{center}
\legende{Arbre de décision simplifié : « Quelle forme juridique choisir ? ». Critères successifs : nombre d'associés, montant du capital, niveau de responsabilité accepté, mode de gouvernance et de transmission.}
\end{popfigure}

\courssub{Le critère de la propriété et de l'origine du capital}

\textbf{La propriété du capital} distingue :

\begin{itemize}
\item Les entreprises \cle{privées} : le capital appartient à des personnes privées, nationales (ex. un groupe camerounais) ou étrangères (ex. filiales de multinationales).
\item Les entreprises \cle{publiques} : le capital est détenu majoritairement ou totalement par l'État. Ce sont souvent des \cle{EPIC} (Établissements Publics à caractère Industriel et Commercial) ou des sociétés à capital public (SA à capitaux publics) : CAMWATER (eau potable), CAMTEL (télécommunications), CAMRAIL (concession ferroviaire), ENEO (concession électrique, capital mixte).
\item Les entreprises \cle{mixtes} (Sociétés d'Économie Mixte, SEM) : le capital est partagé entre l'État (ou des collectivités) et des partenaires privés. Exemple : la SEM chargée de la gestion d'un marché ou d'un port, où la puissance publique garde un droit de regard stratégique.
\end{itemize}

\textbf{L'origine du capital} recoupe partiellement la question précédente mais répond à une autre interrogation : d'où vient l'argent investi ? On distingue le capital \cle{local} (épargne nationale, banques camerounaises, diaspora) et les \cle{IDE} (Investissements Directs Étrangers) : un investisseur étranger crée ou rachète durablement une entreprise au Cameroun. Exemples : Bolloré (logistique portuaire), Orange (télécommunications), Guinness (brasserie via Diageo/SABC). Les IDE apportent capitaux, technologies et emplois qualifiés, mais posent la question des rapatriements de bénéfices et de la dépendance économique — débat classique au programme de la réflexion sur le développement.

\begin{exoresolu}[Classer une entreprise réelle — Modèle de réponse attendue]
\textbf{Énoncé.} CAMWATER (Cameroon Water Utilities) est une société de droit camerounien, au capital entièrement détenu par l'État, qui emploie plus de 1 000 salariés et produit et distribue l'eau potable. Classer CAMWATER selon les critères de taille, de secteur, de forme et de propriété du capital, en justifiant chaque réponse.
\tcblower
\textbf{Solution détaillée.}
\begin{enumerate}
\item \textbf{Taille} : avec plus de 1 000 salariés permanents, CAMWATER dépasse le seuil de 200 salariés du Décret n°2010/0882/PM : c'est une \textbf{Grande Entreprise} (condition effectif remplie).
\item \textbf{Secteur} : la production et la distribution d'eau potable relèvent des services aux collectivités : \textbf{secteur tertiaire} (certains classements parlent d'« industrie des réseaux », mais la prestation délivrée est un service).
\item \textbf{Forme juridique} : société de droit camerounien à capital public, structurée en société commerciale (SA à capitaux publics) relevant du droit OHADA.
\item \textbf{Propriété du capital} : capital détenu à 100 \% par l'État : entreprise \textbf{publique}.
\end{enumerate}
\textbf{Formule de synthèse attendue} : « CAMWATER est une grande entreprise publique du secteur tertiaire, de forme sociétaire de droit camerounais (SA à capitaux publics). » On voit ici qu'une seule entreprise mobilise quatre critères à la fois.
\end{exoresolu}

\courssub{Le secteur informel : l'entreprise hors des cadres}

\textbf{L'observation.} À Douala ou à Yaoundé, le cordonnier du carrefour, la couturière à domicile, le « bayam-sellam » du marché produisent bien des biens et des services : ce sont économiquement des entreprises. Mais elles échappent aux cadres juridiques et statistiques.

\begin{definition}[Secteur informel]
Une unité de production \cle{informelle} est une unité non enregistrée auprès des administrations : pas d'immatriculation au Registre du Commerce et du Crédit Mobilier (RCCM), pas de numéro de contribuable auprès de la Direction Générale des Impôts, pas d'affiliation des travailleurs à la CNPS. Elle fonctionne donc en dehors du droit formel, sans pour autant être nécessairement illicite dans son activité (vente de beignets, couture, réparation).
\end{definition}

\textbf{Les enjeux sont ambivalents.} D'un côté, l'informel est un \emph{problème} : précarité des travailleurs (pas de contrat, pas de couverture sociale, pas de retraite), échappement fiscal qui prive l'État de ressources pour les écoles et les hôpitaux, concurrence parfois déloyale envers les entreprises formelles (qui paient charges et impôts), accès impossible au crédit bancaire. De l'autre, c'est une \emph{réalité structurante} : le secteur informel fournit environ 90 \% de l'emploi urbain non agricole au Cameroun ; il est le vivier d'emploi des jeunes et des non-diplômés, et beaucoup d'entreprises formelles d'aujourd'hui ont commencé dans l'informel. La question philosophique et politique est donc : faut-il réprimer l'informel ou l'accompagner vers la formalisation ? Les politiques récentes (statut simplifié d'entreprenant, centres de formalisation des entreprises) parient sur la seconde voie.

\begin{methode}[Classer une entreprise dans un dossier d'examen — Méthode en 5 étapes]
\begin{enumerate}
\item Identifier l'entreprise et relever les informations du document (effectif, CA, activité, capital, statut).
\item Appliquer \textbf{au moins deux critères} croisés : taille + secteur + juridique + propriété.
\item Rédiger une formule complète : par exemple « PME industrielle privée de droit camerounais » ou « GE publique du secteur tertiaire ».
\item Justifier \textbf{chaque} choix par un seuil ou une définition du cours (Décret 2010, Acte uniforme OHADA).
\item Ne jamais se limiter à la taille : une classification à un seul critère est une copie incomplète.
\end{enumerate}
\end{methode}

\begin{attention}[Pièges fréquents — Trois erreurs à ne pas commettre]
Trois erreurs reviennent souvent : (1) confondre taille et forme juridique (voir l'encadré plus haut) ; (2) croire que « informel » signifie « illégal » — vendre des beignets n'est pas un délit, c'est l'absence d'enregistrement qui définit l'informalité ; (3) oublier que les seuils de taille combinent effectif \emph{et} chiffre d'affaires, et qu'une \textbf{seule} des conditions peut suffire à classer une entreprise en GE (effectif $\ge$ 200 \textbf{ou} CA $>$ 10 Mds).
\end{attention}

\textbf{Conclusion de la section.} Classer les entreprises, ce n'est pas un exercice scolaire gratuit : c'est l'outil qui permet à l'État de fixer des seuils d'aide, aux banques d'évaluer les risques, aux statisticiens de photographier l'économie et à l'entrepreneur de choisir sa structure. Nous savons désormais \emph{ce qu'est} une entreprise et \emph{comment la classer}. Reste à ouvrir la boîte noire : comment une entreprise fonctionne-t-elle de l'intérieur ? C'est l'objet de la section suivante, consacrée aux facteurs de production et à l'organisation.

\coursec{Le fonctionnement interne : Facteurs de production et Organisation}

Dans la section précédente, nous avons classé les entreprises camerounaises selon leur taille, leur secteur et leur statut juridique. Mais qu'il s'agisse d'une micro-entreprise de quartier à Douala ou d'une grande société comme la SABC, toutes partagent une même question : \emph{comment transformer des ressources en biens et services ?} Autrement dit, comment l'entreprise \cle{fonctionne}-t-elle de l'intérieur ? C'est l'objet de cette section : comprendre d'abord \emph{avec quoi} l'entreprise produit (les facteurs de production), puis \emph{comment} elle organise ces ressources (structure, flux, pilotage).

\courssub{Les facteurs de production : les ingrédients de l'entreprise}

\begin{definition}[Facteur de production]
Un \cle{facteur de production} est une ressource utilisée par l'entreprise pour produire des biens ou des services. La théorie économique classique en distingue trois principaux : le \cle{travail}, le \cle{capital} et l'\cle{organisation} (aussi désignée comme facteur entrepreneurial ou direction d'entreprise).
\end{definition}

\textbf{Le travail (facteur humain).} C'est l'ensemble des efforts physiques et intellectuels fournis par les salariés. Sa valeur dépend de deux éléments : la \cle{qualification} (un soudeur certifié coûte plus cher qu'une manœuvre non qualifiée, mais produit davantage et mieux) et le \cle{coût} (salaire brut + charges sociales patronales). Au Cameroun, le paradoxe est connu : le pays manque de main-d'œuvre qualifiée dans certains métiers techniques (froid industriel, maintenance, usinage) alors que le chômage des jeunes reste élevé — c'est le problème de l'\emph{adéquation formation-emploi}, que l'enseignement technique cherche précisément à résoudre.

\textbf{Le capital.} On distingue :
\begin{itemize}
\item le \cle{capital fixe} : les biens durables qui servent plusieurs cycles de production (machines, bâtiments, camions) ;
\item le \cle{capital circulant} : les biens consommés ou transformés à chaque cycle (matières premières, stock de bois pour un menuisier) ;
\item le \cle{capital matériel} (machines, véhicules) et le \cle{capital immatériel} : brevets, logiciels, marque, réputation. La marque « Top » ou « 33 Export » vaut parfois plus que les usines qui la produisent !
\end{itemize}

\textbf{L'organisation.} C'est le facteur souvent oublié, et pourtant décisif : deux ateliers avec les mêmes machines et les mêmes ouvriers n'ont pas les mêmes résultats. L'organisation, c'est le \emph{savoir-faire managérial} : la manière de répartir les tâches, de faire circuler l'information, de décider.

\begin{definition}[L'organisation]
L'\cle{organisation} est l'agencement stable des ressources (humaines, matérielles, informationnelles) en une structure cohérente pour atteindre les objectifs de l'entreprise.
\end{definition}

\courssub{La fonction de production et la productivité}

On peut résumer l'activité productive par une relation simple : la quantité produite $Q$ dépend des quantités de travail $T$, de capital $K$ et de la qualité de l'organisation $O$ :
\[ Q = f(T, K, O) \]
Cette écriture signifie simplement : « la production est une fonction des facteurs employés ». À organisation égale, plus on emploie de travail et de capital, plus on produit — mais pas indéfiniment, et pas toujours proportionnellement.

\begin{definition}[Productivité]
La \cle{productivité} mesure l'efficacité d'un facteur : c'est le rapport entre la production obtenue et la quantité de facteur utilisée.
\[ \text{Productivité du travail} = \frac{Q}{T} \qquad ; \qquad \text{Productivité du capital} = \frac{Q}{K} \]
\end{definition}

\begin{exemplebox}[Productivité dans une menuiserie de Yaoundé]
L'atelier « Bois \& Design » emploie 5 menuisiers qui produisent 100 chaises par mois. La productivité du travail est $Q/T = 100/5 = 20$ chaises par menuisier et par mois. Après l'achat d'une scie numérique (augmentation de $K$), les mêmes 5 menuisiers produisent 150 chaises : la productivité passe à 30 chaises par menuisier. Le capital a \emph{augmenté l'efficacité du travail}.
\end{exemplebox}

\begin{propriete}[Rendements d'échelle]
Lorsqu'on augmente \emph{tous} les facteurs de production dans la même proportion, la production peut varier de trois façons :
\begin{itemize}
\item \cle{rendements croissants} : doubler $T$ et $K$ fait plus que doubler $Q$ (gains d'échelle, spécialisation accrue) ;
\item \cle{rendements constants} : doubler $T$ et $K$ double exactement $Q$ ;
\item \cle{rendements décroissants} : doubler $T$ et $K$ fait moins que doubler $Q$ (l'atelier devient trop grand, la coordination se dégrade, les files d'attente apparaissent).
\end{itemize}
\end{propriete}

C'est pourquoi « grandir » n'est pas toujours rentable : au-delà d'une certaine taille, une PME mal organisée perd en efficacité ce qu'elle gagne en volume.

\courssub{L'organisation du travail : divisions verticale et horizontale}

Pour organiser, il faut d'abord \emph{diviser} le travail, puis \emph{coordonner} ce qui a été divisé. On distingue deux divisions complémentaires.

\textbf{La division verticale (hiérarchique).} Elle répartit le pouvoir de décision en trois niveaux :
\begin{itemize}
\item la \cle{direction} : fixe les objectifs et la stratégie (le « pourquoi » et le « quoi ») ;
\item l'\cle{encadrement} (cadres intermédiaires) : traduit les objectifs en plans et contrôle leur exécution (le « comment ») ;
\item l'\cle{exécution} : réalise concrètement les tâches (le « faire »).
\end{itemize}

\textbf{La division horizontale (fonctionnelle).} À un même niveau hiérarchique, on répartit les activités en \cle{services} spécialisés : Production, Commercial, Finance-Comptabilité, Ressources Humaines (RH), Logistique, Systèmes d'Information (SI).

\begin{propriete}[Les principes classiques de l'organisation — Taylor et Fayol]
Au début du \textsc{xx}\textsuperscript{e} siècle, Frederick \textbf{Taylor} (organisation scientifique du travail) et Henri \textbf{Fayol} (administration industrielle) ont posé des principes encore enseignés :
\begin{itemize}
\item la \cle{division du travail} : la spécialisation augmente la productivité ;
\item l'\cle{unité de commandement} : un salarié ne reçoit d'ordres que d'un seul chef ;
\item la \cle{hiérarchie} : une chaîne claire de commandement, du sommet à la base ;
\item l'\cle{étendue de contrôle} : un chef ne peut superviser efficacement qu'un nombre limité de subordonnés.
\end{itemize}
Ces principes sont aujourd'hui nuancés par les modèles modernes : organisation \emph{matricielle} (un salarié peut dépendre d'un chef de service \emph{et} d'un chef de projet) et organisation \emph{agile} (petites équipes autonomes, hiérarchie allégée). L'essentiel demeure : qui décide, qui fait, qui contrôle ?
\end{propriete}

\begin{methode}[Construire un organigramme]
\begin{enumerate}
\item \textbf{Lister} toutes les activités de l'entreprise (produire, vendre, payer, recruter, livrer...).
\item \textbf{Regrouper} ces activités en services homogènes (par fonction, par produit, par région ou par type de client).
\item \textbf{Définir les liens} : liens hiérarchiques (qui commande qui, le « N+1 ») et liens fonctionnels (coopération sans autorité directe).
\item \textbf{Dessiner} : un rectangle par service, un trait plein pour la hiérarchie, des pointillés pour les relations fonctionnelles.
\end{enumerate}
\end{methode}

\begin{popfigure}
\begin{center}
\begin{tikzpicture}[
  node distance=0.55cm and 0.35cm,
  dir/.style={rectangle, draw=popdark, fill=popblue!25, rounded corners=2pt, minimum width=3.4cm, minimum height=0.8cm, align=center, font=\small\bfseries},
  serv/.style={rectangle, draw=popdark, fill=popcream, rounded corners=2pt, minimum width=2.5cm, minimum height=0.75cm, align=center, font=\footnotesize},
  hier/.style={-{Stealth[length=2mm]}, thick, popdark},
  fonc/.style={dashed, poppurple, thick}
]
\node[dir] (dg) {Direction Générale};
\node[serv, below left=1.1cm and 2.6cm of dg, align=center] (daf){DAF\\ (Finances)};
\node[serv, right=0.3cm of daf, align=center] (drh){DRH\\ (Personnel)};
\node[serv, right=0.3cm of drh, align=center] (dcom){D. Commercial\\ (Ventes)};
\node[serv, right=0.3cm of dcom, align=center] (dprod){D. Production\\ (Atelier)};
\node[serv, right=0.3cm of dprod, align=center] (dsi){DSI\\ (Informatique)};
\draw[hier] (dg.south) -- ++(0,-0.35) -| (daf.north);
\draw[hier] (dg.south) -- ++(0,-0.35) -| (drh.north);
\draw[hier] (dg.south) -- ++(0,-0.35) -| (dcom.north);
\draw[hier] (dg.south) -- ++(0,-0.35) -| (dprod.north);
\draw[hier] (dg.south) -- ++(0,-0.35) -| (dsi.north);
\draw[fonc] (dcom.south) .. controls +(0,-0.5) and +(0,-0.5) .. (dprod.south);
\draw[fonc] (drh.south) .. controls +(0,-0.9) and +(0,-0.9) .. (dsi.south);
\end{tikzpicture}
\end{center}
\legende{Organigramme fonctionnel simplifié d'une PME : les traits pleins indiquent les liens hiérarchiques (autorité directe), les pointillés violets les liens fonctionnels (coopération : le Commercial transmet les commandes à la Production ; la DRH et la DSI coordonnent la formation aux outils numériques).}
\end{popfigure}

\begin{attention}[Organigramme $\neq$ processus]
Ne pas confondre l'\cle{organigramme} (la structure \emph{statique} : qui fait quoi, qui commande qui) et le \cle{processus} (la dynamique : comment circulent les commandes, les matières, l'argent). Un bel organigramme ne garantit rien si les processus sont bloqués. Une entreprise efficace \emph{aligne} sa structure et ses processus.
\end{attention}

\courssub{Les flux dans l'entreprise}

L'entreprise est traversée en permanence par trois types de \cle{flux} :
\begin{itemize}
\item les \cle{flux physiques} : matières premières $\rightarrow$ production $\rightarrow$ produits finis $\rightarrow$ livraison au client ;
\item les \cle{flux financiers} : encaissements des ventes, paiements des fournisseurs et des salaires, investissements (achat de machines) — leur gestion s'appelle la \emph{trésorerie} ;
\item les \cle{flux d'information} : ordres de fabrication, fiches de paie, reporting vers la direction, données clients. Sans information fiable et rapide, aucune décision correcte n'est possible.
\end{itemize}

\begin{popfigure}
\begin{center}
\begin{tikzpicture}[
  node distance=0.9cm,
  etape/.style={rectangle, draw=popdark, fill=popblueL, rounded corners=3pt, minimum width=2.6cm, minimum height=0.9cm, align=center, font=\small},
  flux/.style={-{Stealth[length=2.5mm]}, very thick, poporange},
  info/.style={-{Stealth[length=2mm]}, dashed, poppurple}
]
\node[etape, align=center] (c){Commande\\ client};
\node[etape, right=of c, align=center] (l){Livraison\\ (flux physique)};
\node[etape, right=of l] (f) {Facturation};
\node[etape, right=of f, align=center] (e){Encaissement\\ (flux financier)};
\draw[flux] (c) -- (l);
\draw[flux] (l) -- (f);
\draw[flux] (f) -- (e);
\draw[info] (e.south) .. controls +(0,-0.9) and +(0,-0.9) .. node[below, font=\footnotesize, poppurple]{information : solde du compte client} (c.south);
\end{tikzpicture}
\end{center}
\legende{Le cycle « Order-to-Cash » (de la commande à l'encaissement) : les flèches orange représentent le flux principal (physique puis financier), la flèche violette en pointillés le flux d'information qui boucle le processus.}
\end{popfigure}

\courssub{Les outils de pilotage}

Piloter une entreprise, c'est comme conduire un taxi à Douala : il faut un tableau de bord (vitesse, essence) et pas seulement regarder par la fenêtre.

\begin{itemize}
\item le \cle{tableau de bord} : synthèse régulière (mensuelle) des chiffres clés présentée à la direction ;
\item les \cle{indicateurs} (KPI, \emph{Key Performance Indicators}) : grandeurs suivies pour mesurer la performance — chiffre d'affaires, taux de défauts, délai moyen de livraison, taux d'absentéisme ;
\item le \cle{budget} : prévision chiffrée des recettes et dépenses, qui sert de feuille de route et de référence pour comparer le prévu et le réalisé ;
\item la \cle{comptabilité analytique} : elle calcule les coûts par produit ou par activité. On distingue les \emph{coûts variables} (qui varient avec la production : bois, clous, électricité des machines) et les \emph{coûts fixes} (loyer, salaires administratifs). Le \emph{coût complet} d'un produit = coûts variables + part des coûts fixes.
\end{itemize}

\begin{exoresolu}[Coût complet d'une chaise chez « Bois \& Design »]
L'atelier produit 100 chaises par mois. Coûts variables par chaise : bois et quincaillerie 4\,000 FCFA. Coûts fixes mensuels : loyer 100\,000 FCFA, salaires 250\,000 FCFA, énergie 50\,000 FCFA. (1) Calculer le coût complet d'une chaise. (2) L'atelier vend la chaise 8\,000 FCFA : est-ce rentable ?
\tcblower
(1) Coûts fixes totaux : $100\,000 + 250\,000 + 50\,000 = 400\,000$ FCFA, soit $400\,000 / 100 = 4\,000$ FCFA par chaise.
Coût complet unitaire : $4\,000 + 4\,000 = 8\,000$ FCFA.
(2) À 8\,000 FCFA, l'atelier couvre exactement ses coûts : résultat nul (seuil de rentabilité atteint). Pour dégager un profit, il faut soit vendre plus cher, soit produire davantage (à 150 chaises, la part de fixe tombe à $400\,000/150 \approx 2\,667$ FCFA et le coût complet à $\approx 6\,667$ FCFA, dégageant $\approx 1\,333$ FCFA de marge par chaise). On retrouve l'intérêt des rendements d'échelle.
\end{exoresolu}

\begin{exemplebox}[Le coût du travail au Cameroun]
Le SMIG est d'environ 41\,875 FCFA mensuels (revalorisation 2023). Mais pour l'employeur, le coût réel est supérieur : les charges patronales (CNPS, Crédit Foncier, taxe d'apprentissage et formation professionnelle) représentent environ 20 à 25\,\% du salaire brut. Un salarié payé 50\,000 FCFA brut coûte donc environ 60\,000 à 62\,500 FCFA à l'entreprise. S'y ajoutent le coût élevé de l'énergie (facteur capital : groupes électrogènes face aux délestages) et les lourdeurs administratives (procédures de création, fiscalité), qui pèsent surtout sur les PME.
\end{exemplebox}

\begin{savaistu}[Le Lean Management chez SOSUCAM]
Le \emph{Lean Management}, né chez Toyota au Japon, consiste à éliminer les gaspillages (les \emph{Muda}) : surproduction, attentes, transports inutiles, stocks excessifs, défauts. Au Cameroun, des agro-industries comme la SOSUCAM (Société Sucrière du Cameroun) appliquent ces méthodes pour fluidifier la récolte et l'usinage de la canne : moins d'attente entre coupe et broyage, c'est plus de sucre récupéré. Preuve que les outils d'organisation ne sont pas un luxe de pays riches !
\end{savaistu}

\begin{aretenir}[L'essentiel de la section]
\begin{itemize}
\item Trois facteurs de production : \cle{travail}, \cle{capital} (fixe/circulant, matériel/immatériel), \cle{organisation} ; $Q = f(T, K, O)$.
\item Productivité = $Q/T$ ou $Q/K$ ; les rendements d'échelle peuvent être croissants, constants ou décroissants.
\item Division verticale (direction, encadrement, exécution) et horizontale (services fonctionnels), formalisées par l'organigramme.
\item Trois flux : physiques, financiers, informationnels ; le pilotage s'appuie sur tableau de bord, indicateurs, budget et comptabilité analytique.
\item Au Cameroun : pénurie de main-d'œuvre qualifiée, coût de l'énergie et lourdeurs administratives pèsent sur le fonctionnement des entreprises.
\end{itemize}
\end{aretenir}

\begin{exercice}[Application]
Une imprimerie de Bafoussam emploie 8 ouvriers et produit 12\,000 cahiers par mois. Elle envisage d'acheter une presse numérique qui permettrait de produire 20\,000 cahiers avec 6 ouvriers. (1) Calculer la productivité du travail avant et après investissement. (2) Identifier deux coûts fixes et deux coûts variables de cette imprimerie. (3) Expliquer pourquoi cette décision relève à la fois du facteur capital et du facteur organisation.
\end{exercice}

\begin{corrige}[Exercice 1]
(1) Avant : $12\,000/8 = 1\,500$ cahiers par ouvrier. Après : $20\,000/6 \approx 3\,333$ cahiers par ouvrier, soit une productivité plus que doublée.
(2) Coûts fixes : loyer de l'atelier, salaire du comptable, amortissement de la presse. Coûts variables : papier, encre, électricité de fonctionnement.
(3) Capital : achat d'un bien d'équipement durable (capital fixe matériel). Organisation : la nouvelle machine impose de redéfinir les tâches, de former les opérateurs et de réorganiser le flux de production — un investissement sans réorganisation ne donne pas les gains espérés.
\end{corrige}

\coursec{Gouvernance, Stratégie et Responsabilité Sociétale (RSE)}

Après avoir analysé le fonctionnement interne de l'entreprise — facteurs de production, organisation des services et processus opérationnels — il convient d'examiner comment l'entreprise est \emph{dirigée} et \emph{contrôlée}, comment elle définit sa \emph{stratégie} à long terme et comment elle intègre les enjeux \emph{sociétaux} et \emph{environnementaux} dans son modèle d'affaires. Cette section articule donc gouvernance, stratégie et RSE, en mobilisant des exemples camerounais concrets.

\courssub{Organes de gouvernance : Assemblée Générale, Conseil d'Administration, Direction Générale, Commissaires aux Comptes}

\begin{definition}[Gouvernance d'entreprise]
La gouvernance d'entreprise est l'ensemble des règles, processus et pratiques par lesquels une entreprise est dirigée et contrôlée, définissant la répartition des droits et responsabilités entre les parties prenantes (actionnaires, dirigeants, salariés, créanciers, État, communautés locales).
\end{definition}

Dans le droit OHADA (Acte uniforme sur les sociétés commerciales), la structure de gouvernance varie selon la forme juridique. Pour une \textbf{Société Anonyme (SA)}, quatre organes sont obligatoires. Pour une \textbf{Société par Actions Simplifiée (SAS)}, la structure est plus souple : un Président unique est obligatoire, le Conseil d'Administration est facultatif, et le Commissaire aux Comptes n'est requis que si des seuils sont franchis.

\begin{itemize}
    \item \textbf{Assemblée Générale (AG)} : organe souverain réunissant les actionnaires/associés. Elle nomme et révoque les administrateurs (SA) ou le Président (SAS), approuve les comptes, décide des augmentations de capital, de la dissolution, etc. Principe : « une action = une voix ».
    \item \textbf{Conseil d'Administration (CA) / Président (SAS)} : dans la SA, le CA (3 à 12 membres) définit la stratégie, nomme le Directeur Général (DG), contrôle la gestion. Dans la SAS, le Président (personne physique ou morale) exerce ces pouvoirs, éventuellement assisté d'organes consultatifs.
    \item \textbf{Direction Générale (DG)} : assure l'exécution opérationnelle de la stratégie, manage les ressources, rend des comptes au CA (SA) ou au Président (SAS).
    \item \textbf{Commissaires aux Comptes (CAC)} : organe de contrôle légal et financier, vérifie la régularité et la sincérité des comptes, certifie les états financiers. Obligatoire en SA ; en SAS, obligatoire si 2 seuils sur 3 sont dépassés (50 salariés, 125 millions FCFA de CA, 62,5 millions FCFA de total bilan).
\end{itemize}

\begin{popfigure}
\begin{tikzpicture}[
    node distance=2.5cm and 3cm,
    box/.style={draw=popblue, thick, rounded corners, fill=popblueL, minimum width=2.8cm, minimum height=1.2cm, align=center},
    arrow/.style={-Stealth, thick, popdark},
    labelarrow/.style={midway, font=\small, text=popdark, fill=popcream, inner sep=2pt, rounded corners}
]
\node[box, align=center] (actionnaires){Actionnaires / Associés\\Assemblée Générale};
\node[box, right=of actionnaires, align=center] (ca){Conseil d'Administration\\ (SA) / Président (SAS)};
\node[box, below=of ca, align=center] (dg){Direction Générale\\ / DG};
\draw[arrow] (actionnaires) -- node[labelarrow] {Nomination / Révocation} (ca);
\draw[arrow] (ca) -- node[labelarrow] {Stratégie / Contrôle} (dg);
\draw[arrow] (dg) -- node[labelarrow] {Redevabilité / Reporting} (ca);
\draw[arrow] (ca) -- node[labelarrow] {Information / Comptes} (actionnaires);
% Commissaires aux comptes
\node[box, left=of dg, xshift=-1.5cm, align=center] (cac){Commissaires\\aux Comptes};
\draw[arrow, dashed] (cac) -- node[labelarrow, rotate=90] {Contrôle légalité} (dg);
\draw[arrow, dashed] (cac) -- node[labelarrow, rotate=-90] {Certification comptes} (actionnaires);
\end{tikzpicture}
\legende{Schéma simplifié de la gouvernance SA (organes obligatoires) / SAS (CA facultatif, CAC selon seuils). Flèches pleines : relations de pouvoir et de reddition de comptes ; flèches pointillées : contrôle externe des Commissaires aux Comptes.}
\end{popfigure}

\courssub{Conflits d'agence et mécanismes d'alignement (Jensen \& Meckling)}

La séparation entre \emph{propriete} (actionnaires) et \emph{contrôle} (dirigeants) engendre des \textbf{conflits d'agence} : les intérêts des parties ne coïncident pas spontanément.

\begin{itemize}
    \item \textbf{Actionnaires} : maximisation de la valeur de l'action et des dividendes.
    \item \textbf{Dirigeants} : sécurité d'emploi, rémunération, pouvoir, prestige, aversion au risque excessif.
    \item \textbf{Salariés} : salaires, conditions de travail, formation, stabilité.
\end{itemize}

\emph{Michael Jensen et William Meckling (1976)} modélisent ces conflits comme des coûts d'agence : coûts de surveillance, coûts de garantie (bonding) et pertes résiduelles. Pour les réduire, plusieurs \textbf{mécanismes d'alignement} sont mobilisés :

\begin{itemize}
    \item \textbf{Incitations financières} : stock-options, actions gratuites, bonus indexés sur la performance boursière ou les résultats (EBITDA, ROE).
    \item \textbf{Contrôle interne} : comité d'audit, comité des rémunérations, procédures de compliance.
    \item \textbf{Marché du contrôle des sociétés} : menace d'OPA hostile, vote des actionnaires institutionnels.
    \item \textbf{Gouvernance renforcée} : indépendance des administrateurs, séparation Président / DG, limitation des mandats croisés.
\end{itemize}

\begin{attention}[Spécificité camerounaise : PME familiales]
Au Cameroun, la distinction « Actionnaire / Dirigeant » est souvent floue dans les PME familiales : le fondateur est simultanément PDG, actionnaire majoritaire et parfois unique administrateur. Le conflit d'agence classique se déplace alors vers :
\begin{itemize}
    \item \textbf{Famille / Managers extérieurs} : risque d'expropriation des managers non familiaux.
    \item \textbf{Actionnaires majoritaires / minoritaires} : décisions unilatérales (distribution de dividendes, rémunérations excessives) au détriment des minoritaires.
\end{itemize}
La professionnalisation de la gouvernance (conseil consultatif, CAC indépendant, charte de gouvernance familiale) devient alors cruciale pour la pérennité et l'accès au financement.
\end{attention}

\courssub{Stratégie d'entreprise : Analyse SWOT et choix stratégiques}

La stratégie est l'art d'allouer les ressources pour atteindre des objectifs durables dans un environnement incertain. L'outil pédagogique de base est l'\textbf{analyse SWOT} (Strengths, Weaknesses, Opportunities, Threats).

\begin{methode}[Analyse SWOT — 3 étapes]
\begin{enumerate}
    \item \textbf{Diagnostiquer l'interne} : identifier les \textbf{Forces} (ressources rares, compétences clés, marque, trésorerie) et les \textbf{Faiblesses} (dette, technologie obsolète, turnover, dépendance à un client).
    \item \textbf{Analyser l'externe} : repérer les \textbf{Opportunités} (croissance du marché, évolution réglementaire favorable, innovation technologique) et les \textbf{Menaces} (concurrence agressive, instabilité politique, hausse des matières premières). On utilise souvent le cadre \textbf{PESTEL} (Politique, Économique, Sociologique, Technologique, Écologique, Légal) pour structurer cette veille.
    \item \textbf{Croiser les quatre quadrants} pour définir des axes stratégiques :
    \begin{itemize}
        \item \textbf{FO (Forces–Opportunités)} : stratégies offensives (exploiter une force pour saisir une opportunité).
        \item \textbf{FA (Forces–Menaces)} : stratégies de protection (utiliser une force pour parer une menace).
        \item \textbf{MO (Faiblesses–Opportunités)} : stratégies de rattrapage (corriger une faiblesse pour profiter d'une opportunité).
        \item \textbf{MA (Faiblesses–Menaces)} : stratégies de survie / désengagement (réduire l'exposition).
    \end{itemize}
\end{enumerate}
\end{methode}

\begin{popfigure}
\begin{tikzpicture}[
    cell/.style={draw=popblue, thick, minimum width=7.5cm, minimum height=6cm, fill=popcream, align=left, text width=7cm, font=\small},
    header/.style={fill=popblue, text=white, font=\bfseries, minimum height=1cm, align=center},
    subheader/.style={fill=popblueL, font=\bfseries\small, minimum height=0.7cm, align=center},
]
\matrix (swot) [matrix of nodes, nodes in empty cells, column sep=-\pgflinewidth, row sep=-\pgflinewidth,
    column 1/.style={nodes={cell}},
    column 2/.style={nodes={cell}},
    row 1/.style={nodes={header}},
    row 2/.style={nodes={subheader}},
    row 3/.style={nodes={cell}},
    row 4/.style={nodes={cell}}
]{
    \textbf{Matrice SWOT — AgroTransfo SA (Ngaoundéré)} & \\
    \textbf{INTERNE} & \textbf{EXTERNE} \\
    {\textbf{Forces (F)} :\newline
    • Maîtrise transformation mangue (procédé breveté)\newline
    • Réseau 120 planteurs contractuels (fidélisés)\newline
    • Marque « Jus du Sahel » reconnue localement\newline
    • Trésorerie positive (autofinancement 65\%)} &
    {\textbf{Opportunités (O)} :\newline
    • Demande croissante jus naturel (classe moyenne urbaine)\newline
    • ZLECAf : accès marché CEMAC + Nigeria sans droits\newline
    • Programme « Agro-Industrie » MINADER (subventions équipement)\newline
    • Digitalisation distribution (e-commerce, livraison)} \\
    {\textbf{Faiblesses (W)} :\newline
    • Capacité usine saturée (85\%)\newline
    • Dépendance unique fruit (mangue, saison 4 mois)\newline
    • Logistique routière aléatoire (pistes en saison des pluies)\newline
    • Absence certification HACCP / ISO 22000} &
    {\textbf{Menaces (T)} :\newline
    • Concurrence importés (jus concentrés Chine, Europe)\newline
    • Volatilité prix mangue (aléa climatique)\newline
    • Réglementation douanière changeante (CEMAC)\newline
    • Pénurie main-d'œuvre qualifiée (techniciens agro)} \\
};
\end{tikzpicture}
\legende{Matrice SWOT pré-remplie pour la PME fictive « AgroTransfo SA » (transformation mangue → jus, Ngaoundéré).}
\end{popfigure}

À partir de cette matrice, AgroTransfo SA peut envisager :

\begin{itemize}
    \item \textbf{FO} : Lancer une gamme « Jus du Sahel Bio » certifiée HACCP pour conquérir le segment premium urbain et exporter vers le Nigeria via la ZLECAf.
    \item \textbf{FA} : Investir dans un séchoir solaire et une unité de purée de mangue (produit stable) pour réduire la dépendance saisonnière (Force : procédé breveté ; Menace : saisonnalité).
    \item \textbf{MO} : Demander une subvention MINADER pour acquérir une ligne de conditionnement aseptique (Faiblesse : pas de certification ; Opportunité : programme subventions).
    \item \textbf{MA} : Diversifier vers la transformation de l'ananas et de la papaye (réduire dépendance mangue) et former 10 techniciens en maintenance (Faiblesse + Menace).
\end{itemize}

\begin{aretenir}[Choix stratégiques classiques]
\begin{itemize}
    \item \textbf{Spécialisation} : concentrer les ressources sur un cœur de métier (ex. : uniquement jus de mangue). Avantages : économies d'échelle, expertise. Risque : dépendance marché/produit.
    \item \textbf{Diversification} : développer de nouveaux produits/marchés (ex. : jus ananas, confitures, cosmétiques à base de mangue). Réduit le risque mais dilue les ressources.
    \item \textbf{Intégration amont} : acquérir des plantations ou contractualiser plus de planteurs (maîtrise approvisionnement, qualité, coûts).
    \item \textbf{Intégration aval} : créer son propre réseau de distribution (kiosques, camions frigorifiques, plateforme e-commerce) pour capter la marge de distribution.
    \item \textbf{Internationalisation} : exporter vers la CEMAC (libre circulation marchandises) et la ZLECAf (1,3 milliard de consommateurs). Nécessite conformité normes (CODEX, ISO), logistique, veille réglementaire.
\end{itemize}
\end{aretenir}

\begin{exemplebox}[Exemple chiffré : ZLECAf]
La \textbf{Zone de Libre-Échange Continentale Africaine (ZLECAf)} entrée en vigueur en 2021 crée un marché unique de \textbf{1,3 milliard de consommateurs} et un PIB combiné d'environ \textbf{3 400 milliards USD}. Pour une entreprise camerounaise :
\begin{itemize}
    \item \textbf{Opportunité} : exporter du cacao transformé (chocolat, beurre de cacao), du coton filé, des services numériques (FinTech, e-learning) sans droits de douane intra-africains.
    \item \textbf{Menace} : concurrence accrue des géants nigérians (Dangote, BUA), égyptiens (Elsewedy) ou sud-africains (Shoprite, MTN) sur le marché camerounais (prix, échelle, marketing).
\end{itemize}
Une PME comme AgroTransfo SA doit donc \emph{certifier} ses produits (HACCP, ISO 22000) et \emph{numériser} sa chaîne logistique pour être compétitive.
\end{exemplebox}

\courssub{Responsabilité Sociétale des Entreprises (RSE) : ISO 26000 et piliers}

La RSE ne se réduit pas à la philanthropie (dons, parrainages). Elle désigne l'\textbf{intégration volontaire} des préoccupations sociales, environnementales et éthiques dans les activités commerciales et les relations avec les \emph{parties prenantes} (stakeholders).

\begin{propriete}[RSE — Définition ISO 26000]
La Responsabilité Sociétale des Entreprises est la responsabilité d'une organisation vis-à-vis des impacts de ses décisions et activités sur la société et l'environnement, se traduisant par un comportement éthique et transparent qui :
\begin{itemize}
    \item contribue au développement durable, y compris la santé et le bien-être de la société ;
    \item prend en compte les attentes des parties prenantes ;
    \item respecte les lois en vigueur et est cohérente avec les normes internationales de comportement ;
    \item est intégrée dans l'ensemble de l'organisation et mise en œuvre dans ses relations.
\end{itemize}
\end{propriete}

La norme \textbf{ISO 26000 (2010)} structure la RSE autour de \textbf{7 piliers (sujets centraux)} :

\begin{center}
\begin{tabularx}{\linewidth}{|l|X|}\hline
\textbf{Pilier} & \textbf{Principaux champs d'action} \\\hline
1. Gouvernance de l'organisation & Transparence, éthique, lutte contre la corruption, dialogue parties prenantes \\\hline
2. Droits de l'homme & Non-discrimination, travail forcé, liberté syndicale, droits des communautés autochtones \\\hline
3. Relations et conditions de travail & Emploi, formation, santé-sécurité, dialogue social, égalité professionnelle \\\hline
4. Environnement & Prévention pollution, usage durable ressources, changement climatique, biodiversité \\\hline
5. Loyauté des pratiques & Concurrence loyale, propriété intellectuelle, lutte contre la corruption, relations fournisseurs responsables \\\hline
6. Questions relatives aux consommateurs & Protection santé/sécurité, information loyale, réclamation, données personnelles \\\hline
7. Communautés et développement local & Implication locale, éducation, culture, création d'emplois, technologie appropriée \\\hline
\end{tabularx}
\end{center}

Au Cameroun, la RSE s'articule entre \textbf{obligations légales} et \textbf{initiatives volontaires} :

\begin{itemize}
    \item \textbf{Cadre légal} : Code du travail (santé-sécurité, salaire minimum, représentation du personnel), Code de l'environnement (Évaluation d'Impact Environnemental et Social — EIES), Acte uniforme OHADA sur la comptabilité (transparence financière), Loi n°2011/002 contre la corruption.
    \item \textbf{Initiatives volontaires} : Fondations d'entreprises (Fondation MTN — éducation, santé ; Fondation Orange — culture, inclusion numérique), certifications sectorielles (RSPO pour l'huile de palme durable, FSC pour la gestion forestière responsable), adhésion au Global Compact ONU.
\end{itemize}

\begin{savaistu}[B Corp : le label « Benefit Corporation »]
La certification \textbf{B Corp} (délivrée par l'ONG américaine B Lab) atteste qu'une entreprise répond à des standards élevés de performance sociale et environnementale, de transparence et de responsabilité légale (statuts modifiés pour inclure la mission). À ce jour (2024), \textbf{aucune entreprise camerounaise ne figure au registre public des B Corp certifiées}, mais plusieurs PME exportatrices (cacao, café, tech) préparent leur dossier pour accéder aux marchés européens et nord-américains exigeants.
\end{savaistu}

\courssub{Éthique des affaires : Corruption, Blanchiment, Concurrence déloyale et institutions camerounaises}

L'éthique des affaires est la boussole morale qui guide les décisions au-delà de la simple conformité légale. Trois fléaux majeurs menacent la confiance et l'efficacité économique :

\begin{itemize}
    \item \textbf{Corruption} : abus de pouvoir public/privé pour un avantage indu. \emph{Convention ONU contre la corruption (2003)}, ratifiée par le Cameroun. \emph{Loi n°2011/002 du 6 mai 2011} réprime la corruption active/passive, le trafic d'influence, l'enrichissement illicite. Peines : 5 à 20 ans de prison, amendes lourdes.
    \item \textbf{Blanchiment de capitaux} : dissimulation de l'origine illicite de fonds. Cadre : \emph{Loi n°2017/006} (transposition directive CEMAC), cellule de renseignement financier (ANIF — Agence Nationale d'Investigation Financière).
    \item \textbf{Concurrence déloyale} : ententes sur les prix, abus de position dominante, pratiques trompeuses. Régie par le \emph{Règlement CEMAC n°05/07-UEAC-064-CM-06} et la \emph{Loi n°98/013} sur la concurrence.
\end{itemize}

Deux institutions clés veillent à l'intégrité :

\begin{itemize}
    \item \textbf{CONAC (Commission Nationale Anti-Corruption)} : organe de prévention, d'éducation et de répression. Publie des rapports annuels, reçoit les dénonciations, mène des enquêtes administratives.
    \item \textbf{CONSUPE (Conseil Supérieur de la Concurrence)} : autorité de régulation de la concurrence, sanctionne les ententes, abus de dominance, contrôle les concentrations.
\end{itemize}

\begin{attention}[Pièges fréquents en rédaction d'examen]
\begin{itemize}
    \item Confondre \textbf{RSE} et \textbf{philanthropie} : la RSE est \emph{intégrée} au métier (ex. : éco-conception, achats responsables), la philanthropie est \emph{externe} (dons).
    \item Oublier que la \textbf{gouvernance} est le \emph{premier pilier} de l'ISO 26000 : sans gouvernance éthique, les autres piliers sont décoratifs.
    \item Négliger le \textbf{contexte OHADA} : les organes de gouvernance (AG, CA, DG, CAC) sont imposés par l'Acte uniforme sur les sociétés commerciales ; toute analyse doit s'y référer.
    \item Traiter la ZLECAf comme une simple « opportunité » sans mentionner les \textbf{conditions de compétitivité} (normes, logistique, financement, propriété intellectuelle).
\end{itemize}
\end{attention}

\begin{exoresolu}[Diagnostic de gouvernance — Cas « Bois \& Design » (PME familiale, Yaoundé)]
\textbf{Énoncé :} Bois \& Design SARL (capital 50 millions FCFA) est dirigée par son fondateur M. Ngono (PDG, 70\% parts), son épouse (Directrice Administrative, 20\%) et un gérant salarié (10\%). L'entreprise n'a pas de Conseil d'Administration, pas de Commissaire aux Comptes (seuil non atteint), et les comptes sont tenus par un expert-comptable externe. Les décisions stratégiques (investissement machine CNC, recrutement) sont prises par M. Ngono seul. Deux cadres clés ont démissionné en 2023 faute de perspectives d'évolution.

\textbf{Question :} Identifiez les faiblesses de gouvernance et proposez trois mesures concrètes pour professionnaliser la structure.

\tcblower
\textbf{Solution détaillée :}
\begin{enumerate}
    \item \textbf{Faiblesses identifiées} :
    \begin{itemize}
        \item \textbf{Concentration des pouvoirs} : PDG = actionnaire majoritaire + décideur unique → risque de décisions opportunistes, absence de contre-pouvoir.
        \item \textbf{Absence d'organe collégial} : pas de CA ni de comité consultatif → pas de débat stratégique, pas de contrôle de la gestion.
        \item \textbf{Conflit d'agence famille / managers} : les cadres non familiaux perçoivent l'opacité des promotions et de la rémunération → turnover.
        \item \textbf{Transparence financière limitée} : pas de CAC, comptes non certifiés → difficulté d'accès au financement bancaire ou investisseurs.
    \end{itemize}
    \item \textbf{Mesures de professionnalisation} :
    \begin{enumerate}
        \item \textbf{Créer un Conseil Consultatif (Advisory Board)} composé de 3 à 5 membres externes (expert sectoriel, financier, juriste, ancien DG) réunis trimestriellement. Rôle : avis stratégique, challenge du PDG, préparation d'un futur CA formel.
        \item \textbf{Nommer un Commissaire aux Comptes volontaire} (même si le seuil légal n'est pas atteint) pour certifier les comptes annuels → crédibilité auprès des banques, signal positif pour d'éventuels investisseurs minoritaires.
        \item \textbf{Formaliser une Charte de Gouvernance Familiale} : règles de succession, critères de nomination des membres de la famille, politique de rémunération transparente, procédure d'évaluation des cadres non familiaux (grilles de salaire, plan de carrière).
    \end{enumerate}
\end{enumerate}
\textbf{Conclusion :} Ces mesures, peu coûteuses, transforment la gouvernance « patriarcale » en gouvernance \emph{responsable}, alignent les intérêts famille/managers et préparent l'ouverture du capital.
\end{exoresolu}

\begin{exercice}[Application : SWOT \& RSE — « Café du Haut-Penja »]
L'entreprise « Café du Haut-Penja » (coopérative de 300 producteurs, unité de torréfaction à Penja) souhaite exporter du café torréfié spécialité vers l'Europe. Elle dispose d'une certification \textbf{Fairtrade} et \textbf{Bio}, mais son emballage est en plastique non recyclable, son site web est inexistant et elle dépend d'un unique exportateur français.
\begin{enumerate}
    \item Réalisez une matrice SWOT (4 items par quadrant) pour cette coopérative.
    \item Identifiez deux piliers ISO 26000 prioritaires et proposez une action concrète pour chacun.
    \item Quelle stratégie d'internationalisation (spécialisation / diversification / intégration) recommandez-vous ? Justifiez en 3 arguments.
\end{enumerate}
\end{exercice}

\begin{corrige}[Exercice — Café du Haut-Penja]
\begin{enumerate}
    \item \textbf{SWOT} :
    \begin{itemize}
        \item \textbf{F} : Café arabica d'altitude (qualité), certifications Fairtrade/Bio, coopérative structurée, terroir réputé.
        \item \textbf{W} : Emballage plastique, absence marketing digital, dépendance unique acheteur, trésorerie tendue (avances campagne).
        \item \textbf{O} : Marché européen café spécialité (+12\%/an), ZLECAf (accès marchés africains), programmes d'appui UE (COLEACP), demande croissante emballages éco-conçus.
        \item \textbf{T} : Concurrence Brésil/Colombie (coûts logistiques), volatilité cours café (bourse), réglementation EUDR (déforestation importée), fluctuations change FCFA/EUR.
    \end{itemize}
    \item \textbf{Piliers prioritaires} :
    \begin{itemize}
        \item \textbf{Environnement (pilier 4)} : Remplacer l'emballage plastique par du sachet kraft compostable certifié OK Compost Home → réduction empreinte carbone, conformité EUDR.
        \item \textbf{Loyauté des pratiques (pilier 5)} : Négocier un contrat cadre pluriannuel avec l'exportateur français incluant clause de révision prix indexée sur le cours international + clause de transparence sur la chaîne de valeur → relation équilibrée, lutte contre la dépendance.
    \end{itemize}
    \item \textbf{Stratégie recommandée : Spécialisation + Intégration aval partielle}.
    \begin{enumerate}
        \item \textbf{Spécialisation} : recentrer sur le café de spécialité (score SCA > 85) pour valoriser le terroir et justifier prime prix.
        \item \textbf{Intégration aval} : créer une marque propre « Penja Peak » vendue en direct via e-commerce (site + marketplaces UE) et ouvrir un showroom à Douala/Yaoundé pour capter la marge de distribution et diversifier les canaux.
        \item \textbf{Argument 3} : La diversification (ex. : cacao, épices) diluerait l'expertise et les certifications ; l'intégration amont (achat cerises) est déjà assurée par la coopérative.
    \end{enumerate}
\end{enumerate}
\end{corrige}

\coursec{TP Intégré : Diagnostic simplifié d'une PME camerounaise (Cas « Bois \& Design »)}

Après avoir analysé les concepts de gouvernance, de stratégie et de responsabilité sociétale (RSE) dans la section précédente, nous passons maintenant à la mise en œuvre concrète de ces outils d'analyse. Ce TP intégré constitue une \textbf{situation d'évaluation formative} : il vous demande de mobiliser l'ensemble des notions du chapitre (statut juridique, facteurs de production, organisation, indicateurs financiers, RSE) pour porter un jugement critique sur la performance et la durabilité d'une entreprise réelle. C'est l'exercice type de l'épreuve du Probatoire : « Appliquer les notions de l'unité à des situations concrètes » et « Rédiger et justifier une démarche ».

\courssub{Présentation du dossier : « Bois \& Design SARL »}

\begin{exemplebox}[Contexte de l'étude]
\textbf{Bois \& Design SARL} est une entreprise camerounaise située à Douala (quartier Bassa), spécialisée dans la \textbf{2\ieme{} et 3\ieme{} transformation du bois local} (essences : Ayous, Iroko, Movingui). Elle conçoit et fabrique du mobilier haut de gamme (bureaux directionnels, literie de luxe, agencement hôtelier).
\begin{itemize}
    \item \textbf{Effectif :} 45 salariés (dont 30 ouvriers qualifiés menuisiers/ébénistes, 5 contremaîtres, 5 administratifs, 3 commerciaux, 1 comptable unique, 1 chef d'atelier).
    \item \textbf{Marchés :} 60\% Export (France, Belgique via conteneurs depuis le Port de Douala), 40\% Local (hôtels, administrations, particuliers aisés).
    \item \textbf{Chiffre d'Affaires 2023 :} 450 millions FCFA.
    \item \textbf{Forme juridique :} SARL (Société à Responsabilité Limitée) au capital de 50 millions FCFA.
    \item \textbf{Dirigeant :} M. Ngono Essomba (Gérant unique, associé majoritaire 60\%), ingénieur bois de formation.
\end{itemize}
\end{exemplebox}

\noindent \textbf{Documents fournis aux élèves (Annexes du TP) :}
\begin{enumerate}
    \item \textbf{Extrait K-bis (Registre du Commerce) :} Immatriculation RC/DLA/2015/B/12345. Objet social : « Transformation du bois, fabrication et commerce de meubles, agencement intérieur ». Associés : M. Ngono (60\%), Mme Ngono (20\%), M. Kouam (20\%).
    \item \textbf{Organigramme 2023 (incomplet à compléter) :} Structure hiérarchique centralisée sur le Gérant.
    \item \textbf{Compte de Résultat Simplifié 2023 (en millions FCFA) :} Voir illustration ci-dessous.
    \item \textbf{Extrait Rapport RSE 2023 :} « Bois \& Design s'engage pour la gestion durable. Démarche FSC en cours. Partenariat avec l'INFOP pour apprentis. Pas de comité d'entreprise. Accidents du travail : 3 en 2023 (coupures, bruit). Turn-over menuisiers : 15\%/an. »
\end{enumerate}

\begin{popfigure}
\centering
\begin{tabularx}{\linewidth}{|l|X|}\hline
\rowcolor{popblueL} \textbf{Compte de Résultat Simplifié 2023 (M FCFA)} & \textbf{Montant} \\ \hline
Chiffre d'Affaires (CA) & 450 \\ \hline
Achats consommés (Matières Premières, Énergie, Sous-traitance) & -180 \\ \hline
\rowcolor{popcream} \textbf{Valeur Ajoutée (VA)} & \textbf{270} \\ \hline
Charges de Personnel (Salaires + Charges sociales) & -120 \\ \hline
Impôts et Taxes sur production (Taxe apprentissage, Patente...) & -5 \\ \hline
\rowcolor{popcream} \textbf{Excédent Brut d'Exploitation (EBE)} & \textbf{145} \\ \hline
Dotations aux Amortissements & -30 \\ \hline
Charges Financières (Intérêts emprunts) & -10 \\ \hline
\rowcolor{popcream} \textbf{Résultat Courant Avant Impôt (RCAI)} & \textbf{105} \\ \hline
Impôt sur les Sociétés (IS \textasciitilde 23,8\%) & -25 \\ \hline
\rowcolor{popgreen} \textbf{Résultat Net (Bénéfice)} & \textbf{80} \\ \hline
\end{tabularx}
\legende{Document support : Extrait du Compte de Résultat Simplifié 2023 de Bois \& Design SARL. Chiffres arrondis pour la lisibilité pédagogique. L'IS est calculé sur le RCAI (105 $\times$ 25\% $\approx$ 26,25M), arrondi à 25M pour donner un RN de 80M cohérent avec les SIG.}
\end{popfigure}

\begin{popfigure}
\centering
\begin{tikzpicture}[
    node distance=1.2cm and 1.5cm,
    every node/.style={font=\small, align=center},
    box/.style={draw=popdark, thick, rounded corners, fill=popcream, minimum width=2.8cm, minimum height=0.9cm},
    boxg/.style={box, fill=popgreenL, text width=3cm},
    boxb/.style={box, fill=popblueL, text width=3cm},
    boxo/.style={box, fill=poporangeL, text width=3cm},
    arrow/.style={-Stealth, thick, popdark}
]
% Niveau 1
\node[boxg, align=center] (gerant){Gérant Unique\\ (M. Ngono)\\ \textit{Direction Générale}\\\textit{ Stratégie, Finance, RH}};
% Niveau 2
\node[boxb, below left=of gerant, xshift=-2cm, align=center] (atelier){Chef d'Atelier\\ \textit{Production, Qualité,}\\\textit{ Maintenance}};
\node[boxb, below=of gerant, align=center] (commercial){Responsable Commercial\\ \textit{Export / Local,}\\\textit{ Devis, Clientèle}};
\node[boxb, below right=of gerant, xshift=2cm, align=center] (admin){Comptable Unique\\ \textit{Comptabilité, Paie,}\\\textit{ Déclarations fiscales}};
% Niveau 3 (Sous Chef Atelier)
\node[boxo, below left=of atelier, xshift=-1.5cm, align=center] (cm1){Contremaître\\ \textit{Usinage / Sciage}};
\node[boxo, below=of atelier, align=center] (cm2){Contremaître\\ \textit{Assemblage / Finition}};
\node[boxo, below right=of atelier, xshift=1.5cm, align=center] (cm3){Contremaître\\ \textit{Conditionnement / Expédition}};
% Niveau 4 Ouvriers
\node[box, below left=of cm1, xshift=-1cm, align=center] (ouv1){Ouvriers Qualifiés\\ (Menuisiers, Ébénistes)\\ $\times$ 10};
\node[box, below=of cm2, align=center] (ouv2){Ouvriers Qualifiés\\ (Assembleurs, Finisseurs)\\ $\times$ 12};
\node[box, below right=of cm3, xshift=1cm, align=center] (ouv3){Ouvriers / Manutention\\ $\times$ 8};
% Flèches
\draw[arrow] (gerant) -- (atelier);
\draw[arrow] (gerant) -- (commercial);
\draw[arrow] (gerant) -- (admin);
\draw[arrow] (atelier) -- (cm1);
\draw[arrow] (atelier) -- (cm2);
\draw[arrow] (atelier) -- (cm3);
\draw[arrow] (cm1) -- (ouv1);
\draw[arrow] (cm2) -- (ouv2);
\draw[arrow] (cm3) -- (ouv3);
% Annotations manquantes (pointillés)
\node[draw=popred, dashed, thick, rounded corners, fit=(admin) (commercial) (atelier), label={[popred]above right:Absence DRH / Service Qualité / Logistique dédié}] {};
\end{tikzpicture}
\legende{Document support : Organigramme 2023 de Bois \& Design SARL (Structure hiérarchique centralisée). Les zones pointillées rouges indiquent les fonctions manquantes ou sous-dimensionnées repérées lors de l'analyse.}
\end{popfigure}

\courssub{Activité 1 : Analyse Juridique et Structurelle (20 min)}

\begin{methode}[Identifier et Classer l'Entreprise]
\textbf{Étape 1 :} Lire l'extrait K-bis. Repérer : Dénomination, Forme, Siège, Capital, Objet social, Dirigeants, Associés.
\textbf{Étape 2 :} Classer selon 3 critères officiels (INS / OHADA) :
\begin{itemize}
    \item \textbf{Taille :} Effectif (45) $\rightarrow$ \cle{PME} (Petite et Moyenne Entreprise : 20 $\le$ Eff $<$ 250). Capital (50M) $\rightarrow$ PME.
    \item \textbf{Secteur d'activité :} Transformation bois $\rightarrow$ \cle{Secteur Secondaire} (Industrie manufacturière). Code NACE/ISIC : 3100 (Fabrication de meubles).
    \item \textbf{Capital :} Capital détenu par des personnes physiques nationales $\rightarrow$ \cle{Secteur Privé} national. Pas d'actionnariat public $\rightarrow$ Pas d'EPIC.
\end{itemize}
\textbf{Étape 3 :} Identifier l'organe de direction : \textbf{Gérant Unique} (M. Ngono). Pouvoirs étendus (Art. 68 Acte Uniforme OHADA SARL) : engage la société en toutes circonstances.
\end{methode}

\begin{exoresolu}[Activité 1 - Réponses attendues]
\textbf{1. Forme Juridique :} SARL (Société à Responsabilité Limitée). \textbf{Capital :} 50 000 000 FCFA. \textbf{Associés :} 3 personnes physiques (M. Ngono 60\%, Mme Ngono 20\%, M. Kouam 20\%). Responsabilité limitée aux apports.
\textbf{2. Objet Social :} Transformation du bois, fabrication et commerce de meubles, agencement intérieur. \textit{Conforme à l'activité réelle (cohérence juridique).}
\textbf{3. Classification :}
\begin{itemize}
    \item Taille : \textbf{PME} (Effectif 45, CA 450M $<$ 1 Md, Bilan $<$ 1 Md).
    \item Secteur : \textbf{Secondaire} (Industrie / Transformation).
    \item Capital : \textbf{Privé} national (familial/associés).
\end{itemize}
\textbf{4. Organe de Direction :} \textbf{Gérant Unique} (M. Ngono). Concentre les pouvoirs de gestion, représentation légale, signature chèques, embauche/licenciement. \textit{Risque : concentration des pouvoirs (voir Activité 3).}
\end{exoresolu}

\courssub{Activité 2 : Analyse Économique et Financière (30 min)}

\begin{methode}[Calculer et Interpréter les Soldes Intermédiaires de Gestion (SIG)]
\textbf{1. Valeur Ajoutée (VA) :} Mesure la richesse \textbf{créée} par l'entreprise (différence entre ce qu'elle vend et ce qu'elle achète à l'extérieur).
\[ \text{VA} = \text{Chiffre d'Affaires} - \text{Consommations Intermédiaires (Achats MP + Énergie + Sous-traitance + Services extérieurs)} \]
\textit{Interprétation :} La VA se partage entre : Salaires (Travail) + Intérêts/Impôts (Capital/État) + EBE (Autofinancement).
\textbf{2. Excédent Brut d'Exploitation (EBE) :} Excédent dégagé par le \textbf{cycle d'exploitation} seul (avant politique d'investissement/amortissement et financement).
\[ \text{EBE} = \text{VA} - \text{Charges de Personnel} - \text{Impôts/taxes sur production} \]
\textbf{3. Résultat Net :} Bénéfice final après \textbf{toutes} charges (amortissements, financiers, IS). Indicateur de performance comptable.
\textbf{4. Taux de Marge Nette :} $\frac{\text{Résultat Net}}{\text{CA}} \times 100$. Rentabilité finale.
\textbf{5. Productivité Apparente du Travail :} $\frac{\text{VA}}{\text{Effectif Moyen}}$. Richesse créée par employé/an.
\end{methode}

\begin{attention}[Piège classique : Résultat Net $\neq$ Trésorerie / CAF]
Ne confondez jamais \textbf{Résultat Net} (comptabilité d'engagement, hors flux) et \textbf{Capacité d'AutoFinancement (CAF)} (flux réel dégagé).
\[ \text{CAF} = \text{Résultat Net} + \text{Amortissements} + \text{Provisions} - \text{Reprises} + \text{Plus-values cession} - \text{Moins-values cession} \]
Ici : CAF $\approx$ 80 + 30 = \textbf{110 M FCFA}. C'est cette somme qui permet de rembourser les emprunts, investir, verser des dividendes. Le Résultat Net (80M) surestime la capacité de paiement s'il y a de gros investissements (BFR, CAPEX).
\end{attention}

\begin{exoresolu}[Activité 2 - Calculs et Interprétation (Chiffres du dossier)]
\textbf{Données :} CA = 450M ; Achats = 180M ; Charges Perso = 120M ; Amort. = 30M ; Ch. Fin = 10M ; IS = 25M ; Impôts Prod = 5M (estimé). Effectif = 45.

\begin{enumerate}
    \item \textbf{VA} = 450 - 180 = \textbf{270 M FCFA}.
    \item \textbf{EBE} = 270 - 120 - 5 = \textbf{145 M FCFA}. \textit{Taux EBE/CA = 32\% (Excellent, norme secteur \textasciitilde 15-20\%). L'exploitation est très performante.}
    \item \textbf{Résultat Net} = 80 M FCFA (vérifié : 145 - 30 - 10 - 25 = 80M).
    \item \textbf{Taux de Marge Nette} = 80 / 450 = \textbf{17,8\%}. \textit{Très bon niveau (moyenne secteur meuble \textasciitilde 8-12\%). L'entreprise est très rentable.}
    \item \textbf{Productivité Apparente du Travail} = 270M / 45 = \textbf{6 M FCFA / employé / an}.
    \textit{Analyse :} Moyen pour de l'industrie bois haut de gamme (cible souvent $>$ 10M). Explique par : main d'œuvre nombreuse peu automatisée, turn-over (perte de savoir-faire), organisation centralisée freinant l'efficacité.
\end{enumerate}
\textbf{Synthèse Financière :} Entreprise \textbf{rentable} (marges fortes), \textbf{solvable} (EBE couvre largement charges financières), \textbf{autofinancement} solide (CAF 110M). Point de vigilance : Productivité perfectible, dépendance aux matières premières (achats 40\% CA).
\end{exoresolu}

\courssub{Activité 3 : Analyse de l'Organisation et des Ressources Humaines (20 min)}

\begin{experience}[Analyse de l'Organigramme : Structure et Dysfonctionnements]
\textbf{Observation du document (Figure 2) :}
\begin{enumerate}
    \item \textbf{Type de structure :} \textbf{Hiérarchique (Pyramidale) / Fonctionnelle (par métier : Prod, Com, Admin)}. Liens hiérarchiques clairs (linéaires). Pas de transversalité visible (pas de chef de projet, pas de comité de direction).
    \item \textbf{Span of Control (Étendue de contrôle) :} Gérant $\rightarrow$ 3 directs (Correct). Chef Atelier $\rightarrow$ 3 Contremaîtres (Correct). Contremaîtres $\rightarrow$ $\sim$10 ouvriers (Correct).
    \item \textbf{Point critique 1 : Centralisation excessive.} Le Gérant est \textbf{unique} décideur : Stratégie + Finance + RH + Commercial + Production. \textit{Risque : Goulet d'étranglement (bottleneck), fatigue décisionnelle, absence de relais en cas d'absence.}
    \item \textbf{Point critique 2 : Absence de DRH / Service RH.} Gestion du personnel (paie, conflits, formation, recrutement, sécurité) gérée par le Comptable Unique ou le Gérant. \textit{Conséquence directe : Turn-over 15\% menuisiers qualifiés (coût de remplacement élevé, perte qualité), 3 accidents/an (prévention défaillante), pas de plan de formation formel.}
    \item \textbf{Point critique 3 : Absence Service Qualité / Logistique dédié.} Export exige traçabilité, normes emballage (NIMP15), gestion conteneurs. Confié au Commercial ou Chef Atelier $\rightarrow$ risques litiges clients, retards.
\end{enumerate}
\end{experience}

\begin{aretenir}[Diagnostic Organisationnel]
Bois \& Design souffre du \textbf{syndrome du « Patron tout-puissant »} typique des PME fondatrices en phase de croissance. La structure n'a pas suivi la complexité (Export, 45 pers). L'absence de \cle{middle management} (Directeur Opérationnel, DRH) et de \cle{transversalité} (Réunions production/commercial, Qualité) freine la productivité (6M/employé) et crée des risques RH (turn-over, accidents) et juridiques (conformité export).
\end{aretenir}

\courssub{Activité 4 : Diagnostic Stratégique et RSE — Matrice SWOT (25 min)}

\begin{methode}[Construire une Matrice SWOT (FFOM en français)]
\begin{itemize}
    \item \textbf{Forces (Interne +) :} Atouts, ressources, compétences distinctives.
    \item \textbf{Faiblesses (Interne -) :} Manques, dysfonctionnements, ressources insuffisantes.
    \item \textbf{Opportunités (Externe +) :} Marchés en croissance, évolutions réglementaires favorables, technologies.
    \item \textbf{Menaces (Externe -) :} Concurrence, réglementation contraignante, instabilité, matières premières.
\end{itemize}
\textbf{Règle :} Croiser pour définir la stratégie : FO (Offensive), FA (Renforcement), MO (Vigilance), MA (Survie/Désinvestissement).
\end{methode}

\begin{exemplebox}[Matrice SWOT — Bois \& Design SARL (Proposition corrigée)]
\begin{center}
\begin{tabularx}{\linewidth}{|X|X|}\hline
\rowcolor{popblueL} \textbf{FORCES (Interne +)} & \textbf{FAIBLESSES (Interne -)} \\ \hline
\begin{itemize}
    \item Savoir-faire ébénisterie haut de gamme (niche).
    \item Rentabilité forte (Marge 17,8\%, EBE 32\%).
    \item Trésorerie saine (CAF 110M, peu d'endettement).
    \item Bois local certifiable (Ayous/Iroko disponibles).
    \item Références clients Export (portefeuille fidélisé).
\end{itemize} & \begin{itemize}
    \item Gouvernance : Gérant unique, pas de succession.
    \item Organisation : Centralisée, pas DRH, pas Qualité.
    \item RH : Turn-over 15\% (menuisiers), 3 accidents/an.
    \item Productivité travail moyenne (6M/employé).
    \item Dépendance 1 client export / 1 fournisseur bois ?
\end{itemize} \\ \hline
\rowcolor{popgreenL} \textbf{OPPORTUNITÉS (Externe +)} & \textbf{MENACES (Externe -)} \\ \hline
\begin{itemize}
    \item \textbf{RDUE (Règlement UE Déforestation) :} Exige traçabilité $\rightarrow$ Avantage si certifié FSC.
    \item \textbf{ZLECAf :} Marché africain 1,3 Md conso, droits de douane 0\%.
    \item Demande locale croissante (immobilier, hôtels standing).
    \item Financements verts (Banques, AFD, BEAC) pour certification.
    \item Digitalisation (Configurateur 3D, e-commerce B2B).
\end{itemize} & \begin{itemize}
    \item \textbf{RDUE :} Interdiction entrée UE si pas traçabilité (risque perte 60\% CA).
    \item Concurrence asiatique (bois composite, prix bas).
    \item Pénurie/vol bois local (braconnage, pisteurs).
    \item Inflation coûts énergie/transport (Douala portuaire).
    \item Instabilité fiscale / sociale (grèves, routes).
\end{itemize} \\ \hline
\end{tabularx}
\end{center}
\end{exemplebox}

\begin{definition}[Certification FSC (Forest Stewardship Council)]
Label international indépendant garantissant que le bois provient de forêts gérées \textbf{durablement} sur 3 piliers : \textbf{Écologique} (biodiversité, produits chimiques), \textbf{Social} (droits travailleurs, populations autochtones, communautés locales), \textbf{Économique} (viabilité). \textbf{Prérequis majeur} pour le marché européen depuis le \textbf{RDUE (Règlement Déforestation UE, applicable 2025)} : tout opérateur doit prouver (géolocalisation parcelles) que le bois ne cause pas de déforestation. Sans FSC (ou équivalent PEFC/PAFC), \textbf{Bois \& Design perd son marché export principal}.
\end{definition}

\begin{savaistu}[Le paradoxe camerounais du bois]
Le Cameroun possède \textbf{22 millions d'hectares} de forêts (2\ieme{} bassin du Congo). Pourtant, l'export est historiquement dominé par les \textbf{grumes et bouvet (1\iere{} transformation)}, à faible valeur ajoutée. L'industrialisation locale (2\ieme/3\ieme transformation : meubles, parquets, panneaux) est un pilier de la \textbf{SND30 (Stratégie Nationale de Développement 2020-2030)}. « Bois \& Design » illustre cette transition difficile : une PME qui crée de la valeur (VA 270M) mais bute sur l'organisation, la certification et la formation. Le \textbf{ZLECAf} (Zone de Libre-Échange Continentale Africaine) offre une opportunité historique de vendre du meuble « Made in Cameroon » en Afrique sans droits de douane.
\end{savaistu}

\courssub{Activité 4 (Suite) : Axe Stratégiques Prioritaires et Conclusion Argumentée}

À partir du croisement SWOT, deux axes s'imposent en \textbf{stratégie FO (Forces-Opportunités)} et \textbf{FA (Forces-Faiblesses)} :

\begin{propriete}[Axe Stratégique 1 : « Certification FSC \& Traçabilité Numérique » (Urgence Vitale - Menace RDUE)]
\textbf{Objectif :} Sécuriser le CA Export (270M FCFA) et ouvrir marchés premium.
\textbf{Actions :}
\begin{enumerate}
    \item Finaliser audit FSC (coût estimé 5-10M FCFA, durée 6-12 mois). Chercher co-financement (PME Vert, AFD).
    \item Investir dans logiciel traçabilité (blockchain/QR code) : de la parcelle forêt $\rightarrow$ scierie $\rightarrow$ atelier $\rightarrow$ conteneur. Répondre exigence RDUE (géolocalisation).
    \item Valoriser marketing : « Meuble 100\% Légal, Durable, Traçable » $\rightarrow$ Prix premium +5-10\%.
\end{enumerate}
\textbf{Indicateur :} Certificat FSC obtenu avant déc. 2025. 100\% lots export traçables.
\end{propriete}

\begin{propriete}[Axe Stratégique 2 : « Capital Humain \& Organisation : Académie Bois \& Design » (Levier Productivité - Faiblesse RH)]
\textbf{Objectif :} Faire passer Productivité de 6M $\rightarrow$ 10M FCFA/employé/an. Réduire turn-over $<$ 5\%. Zéro accident.
\textbf{Actions :}
\begin{enumerate}
    \item Créer poste \textbf{DRH / Responsable Formation} (recrutement 1 cadre ou promotion interne Comptable + formation).
    \item Lancer \textbf{École interne / Apprentissage} : Partenariat INFOP / Lycée Technique de Douala. Former 5 apprentis/an (Menuisier/Ébéniste CNC/Traditionnel). Coût partagé État/Entreprise.
    \item Instaurer \textbf{Gestion par les compétences / Prime au mérite / Qualité} : Réunions hebdo Prod/Com (transversalité), Cercles de qualité (Kaizen), Prime sécurité (0 accident).
    \item Délégation : Nommer \textbf{Directeur Opérationnel Adjoint} (Chef Atelier promu ?) pour libérer Gérant $\rightarrow$ Stratégie / Finance / RSE.
\end{enumerate}
\textbf{Indicateur :} Turn-over $<$ 5\% à 2 ans. Productivité $>$ 9M. Certification ISO 45001 (SST) visée 2027.
\end{propriete}

\begin{methode}[Rédaction Bac : Structure « Diagnostic $\rightarrow$ Analyse $\rightarrow$ Recommandation »]
Pour la conclusion écrite (15 lignes) ou l'oral (5 min), suivez ce canevas rigoureux :
\begin{enumerate}
    \item \textbf{Introduction (2-3 lignes) :} Accroche (contexte bois Cameroun) + Présentation objet (Bois \& Design, PME 45p, 450M CA) + \textbf{Problématique} : « Dans quelle mesure Bois \& Design, malgré sa performance financière, est-elle structurellement fragile face aux enjeux de durabilité (RDUE) et de capital humain, et quels leviers actionner ? »
    \item \textbf{Développement (2 parties argumentées, 8-10 lignes) :}
        \begin{itemize}
            \item \textbf{Partie 1 : Performance réelle mais modèle fragile.} Appui chiffres : VA 270M, EBE 145M, Marge 17,8\%, CAF 110M (Force). MAIS : Centralisation Gérant, Turn-over 15\%, Productivité 6M (Faiblesse). Risque majeur : RDUE (Menace) = Mort export si pas FSC.
            \item \textbf{Partie 2 : Deux leviers indissociables pour la durabilité.} (Axe 1) Certification FSC + Traçabilité = Condition \textit{sine qua non} survie export (RSE Environnement/Économie). (Axe 2) DRH + Formation interne + Délégation = Condition efficacité productive \& sociale (RSE Social/Gouvernance). Synergie : Bois certifié $\rightarrow$ Main d'œuvre qualifiée pour le travailler.
        \end{itemize}
    \item \textbf{Conclusion (2-3 lignes) :} Réponse synthétique : « Bois \& Design est performante \textit{aujourd'hui} mais \textbf{non durable} \textit{demain} sans transformation profonde. » Ouverture : « Réussir ce double pari (FSC + RH) ferait de l'entreprise un modèle de l'industrialisation verte camerounaise, duplicable via ZLECAf. »
\end{enumerate}
\end{methode}

\begin{exoresolu}[Exemple de Conclusion Argumentée (15 lignes)]
\textit{« Bois \& Design affiche une santé financière enviable (Marge nette 17,8\%, CAF 110M), preuve d'un positionnement « niche haut de gamme » maîtrisé techniquement. Pourtant, le diagnostic révèle une entreprise « à haut risque » structurellement. La centralisation extrême sur le Gérant unique, l'absence de DRH et la productivité moyenne (6M/employé) traduisent une organisation inadaptée à la complexité export. Surtout, la menace RDUE agit comme une épée de Damoclès : sans certification FSC et traçabilité numérique d'ici 2025, 60\% du CA s'évaporent. L'entreprise n'est donc pas durable \textit{en l'état}. Sa pérennité exige un double basculement stratégique indissociable. Premièrement, l'urgence environnementale et commerciale : investir massivement dans la certification FSC et la traçabilité blockchain pour transformer la contrainte RDUE en avantage concurrentiel « premium » sur les marchés européen et africain (ZLECAf). Deuxièmement, le pari humain : créer une fonction RH, lancer une académie interne de formation aux métiers du bois (CNC, ébénisterie) et déléguer l'opérationnel pour fidéliser les talents et sécuriser la qualité. Seule cette alliance « Bois Certifié / Main-d'Œuvre Qualifiée » ancrera Bois \& Design dans une performance durable, modèle de la 3\ieme{} transformation locale prônée par la SND30. »}
\end{exoresolu}

\courssub{Restitution : Présentation Orale par Groupes (5 min / groupe)}

\begin{methode}[Consignes de l'Oral « Grand Oral » Style Probatoire]
\begin{enumerate}
    \item \textbf{Constitution :} Groupes de 4-5 élèves. Rôles : \textbf{Présentateur principal} (2 min), \textbf{Expert Financier} (1 min), \textbf{Expert RH/Org} (1 min), \textbf{Expert Stratégie/RSE} (1 min), \textbf{Chronométreur/Relance}.
    \item \textbf{Sujet imposé :} \textbf{« Bois \& Design est-elle une entreprise performante et durable ? Justifiez. »}
    \item \textbf{Structure imposée :} 1) Diagnostic Flash (Chiffres clés) 2) Le « Mais » critique (Point de rupture) 3) Les 2 Recommandations Prioritaires (FSC + RH) 4) Verdict final (Performante OUI / Durable NON $\rightarrow$ SI transformation).
    \item \textbf{Critères d'évaluation (Grille) :} Justesse concepts (VA, EBE, SWOT, RSE, Gouvernance) / 5 ; Qualité argumentation chiffrée / 5 ; Structure logique / 3 ; Expression orale / 2 ; Réponse aux questions (Jury/Classe) / 5. \textbf{Total /20.}
\end{enumerate}
\end{methode}

\begin{exercice}[Entraînement Individuel — Sujet Type Bac]
\textbf{Sujet :} « À partir du dossier Bois \& Design et de vos connaissances, montrez que la performance économique d'une PME industrielle camerounaise ne saurait se réduire à son seul résultat net. »
\textbf{Consignes :} Rédigez une réponse structurée (Intro, Développement 2 parties, Conclusion) en 1h30. Mobilisez : SIG (VA, EBE, CAF), Organisation (Gouvernance, RH), Environnement (RDUE, ZLECAf, RSE), Concepts (Durabilité, Valeur Ajoutée, Capital Humain).
\end{exercice}

\begin{corrige}[Exercice — Éléments de Corrigé Attendus]
\textbf{Introduction :} Accroche sur enjeu industrialisation Cameroun (SND30). Définition PME / Résultat Net. Problématique : Le RN est indicateur nécessaire mais insuffisant ; la durabilité exige analyse de la création de valeur (VA), du modèle organisationnel et de l'ancrage territorial/RSE. Annonce plan.
\textbf{Partie I : Le Résultat Net, indicateur comptable réducteur de la performance.}
\begin{itemize}
    \item RN = 80M (Bois \& Design) $\rightarrow$ Image statique, passé, manipulable (amortissements, provisions).
    \item Nécessité d'analyser la \textbf{VA (270M)} : Richesse réellement créée, partage (Salaires 120M, EBE 145M, État 30M). Taux VA/CA = 60\% (fort, peu sous-traitance).
    \item \textbf{EBE (145M) \& CAF (110M)} : Vrais indicateurs de performance \textit{opérationnelle} et de \textit{solvabilité}/investissement futur. RN ne dit rien sur la trésorerie dégagée.
\end{itemize}
\textbf{Partie II : La Durabilité exige d'aller au-delà du compte de résultat : Organisation, Gouvernance, RSE.}
\begin{itemize}
    \item \textbf{Gouvernance / RH :} Gérant unique = Risque systémique (clé homme). Turn-over 15\% = Destruction capital humain implicite (coût recrutement, perte savoir-faire). Absence DRH = Défaillance gouvernance (RSE Social). Productivité 6M = Sous-optimum organisationnel.
    \item \textbf{Environnement / Marché :} RDUE = Menace existentielle (60\% CA). Performance actuelle \textit{dépend} d'une faille réglementaire future. RSE Environnement (FSC) = Condition \textit{sine qua non} pérennité.
    \item \textbf{Ancrage Territorial :} Achats bois local (180M) = Soutien filière amont. Export = Devises. Mais faible transformation locale amont (grumes) = Dépendance fournisseurs.
\end{itemize}
\textbf{Conclusion :} Le RN de Bois \& Design masque une fragilité structurelle (Gouvernance, RH) et conjoncturelle (RDUE). Une PME performante \textit{durablement} est celle qui aligne \textbf{Performance Économique (VA, EBE)}, \textbf{Efficience Organisationnelle (RH, Processus)} et \textbf{Responsabilité Sociétale (FSC, SST, Gouvernance partagée)}. Ouverture : Rôle de l'État (incitations fiscales FSC, formation INFOP) et ZLECAf pour changer d'échelle.
\end{corrige}

\newpage

\coursec{Activité d'intégration}

\coursec{TP Intégré : Diagnostic simplifié d'une PME camerounaise — Cas « GreenEnergy Cameroon »}

\courssub{Objectifs de l'activité}

Cette activité d'intégration clôt la leçon sur \cle{Les Entreprises}. Elle invite chaque groupe d'élèves à mobiliser, dans une situation authentique de création d'entreprise au Cameroun, l'ensemble des savoirs de l'unité :

\begin{itemize}
\item la \cle{définition} et la nature de l'entreprise (unité économique et sociale combinant des facteurs de production pour produire des biens ou services marchands) ;
\item la \cle{typologie} des entreprises (secteur d'activité, taille, forme juridique selon le droit OHADA) ;
\item le \cle{fonctionnement interne} : facteurs de production (travail, capital, technique), organisation, prévisionnel financier simple ;
\item la \cle{gouvernance}, la stratégie (matrice SWOT) et la \cle{Responsabilité Sociétale des Entreprises} (RSE).
\end{itemize}

Les savoir-faire visés sont ceux du programme officiel : \emph{appliquer les notions de l'unité à des situations concrètes de la vie courante} et \emph{rédiger et justifier une démarche, une réponse ou une conclusion}. Au terme de l'activité, chaque groupe produit une \textbf{Note de cadrage stratégique} argumentée, évaluée sur 20 points.

\courssub{Situation}

\textbf{Le dilemme du promoteur : choisir la forme et la stratégie pour « GreenEnergy Cameroon ».}

Trois jeunes diplômés camerounais — Amina (ingénieure informaticienne), Basile (électrotechnicien) et Carine (gestionnaire) — décident de créer à Yaoundé l'entreprise \emph{GreenEnergy Cameroon}, spécialisée dans l'installation et la maintenance de kits solaires photovoltaïques pour les ménages et les PME des villes de Yaoundé et Douala. Leur constat de départ est simple : les coupures d'électricité restent fréquentes, le coût des groupes électrogènes explose, et la demande en énergie solaire croît chaque année.

\textbf{Données du dossier :}

\begin{itemize}
\item \textbf{Capital d'apport initial :} 5\,000\,000 FCFA (apports personnels des trois associés).
\item \textbf{Objectifs :} rentabilité atteinte en 3 ans ; impact social mesurable (accès à l'énergie des ménages à revenus modestes) ; croissance rapide avec entrée possible d'investisseurs.
\item \textbf{Étude de marché simplifiée :} forte demande urbaine ; concurrence des kits chinois bon marché (qualité inégale, sans service après-vente) et de quelques installateurs locaux informels ; pouvoir d'achat limité, d'où l'idée d'un paiement échelonné sur 12 mois.
\item \textbf{Prévisionnel simplifié sur 3 ans :}
\end{itemize}

\begin{center}
\begin{tabularx}{\linewidth}{|l|X|X|X|}
\hline
\rowcolor{popblueL} \textbf{Élément (en FCFA)} & \textbf{Année 1} & \textbf{Année 2} & \textbf{Année 3} \\
\hline
Investissements (matériel, véhicule, stock) & 3\,500\,000 & 1\,000\,000 & 1\,500\,000 \\
\hline
Charges fixes (loyer, salaires, assurances) & 4\,800\,000 & 7\,200\,000 & 9\,600\,000 \\
\hline
Chiffre d'affaires prévisionnel & 6\,000\,000 & 11\,000\,000 & 17\,000\,000 \\
\hline
Résultat prévisionnel (CA $-$ charges $-$ invest.) & $-2\,300\,000$ & $+2\,800\,000$ & $+5\,900\,000$ \\
\hline
\end{tabularx}
\end{center}

\begin{itemize}
\item \textbf{Formes juridiques OHADA envisageables :} Entreprise Individuelle, SNC (Société en Nom Collectif), SARL (Société à Responsabilité Limitée), SAS (Société par Actions Simplifiée). Le tableau comparatif fourni en ressource précise : capital minimum, responsabilité des associés/actionnaires (limitée aux apports ou illimitée), fiscalité (Impôt sur les Sociétés ou Impôt sur le Revenu), formalisme juridique, possibilité d'entrée d'investisseurs.
\item \textbf{Extrait de presse :} « Cameroun : la fiscalité verte et les défis du financement des PME vertes » — l'article souligne les avantages fiscaux possibles pour les entreprises des énergies renouvelables et la difficulté d'accès au crédit bancaire, qui pousse les PME vertes vers les \emph{Business Angels}.
\end{itemize}

Les trois promoteurs hésitent : Basile veut la simplicité de l'entreprise individuelle, Carine défend la SARL, Amina rêve d'une SAS pour attirer rapidement des investisseurs. Votre groupe est missionné comme \textbf{cabinet-conseil} pour trancher ce dilemme et rédiger la note destinée à un Business Angel.

\courssub{Consignes}

Travail en groupes de 4 élèves (2 heures + 1 heure de restitution orale).

\begin{enumerate}
\item \textbf{Analyser} le projet : montrez que \emph{GreenEnergy Cameroon} est bien une entreprise au sens de la leçon. Précisez son secteur d'activité et la taille visée à 3 ans.
\item \textbf{Choisir} la forme juridique la plus adaptée parmi les quatre proposées. Justifiez votre choix par \textbf{3 arguments juridiques, 2 arguments financiers et 1 argument stratégique}, en vous appuyant sur le tableau comparatif OHADA.
\item \textbf{Organiser} : proposez un organigramme prévisionnel à 3 ans (15 personnes). Identifiez \textbf{2 risques RH majeurs} (par exemple : fuite des techniciens qualifiés, surcharge du gestionnaire) et proposez une parade pour chacun.
\item \textbf{Stratégie} : construisez la matrice SWOT du projet (2 forces, 2 faiblesses, 2 opportunités, 2 menaces minimum). Déduisez-en \textbf{la} stratégie prioritaire : pénétration de marché, stratégie de niche ou partenariat ? Justifiez.
\item \textbf{Rédiger} la \emph{Note de cadrage stratégique} (2 pages maximum) destinée au Business Angel : présentation du projet, choix juridique argumenté, modèle économique, synthèse SWOT, besoin de financement chiffré, impact RSE.
\end{enumerate}

\begin{attention}[Critères d'évaluation de la note (20 points)]
Justification juridique (4 pts) ; cohérence économique et financière (4 pts) ; qualité de l'argumentation et de l'écriture (4 pts) ; intégration de la RSE (3 pts) ; présentation générale (3 pts). Chaque argument non justifié par une donnée du dossier perd la moitié de sa valeur.
\end{attention}

\courssub{Ressources à mobiliser}

\begin{aretenir}[Notions utiles de la leçon]
\begin{itemize}
\item \textbf{Entreprise} : unité économique, juridiquement autonome, qui combine des facteurs de production (travail, capital, technique) pour produire des biens ou services marchands destinés à être vendus sur un marché en vue d'un profit.
\item \textbf{Classification} : par secteur (primaire, secondaire, tertiaire), par taille (TPE, PME, grande entreprise), par forme juridique (droit OHADA : entreprise individuelle, SNC, SARL, SAS, SA).
\item \textbf{Responsabilité} : limitée aux apports (SARL, SAS) ou illimitée et solidaire (SNC, entreprise individuelle).
\item \textbf{Résultat prévisionnel} $=$ Chiffre d'affaires $-$ charges $-$ investissements. Le \textbf{seuil de rentabilité} est atteint quand le résultat devient positif de façon durable.
\item \textbf{Matrice SWOT} : Forces et Faiblesses (internes) ; Opportunités et Menaces (externes).
\item \textbf{RSE} : démarche par laquelle l'entreprise intègre volontairement les préoccupations sociales, environnementales et éthiques dans ses activités.
\end{itemize}

\textbf{Tableau comparatif des formes juridiques OHADA (extrait)}

\begin{center}
\begin{tabularx}{\linewidth}{|l|X|X|X|X|}
\hline
\rowcolor{popblueL} \textbf{Critère} & \textbf{Entreprise Individuelle} & \textbf{SNC} & \textbf{SARL} & \textbf{SAS} \\
\hline
\textbf{Personnalité morale} & Non & Oui & Oui & Oui \\
\hline
\textbf{Nombre d'associés/actionnaires} & 1 exploitant & Min 2 associés & Min 1 associé & Min 1 actionnaire \\
\hline
\textbf{Responsabilité} & Illimitée (biens propres) & Illimitée \& solidaire & Limitée aux apports & Limitée aux apports \\
\hline
\textbf{Capital minimum} & Aucun & Aucun & Libre (fixé par les statuts) & Libre (fixé par les statuts) \\
\hline
\textbf{Fiscalité (défaut)} & IR (BIC) & IR (transparence) & IS & IS \\
\hline
\textbf{Formalisme} & Très simple (déclaration) & Statuts, immatriculation & Statuts, immatriculation, gérant & Statuts sur mesure, Président obligatoire, commissaire aux comptes si seuils franchis \\
\hline
\textbf{Ouverture aux investisseurs} & Impossible & Difficile (agrément) & Possible (augmentation capital, agrément) & Facile (actions, clauses de variabilité) \\
\hline
\end{tabularx}
\end{center}
\end{aretenir}

\courssub{Démarche et corrigé commenté}

\begin{exoresolu}[Corrigé commenté de l'activité « GreenEnergy Cameroon »]
\textbf{Étape 1 — Analyser.} \emph{GreenEnergy Cameroon} réunit les trois caractères de l'entreprise : elle \textbf{combine des facteurs de production} (travail des techniciens, capital de 5 millions FCFA, savoir-faire technique), elle \textbf{produit un service marchand} (installation et maintenance de kits solaires vendus aux ménages et PME) et elle \textbf{vise le profit} (rentabilité à 3 ans). Secteur : \textbf{tertiaire} (services d'installation et de maintenance), avec une composante commerciale (vente de kits). Taille visée à 3 ans : 15 personnes et un CA de 17 millions FCFA, donc une \textbf{PME} (très petite entreprise devenant petite entreprise selon la nomenclature camerounaise).

\tcblower

\textbf{Étape 2 — Choisir la forme juridique : la SARL.}

\emph{Arguments juridiques :} (1) La \textbf{responsabilité est limitée aux apports} : en cas de faillite, les biens personnels des trois jeunes sont protégés, contrairement à l'entreprise individuelle ou à la SNC où la responsabilité est illimitée et solidaire ; (2) La SARL est une \textbf{personne morale} distincte, adaptée à un projet à plusieurs associés, ce qui exclut l'entreprise individuelle (un seul exploitant) ; (3) Son \textbf{formalisme reste simple} (statuts types, registre du commerce OHADA, gérant majoritaire ou non), contrairement à la SAS qui exige des statuts sur mesure, la nomination d'un Président et éventuellement d'un commissaire aux comptes.

\emph{Arguments financiers :} (1) Le capital de 5 millions FCFA est \textbf{pleinement compatible} avec la SARL, dont le capital minimum est librement fixé par les associés (pas de plancher légal depuis la révision 2014) ; (2) La SARL est soumise à l'\textbf{Impôt sur les Sociétés (IS)}, régime clair et crédible face aux banques et investisseurs, et le dossier de presse signale des \textbf{avantages de la fiscalité verte} (exonérations temporaires, crédits d'impôt) accessibles aux sociétés régulièrement constituées.

\emph{Argument stratégique :} La SARL permet l'\textbf{entrée ultérieure d'investisseurs} par augmentation de capital et ouverture du tour de table (cession de parts sociales avec agrément), ce qui correspond à l'objectif de croissance rapide — avec la possibilité de se transformer en SAS plus tard si le projet grandit et nécessite une structure plus flexible pour lever des fonds massifs.

\textbf{Étape 3 — Organiser.} Organigramme prévisionnel à 3 ans (15 personnes) :
\begin{itemize}
\item \textbf{Direction Générale} (Carine, Gérante) ;
\item \textbf{Pôle Technique} (Basile, Responsable Technique) : 6 techniciens installateurs, 2 techniciens SAV/Itinérance ;
\item \textbf{Pôle Commercial \& Marketing} : 3 commerciaux terrain (prospection, relation client, recouvrement) ;
\item \textbf{Pôle Administratif \& Financier} : 1 comptable, 1 assistant administratif/gestion des stocks ;
\item \textbf{Cellule Informatique \& Suivi Client} (Amina, Responsable Innovation) : 1 développeur (application de paiement échelonné, monitoring installations).
\end{itemize}

\emph{Risques RH majeurs et parades :}
\begin{enumerate}
\item \textbf{Fuite des techniciens qualifiés} vers la concurrence (forte demande) $\rightarrow$ \textbf{Parade :} politique de fidélisation (formation continue certifiante, primes de performance trimestrielles, participation aux bénéfices / intéressement, mutuelle santé).
\item \textbf{Surcharge de la Direction} au démarrage (gestion administrative, RH, trésorerie, relation investisseurs) $\rightarrow$ \textbf{Parade :} recrutement du comptable dès l'année 1 (charge prévue au budget), délégation progressive de la gestion administrative à l'assistant, mise en place d'outils de gestion dématérialisés (Amina).
\end{enumerate}

\textbf{Étape 4 — Stratégie : matrice SWOT.}

\begin{center}
\begin{tabularx}{\linewidth}{|X|X|}
\hline
\rowcolor{popgreenL} \textbf{Forces (Internes)} & \textbf{Faiblesses (Internes)} \\
\hline
Équipe fondatrice pluridisciplinaire et qualifiée (technique, gestion, IT) ; Offre différenciante : SAV inclus + paiement échelonné 12 mois sans frais & Capital initial faible (5 M FCFA) limitant le stock et le marketing ; Résultat négatif en année 1 ($-2,3$ M FCFA) nécessitant un apport externe \\
\hline
\rowcolor{poporangeL} \textbf{Opportunités (Externes)} & \textbf{Menaces (Externes)} \\
\hline
Demande solaire croissante (insécurité énergétique) ; Fiscalité verte favorable (exonérations, subventions) ; Appétence des Business Angels pour l'impact & Concurrence chinoise agressive sur le prix (kits bas de gamme) ; Faible pouvoir d'achat des ménages cibles ; Difficulté d'accès au crédit bancaire classique (taux élevés, garanties) \\
\hline
\end{tabularx}
\end{center}

\emph{Stratégie prioritaire :} La \textbf{stratégie de niche différenciée} (qualité de service et financement). Ne pas jouer la guerre des prix face aux kits chinois, mais cibler les ménages urbains et PME sensibles à la fiabilité, en vendant un \emph{service complet} (installation aux normes + maintenance préventive + paiement en 12 mois via application mobile), ce que la concurrence informelle ne peut pas offrir. Le partenariat avec un Business Angel vient financer le déficit de trésorerie de l'année 1 (BFR + investissements) et valider le modèle.

\textbf{Étape 5 — Rédiger (plan de la Note de cadrage stratégique).}
\begin{enumerate}
\item \textbf{Présentation du projet :} Équipe, mission (accès énergie propre), marché cible (Yaoundé/Douala), modèle (vente + maintenance + crédit).
\item \textbf{Choix juridique argumenté :} SARL (6 arguments structurés : responsabilité, personnalité morale, formalisme, capital, fiscalité IS/verte, ouverture capital).
\item \textbf{Modèle économique :} Marge sur kits (achat/vente), revenus récurrents contrats maintenance (20\% CA/an), revenus financiers (intérêts paiement échelonné). Rentabilité opérationnelle dès l'année 2 ($+2,8$ M FCFA).
\item \textbf{Synthèse SWOT \& Stratégie :} Niche « service premium accessible », barrière à l'entrée par la relation client et l'IT.
\item \textbf{Besoin de financement :} \textbf{3\,000\,000 FCFA} (couvre déficit Année 1 2,3 M + BFR sécurité 0,7 M). Contrepartie : 15-20\% du capital (parts sociales SARL), siège au Conseil de surveillance / Assemblée Générale, reporting trimestriel.
\item \textbf{Impact RSE (mesurable) :} \textbf{Environnement :} ~500 kits/an $\rightarrow$ réduction estimée 200 tonnes CO2/an (vs groupes électrogènes). \textbf{Social :} 15 emplois directs qualifiés (jeunes), accès énergie 500 foyers/an (éclairage, réfrigération, productivité). \textbf{Gouvernance :} Parité équipe fondatrice, transparence comptable (IS), éthique (pas de travail dissimulé).
\end{enumerate}
\end{exoresolu}

\courssub{Bilan de l'activité}

Cette activité a permis de vérifier que les notions de la leçon ne sont pas de simples définitions scolaires mais de véritables \textbf{outils de décision}. Les élèves ont d'abord confirmé qu'un projet concret répond à la définition philosophique et économique de l'entreprise : une organisation humaine qui combine des facteurs de production pour créer de la valeur. Ils ont ensuite découvert que le \textbf{choix juridique n'est jamais neutre} : il engage la responsabilité des personnes, la fiscalité et la capacité de croissance — ce qui illustre concrètement la classification des entreprises. L'organigramme et l'analyse des risques RH ont montré que le fonctionnement interne repose avant tout sur le \textbf{facteur humain} et sa gestion. Enfin, la matrice SWOT et la note de cadrage ont révélé le lien entre \textbf{stratégie, gouvernance et RSE} : une entreprise responsable, qui intègre l'impact social et environnemental dans son modèle d'affaires, devient plus crédible et attractive face aux financeurs. Le savoir-faire essentiel — \emph{rédiger et justifier une démarche, une réponse ou une conclusion} — a été pleinement exercé, en préparation directe des épreuves de dissertation et de commentaire au Probatoire et au Baccalauréat technique.

\newpage

\coursec{Méthodes et exercices résolus}

\courssub{Méthodes et Exercices Résolus}

\begin{methode}[Définir et caractériser une entreprise]
    \textbf{Objectif :} Maîtriser la définition juridique et économique pour qualifier toute structure productive.
    \begin{enumerate}
        \item \textbf{Identifier la définition juridique :} Une personne morale (sauf entreprise individuelle, qui est une personne physique), dotée d'un patrimoine propre, immatriculée au \cle{RCCM} (Registre du Commerce et du Crédit Mobilier), soumise au droit \cle{OHADA} (Acte uniforme relatif au droit des sociétés commerciales et du GIE pour les sociétés).
        \item \textbf{Repérer les 3 critères économiques fondamentaux :}
              \begin{itemize}
                  \item \cle{Autonomie de décision} : pouvoir de gestion propre (choix de production, d'achat, de vente).
                  \item \cle{Prise de risque} : support des pertes éventuelles sur le capital engagé.
                  \item \cle{Finalité lucrative} : recherche de profit (sauf associations et ONG, qui poursuivent un but non lucratif).
              \end{itemize}
        \item \textbf{Distinguer les 3 composantes (réflexe Bac) :}
              \begin{itemize}
                  \item \cle{Le capital} (apports : en numéraire, en nature, en industrie).
                  \item \cle{Le travail} (salariés, dirigeant).
                  \item \cle{La technologie / l'organisation} (savoir-faire, processus).
              \end{itemize}
        \item \textbf{Qualifier la nature :} Marchande (vente de biens/services) vs Non marchande (administrations, associations) ; Privée vs Publique (capital majoritairement détenu par l'État).
    \end{enumerate}
    \textbf{Formule à retenir :} Entreprise = \textbf{Unité de production} + \textbf{Autonomie de décision} + \cle{Prise de risque} + \textbf{Finalité lucrative} (généralement).
\end{methode}

\begin{exoresolu}[Caractérisation de l'entreprise « Bois \& Design »]
    \textbf{Énoncé :} M. Ngono crée « Bois \& Design », menuiserie moderne à Douala. Il apporte 5 millions FCFA, achète des machines (10 millions FCFA, grâce à un emprunt bancaire), embauche 3 menuisiers et un comptable. Il gère seul les commandes et la trésorerie.
    \begin{enumerate}
        \item Cette structure est-elle une entreprise au sens économique ? Justifier par 3 arguments.
        \item Quelle est sa forme juridique probable ? Quel texte encadre son immatriculation ?
        \item Identifier les facteurs de production présents.
    \end{enumerate}
    \tcblower
    \textbf{Solution détaillée :}
    \begin{enumerate}
        \item \textbf{OUI, c'est une entreprise.}
              \begin{itemize}
                  \item \textbf{Autonomie de décision :} M. Ngono décide seul de la production, des achats et des ventes (critère de l'unité de décision).
                  \item \textbf{Prise de risque :} Il a investi 5 millions FCFA (apport personnel) et contracté un emprunt de 10 millions FCFA. En cas d'échec, il perd son apport et reste redevable envers la banque (risque sur son patrimoine personnel, car en Entreprise Individuelle les patrimoines sont confondus).
                  \item \textbf{Finalité lucrative :} La menuiserie vise à vendre des meubles pour dégager un excédent (bénéfice) rémunérant le capital et le travail de l'entrepreneur.
              \end{itemize}
        \item \textbf{Forme juridique :} \cle{Entreprise Individuelle (EI)} / entrepreneur individuel. Il n'y a pas de personne morale distincte : le patrimoine de l'entreprise est confondu avec celui de M. Ngono.
              \textbf{Encadrement juridique :} L'immatriculation au \cle{RCCM} est obligatoire (Acte uniforme OHADA portant organisation des procédures simplifiées de recouvrement et des voies d'exécution, et Acte uniforme sur le droit commercial général). \emph{Note : si M. Ngono créait une SARL ou une SA, l'Acte uniforme relatif au droit des sociétés commerciales et du GIE (AUSCGIE) s'appliquerait pleinement.}
        \item \textbf{Facteurs de production :}
              \begin{itemize}
                  \item \cle{Travail} : M. Ngono (travail de direction/gestion) + 3 menuisiers (exécution) + 1 comptable (administration).
                  \item \cle{Capital} : machines (capital fixe), bois et vernis (capital circulant), trésorerie (emprunt + apport).
                  \item \cle{Nature} : le bois (matière première), le local (foncier).
                  \item \cle{Technologie / Organisation} : plans de fabrication, processus de commande/livraison, logiciel comptable.
              \end{itemize}
    \end{enumerate}
    \textbf{Conclusion :} « Bois \& Design » remplit tous les critères de l'entreprise marchande privée individuelle.
\end{exoresolu}

\begin{methode}[Classer une entreprise selon les critères officiels (Cameroun / OHADA)]
    \textbf{Objectif :} Appliquer la nomenclature statistique (INS) et juridique (OHADA) pour situer l'entreprise dans son environnement.
    \begin{enumerate}
        \item \textbf{Critère de la TAILLE (INS / Loi n° 2010/001 du 13 avril 2010 de promotion des PME) :} Calculer l'\cle{Effectif} et le \cle{Chiffre d'Affaires (CA) HT} ou le \cle{Total Bilan}.
              \begin{itemize}
                  \item \cle{TPE (Très Petite Entreprise)} : effectif faible (moins de 10 salariés) et CA modeste.
                  \item \cle{PME (Petite et Moyenne Entreprise)} : effectif intermédiaire et CA plafonné (seuils fixés par la réglementation en vigueur).
                  \item \cle{GE (Grande Entreprise)} : effectif élevé (100 salariés et plus) ou CA supérieur au milliard de FCFA.
              \end{itemize}
              \emph{Attention :} retenir les ordres de grandeur (effectif, CA) plutôt que des seuils chiffrés qui évoluent par textes d'application.
        \item \textbf{Critère du SECTEUR d'activité :}
              \begin{itemize}
                  \item \cle{Primaire} : agriculture, pêche, forêt, mines.
                  \item \cle{Secondaire} : industrie de transformation, BTP, énergie, eau.
                  \item \cle{Tertiaire} : commerce, transport, services, finance.
              \end{itemize}
        \item \textbf{Critère JURIDIQUE (Acte uniforme OHADA) :}
              \begin{itemize}
                  \item \cle{Personne physique} : Entreprise Individuelle (EI).
                  \item \cle{Personnes morales} : SARL (responsabilité limitée aux apports), SA (capital divisé en actions, minimum 10 millions FCFA), SNC, SCS, GIE, Coopérative (régi par l'Acte uniforme relatif au droit des sociétés coopératives).
              \end{itemize}
        \item \textbf{Critère du CAPITAL / PROPRIÉTÉ :} Privée (nationaux/étrangers), Publique (État/collectivités majoritaire), Mixte.
    \end{enumerate}
    \textbf{Réflexe Bac :} Toujours croiser Taille + Secteur + Forme juridique pour une classification complète.
\end{methode}

\begin{exoresolu}[Classification croisée d'entreprises camerounaises]
    \textbf{Énoncé :} Classer les structures suivantes en justifiant \textbf{Taille}, \textbf{Secteur} et \textbf{Forme juridique} :
    \begin{enumerate}
        \item \textbf{SOCAPALM} : plusieurs milliers d'employés, palmeraies et usines d'huilerie, société anonyme cotée en bourse, actionnariat mixte (groupe privé majoritaire, État actionnaire).
        \item \textbf{« Chez Mama » (restaurant à Yaoundé)} : 4 employés, CA 30 millions FCFA, propriétaire Mme Mendo (patrimoine confondu).
        \item \textbf{MTN Cameroon} : plus d'un millier d'employés, télécommunications, SA de droit camerounais, actionnariat mixte (MTN Group + public).
        \item \textbf{GIC « Paysans Unis »} : 50 agriculteurs, mutualisation du matériel et vente collective du maïs, but non lucratif principal.
    \end{enumerate}
    \tcblower
    \textbf{Solution détaillée :}
    \begin{center}
    \begin{tabularx}{\linewidth}{|l|X|X|X|}\hline
    \rowcolor{popblueL}
    \textbf{Structure} & \textbf{Taille} & \textbf{Secteur} & \textbf{Forme juridique} \\ \hline
    \textbf{SOCAPALM} & \textbf{GE} (effectif et CA très élevés) & \textbf{Primaire} (culture des palmiers à huile) + \textbf{Secondaire} (usines d'huileries) $\to$ \emph{agro-industrie} & \textbf{SA (Société Anonyme)} : capital divisé en actions, responsabilité limitée, cotation en bourse, actionnariat mixte. \\ \hline
    \textbf{Chez Mama} & \textbf{TPE} (effectif $< 10$, CA faible) & \textbf{Tertiaire} (restauration = services) & \textbf{Entreprise Individuelle (EI)} : personne physique, patrimoine unique, pas de capital social distinct. \\ \hline
    \textbf{MTN Cameroon} & \textbf{GE} & \textbf{Tertiaire} (télécommunications / services numériques) & \textbf{SA} : forme adaptée aux grandes structures, capital en actions, responsabilité limitée aux apports. \\ \hline
    \textbf{GIC Paysans Unis} & \textbf{PME} (effectif $< 100$) & \textbf{Primaire} (agriculture) & \textbf{Coopérative} (Acte uniforme OHADA relatif au droit des sociétés coopératives) ou \textbf{GIE}. \emph{Note : le GIC camerounais se rapproche juridiquement du GIE ou de la coopérative OHADA.} \\ \hline
    \end{tabularx}
    \end{center}
    \textbf{Conclusion :} La classification permet d'identifier les obligations légales (comptabilité, fiscalité, gouvernance) et les politiques publiques ciblées (ex. : dispositifs d'appui aux PME pour « Chez Mama »).
\end{exoresolu}

\begin{methode}[Analyser le fonctionnement interne : facteurs de production et structure organisationnelle]
    \textbf{Objectif :} Expliquer comment l'entreprise combine les ressources pour créer de la valeur (fonction production).
    \begin{enumerate}
        \item \textbf{Identifier les facteurs de production :}
              \begin{itemize}
                  \item \cle{Travail (L)} : quantité (effectifs, heures) + qualité (compétences, formation, motivation). Coût = salaire brut + charges sociales patronales.
                  \item \cle{Capital (K)} : fixe (machines, bâtiments, logiciels — amortissables) + circulant (stocks, créances, trésorerie — rotation courte). Coût = intérêts (emprunt) + amortissements + dividendes (fonds propres).
                  \item \cle{Nature / Matières premières (M)} : achats consommés (coût variable direct).
                  \item \cle{Organisation / Progrès technique} : la combinaison productive s'écrit $Y = f(K, L, M)$. \cle{Productivité} = $\dfrac{Y}{\text{Facteur utilisé}}$ (ex. : productivité horaire du travail = production / heures travaillées).
              \end{itemize}
        \item \textbf{Analyser la structure organisationnelle (organigramme) :}
              \begin{itemize}
                  \item \cle{Niveaux hiérarchiques} : stratégique (dirigeant/conseil d'administration), tactique (cadres intermédiaires), opérationnel (exécution).
                  \item \cle{Types de structures} : fonctionnelle (par métier : RH, finances, production, commercial), divisionnelle (par produit ou région), par projet.
                  \item \cle{Indicateurs} : \cle{étendue du contrôle} (nombre de subordonnés par supérieur), \cle{centralisation/décentralisation}, \cle{formalisation} (procédures écrites).
              \end{itemize}
        \item \textbf{Lier facteurs et structure :} la structure doit optimiser la combinaison capital/travail et la coordination des tâches.
    \end{enumerate}
    \textbf{Formules clés :} \cle{Productivité apparente du travail} $= \dfrac{\text{Valeur ajoutée}}{\text{Effectif}}$ ; \cle{Taux d'encadrement} $= \dfrac{\text{Cadres}}{\text{Effectif total}}$.
\end{methode}

\begin{exoresolu}[Diagnostic organisationnel et productivité — Cas « Bois \& Design »]
    \textbf{Énoncé :} « Bois \& Design » (M. Ngono, 3 menuisiers, 1 comptable) produit 50 meubles par mois. CA mensuel = 2 500 000 FCFA. Masse salariale (4 salariés, charges comprises) = 600 000 FCFA. Amortissements machines = 100 000 FCFA. Achats bois/consommables = 800 000 FCFA. Loyer/énergie = 150 000 FCFA. Intérêts de l'emprunt = 50 000 FCFA.
    \begin{enumerate}
        \item Calculer la \cle{Valeur Ajoutée (VA)} et le \cle{Résultat Net} (avant impôt).
        \item Calculer la \cle{Productivité Apparente du Travail (PAT)} = VA / Effectif total.
        \item Proposer un organigramme fonctionnel et identifier le risque de \cle{conflit d'agence}.
    \end{enumerate}
    \tcblower
    \textbf{Solution détaillée :}
    \begin{enumerate}
        \item \textbf{Compte de résultat simplifié :}
              \begin{align*}
                  \text{CA} &= 2\,500\,000 \\
                  \text{Consommations intermédiaires (achats + loyer/énergie)} &= 800\,000 + 150\,000 = 950\,000 \\
                  \textbf{VA} &= \text{CA} - \text{CI} = 2\,500\,000 - 950\,000 = \textbf{1\,550\,000 FCFA} \\
                  \text{Charges de personnel} &= 600\,000 \\
                  \text{Amortissements} &= 100\,000 \\
                  \text{Excédent d'exploitation} &= 1\,550\,000 - 600\,000 - 100\,000 = 850\,000 \\
                  \text{Charges financières (intérêts)} &= 50\,000 \\
                  \textbf{Résultat net (avant impôt)} &= 850\,000 - 50\,000 = \textbf{800\,000 FCFA}
              \end{align*}
        \item \textbf{Productivité Apparente du Travail (PAT) :}
              Effectif total = 1 (dirigeant) + 3 (menuisiers) + 1 (comptable) = 5 personnes.
              \[ \text{PAT} = \frac{1\,550\,000}{5} = \textbf{310\,000 FCFA / personne / mois}. \]
              \emph{Interprétation :} chaque personne génère 310 000 FCFA de valeur ajoutée. La part des salaires dans la VA est de $\dfrac{600\,000}{1\,550\,000} \approx 38,7\,\%$, ce qui est raisonnable pour une PME.
        \item \textbf{Organigramme fonctionnel proposé :}
              \begin{popfigure}
              \begin{center}
              \begin{tikzpicture}[node distance=1.2cm, every node/.style={draw, rounded corners, fill=popblueL, minimum width=2.8cm, minimum height=0.9cm, align=center, font=\small}]
              \node[align=center] (dir){Dirigeant\\(M. Ngono)};
              \node (prod) [below left=1cm and 0.5cm of dir, align=center]{Production\\3 menuisiers};
              \node (admin) [below right=1cm and 0.5cm of dir, align=center]{Administration \& Finances\\1 comptable};
              \node (com) [below=2.6cm of dir, align=center]{Commercial\\(assuré par le dirigeant)};
              \draw[->, thick, popblue] (dir) -- (prod);
              \draw[->, thick, popblue] (dir) -- (admin);
              \draw[->, thick, popblue] (dir) -- (com);
              \end{tikzpicture}
              \end{center}
              \legende{Organigramme fonctionnel simplifié de « Bois \& Design » : structure centralisée autour du dirigeant unique.}
              \end{popfigure}
              \textbf{Risque d'agence :} M. Ngono est à la fois \textbf{propriétaire} et \textbf{gérant}.
              \begin{itemize}
                  \item \cle{Conflit actionnaires/dirigeants :} inexistant ici (les deux fonctions sont confondues).
                  \item \cle{Conflit créanciers/propriétaire :} risque réel. M. Ngono peut prendre des risques excessifs (investissements hasardeux, prélèvements personnels sur la trésorerie) alors que la banque supporte une partie du risque de perte. \textbf{Mécanismes de contrôle :} clauses du contrat de crédit, exigence de comptes annuels réguliers.
              \end{itemize}
    \end{enumerate}
    \textbf{Conclusion :} L'entreprise est rentable (résultat net $> 0$) et sa productivité est correcte pour de l'artisanat modernisé. La structure, simple mais très centralisée, fait peser un risque de surcharge sur le dirigeant.
\end{exoresolu}

\begin{methode}[Réaliser un diagnostic stratégique simplifié (SWOT) et évaluer la RSE]
    \textbf{Objectif :} Structurer l'analyse de l'environnement et des ressources internes pour définir une stratégie, en intégrant la RSE (norme ISO 26000).
    \begin{enumerate}
        \item \textbf{Analyse EXTERNE (environnement) :}
              \begin{itemize}
                  \item \cle{Opportunités (O)} : marché en croissance, appui de l'État aux PME, innovation technique, nouveaux segments de clientèle.
                  \item \cle{Menaces (T)} : concurrence (importations, secteur informel), réglementation (droit du travail, fiscalité, environnement), pouvoir de négociation des clients/fournisseurs, instabilité économique (inflation, change).
              \end{itemize}
        \item \textbf{Analyse INTERNE (ressources et compétences) :}
              \begin{itemize}
                  \item \cle{Forces (S)} : savoir-faire unique, marque, trésorerie saine, main-d'œuvre qualifiée, localisation.
                  \item \cle{Faiblesses (W)} : trésorerie tendue, dépendance à un client ou un fournisseur, matériel vétuste, manque de compétences managériales, absence de certification.
              \end{itemize}
        \item \textbf{Matrice SWOT $\rightarrow$ Stratégies :}
              \begin{itemize}
                  \item \cle{SO (offensive)} : utiliser les forces pour saisir les opportunités (investir, exporter).
                  \item \cle{WO (renforcement)} : corriger les faiblesses pour saisir les opportunités (formation, emprunt, partenariat).
                  \item \cle{ST (défensive)} : utiliser les forces pour parer les menaces (fidélisation, diversification).
                  \item \cle{WT (survie)} : réduire les faiblesses face aux menaces (restructuration, cession).
              \end{itemize}
        \item \textbf{RSE (Responsabilité Sociétale des Entreprises) — les 7 questions centrales d'ISO 26000 :}
              \begin{enumerate}
                  \item Gouvernance de l'organisation (éthique, transparence).
                  \item Droits de l'homme (non-discrimination, refus du travail forcé et du travail des enfants).
                  \item Relations et conditions de travail (santé/sécurité, dialogue social, formation).
                  \item Environnement (pollution, économie circulaire, climat).
                  \item Loyauté des pratiques (lutte contre la corruption, concurrence loyale).
                  \item Questions relatives aux consommateurs (sécurité des produits, réclamations).
                  \item Communautés et développement local (emploi local, achats locaux, mécénat).
              \end{enumerate}
    \end{enumerate}
    \textbf{Réflexe Bac :} Ne pas lister O/T/S/W sans lien avec le cas étudié. Toujours conclure par \textbf{UNE} orientation stratégique claire (ex. : « spécialisation haut de gamme + certification qualité »).
\end{methode}

\begin{exoresolu}[SWOT et RSE — « Bois \& Design » face à l'importation de meubles bon marché]
    \textbf{Énoncé :} Le marché du meuble à Douala est envahi par des meubles importés bon marché (prix inférieurs d'environ 30\,\%). M. Ngono souhaite définir sa stratégie 2025-2027. Données : savoir-faire local reconnu (bois durs, essences locales), machines semi-automatisées, 3 menuisiers qualifiés, trésorerie positive (800 000 FCFA), pas de site web, pas de certification, local exigu (200 m\textsuperscript{2}).
    \begin{enumerate}
        \item Construire la matrice SWOT (3 éléments par case minimum).
        \item Déduire la stratégie prioritaire (SO, WO, ST ou WT) et 2 actions concrètes.
        \item Identifier 2 actions RSE prioritaires (ISO 26000) pour se différencier.
    \end{enumerate}
    \tcblower
    \textbf{Solution détaillée :}
    \begin{enumerate}
        \item \textbf{Matrice SWOT « Bois \& Design » :}
              \begin{center}
              \begin{tabularx}{\linewidth}{|X|X|}\hline
              \rowcolor{popgreenL} \multicolumn{2}{|c|}{\textbf{ANALYSE INTERNE}} \\ \hline
              \textbf{FORCES (S)} & \textbf{FAIBLESSES (W)} \\ \hline
              1. Maîtrise des essences locales (padouk, ayous, iroko), durables et esthétiques. & 1. Capacité de production limitée (local de 200 m\textsuperscript{2}, machines semi-automatiques). \\
              2. Équipe stable, qualifiée et polyvalente. & 2. Absence de visibilité numérique (pas de site web, pas de réseaux sociaux professionnels). \\
              3. Trésorerie saine (autofinancement possible) et pas de dettes fournisseurs. & 3. Dépendance au dirigeant unique (commercial + gestion + technique = goulot d'étranglement). \\ \hline
              \rowcolor{poporangeL} \multicolumn{2}{|c|}{\textbf{ANALYSE EXTERNE}} \\ \hline
              \textbf{OPPORTUNITÉS (O)} & \textbf{MENACES (T)} \\ \hline
              1. Demande croissante de meubles « sur mesure / éco-responsables » (classes moyennes, hôtels, entreprises). & 1. Concurrence des prix des importations (production industrielle à bas coût). \\
              2. Programmes publics d'appui aux PME (guichets uniques, dispositifs de certification). & 2. Hausse du prix du bois local (pression sur la ressource forestière). \\
              3. Développement du commerce en ligne et de la livraison urbaine. & 3. Risque de départ des menuisiers qualifiés (ateliers concurrents). \\ \hline
              \end{tabularx}
              \end{center}
        \item \textbf{Stratégie prioritaire : \cle{SO (offensive) — différenciation par la qualité et le service.}}
              \emph{Justification :} les forces (bois noble, compétence) répondent exactement à l'opportunité (demande de sur-mesure éco-responsable). Jouer sur les prix serait perdant face aux importations.
              \textbf{Actions concrètes :}
              \begin{enumerate}
                  \item \textbf{Lancer une gamme « Signature Locale » certifiée :} approvisionnement en bois légal (traçabilité), label de qualité, fiche de traçabilité de l'arbre au meuble. Cibles : hôtels, bureaux, diaspora.
                  \item \textbf{Digitaliser la commercialisation :} création d'un site vitrine avec catalogue, commande en ligne et acompte par mobile money ; délégation progressive de la fonction commerciale.
              \end{enumerate}
        \item \textbf{Actions RSE prioritaires (ISO 26000) :}
              \begin{enumerate}
                  \item \textbf{Environnement :} \cle{économie circulaire} — valorisation des chutes de bois (briquettes de chauffage, panneaux, dons aux écoles de menuiserie) ; argument commercial du bois local face au bois importé.
                  \item \textbf{Relations de travail et communauté :} \cle{formation} — partenariat avec un centre de formation professionnelle pour accueillir 2 apprentis par an (insertion des jeunes du quartier) ; \cle{santé/sécurité} — équipements de protection individuelle obligatoires (bruit, poussières), visite médicale annuelle.
              \end{enumerate}
    \end{enumerate}
    \textbf{Conclusion :} La stratégie de \emph{spécialisation différenciante}, appuyée sur la RSE (traçabilité du bois, formation locale), transforme la menace des « prix bas importés » en argument de « valeur durable locale ».
\end{exoresolu}

\begin{methode}[Rédiger une dissertation philosophique (épreuve du Bac technique)]
    \textbf{Objectif :} Structurer une réponse rigoureuse à un sujet du type « L'entreprise a-t-elle une responsabilité seulement économique ? » ou « Le profit est-il le seul but de l'entreprise ? ».
    \begin{enumerate}
        \item \textbf{ANALYSE DU SUJET (15 min) :}
              \begin{itemize}
                  \item Définir les termes : \cle{entreprise} (unité de production, autonomie, risque), \cle{responsabilité} (obligation de rendre compte de ses actes), \cle{économique} (rentabilité, création de valeur, emploi), \cle{seulement} (exclusif, restrictif).
                  \item Identifier la \cle{problématique} : tension entre la \cle{finalité lucrative} (condition de survie, rémunération des facteurs) et la \cle{finalité sociale et sociétale} (légitimité, durabilité, parties prenantes). \textbf{Question directrice :} la viabilité économique est-elle une condition \emph{suffisante} de la responsabilité de l'entreprise ?
              \end{itemize}
        \item \textbf{PLAN DIALECTIQUE (3 parties, 2 à 3 arguments chacune) :}
              \begin{itemize}
                  \item \textbf{Thèse (OUI, primauté de l'économique) :} 1. Contrainte de survie (concurrence, marché) $\to$ le profit est condition \emph{sine qua non}. 2. Thèse de Friedman : la responsabilité de l'entreprise est d'augmenter ses profits dans le cadre légal. 3. La création de richesse est la base de la redistribution (salaires, impôts, investissement).
                  \item \textbf{Antithèse (NON, responsabilité élargie / RSE) :} 1. Externalités négatives (pollution, santé) non payées par le marché $\to$ devoir de prévention et de réparation. 2. Théorie des parties prenantes (Freeman) : clients, salariés, fournisseurs, territoire, générations futures ont des droits. 3. Performance durable : la RSE est un facteur de compétitivité (marque, innovation, fidélisation, accès aux financements).
                  \item \textbf{Synthèse (dépassement : l'entreprise citoyenne) :} 1. La responsabilité économique est \cle{condition nécessaire mais non suffisante}. 2. Nouveau contrat social : raison d'être inscrite dans les statuts, reporting extra-financier. 3. Gouvernance partagée : association des parties prenantes aux décisions.
              \end{itemize}
        \item \textbf{RÉDACTION (1 h 30) :} introduction (accroche, définitions, problématique, annonce du plan), développement (arguments \textbf{illustrés par des exemples camerounais}), conclusion (réponse synthétique, ouverture).
    \end{enumerate}
    \textbf{Critères d'évaluation :} pertinence et précision des concepts, rigueur logique, qualité des exemples (Cameroun), langage philosophique, orthographe et syntaxe.
\end{methode}

\begin{exoresolu}[Dissertation type Bac : « L'entreprise a-t-elle une responsabilité seulement économique ? »]
    \textbf{Énoncé :} Traiter le sujet en suivant la méthode ci-dessus. Durée : 2 h.
    \tcblower
    \textbf{Plan détaillé et arguments clés (squelette corrigé) :}

    \textbf{INTRODUCTION}
    \begin{itemize}
        \item \textbf{Accroche :} « Il n'y a qu'une seule responsabilité sociale de l'entreprise : utiliser ses ressources pour augmenter ses profits » (M. Friedman, 1970) — face à l'idée que l'entreprise est une communauté de travail au service du bien commun.
        \item \textbf{Définitions :} Entreprise = organisation combinant des facteurs de production pour produire des biens et services marchands en supportant des risques. Responsabilité = obligation juridique et morale de répondre de ses actes devant les tiers. Économique = relatif à la production, à l'échange et à la richesse.
        \item \textbf{Problématique :} si la viabilité économique (le profit) est la condition \emph{sine qua non} de l'existence de l'entreprise, suffit-elle à épuiser sa responsabilité ? Ne doit-elle pas aussi répondre des impacts de son action sur ses \cle{parties prenantes} et sur l'environnement ?
        \item \textbf{Annonce du plan :} nous verrons que la responsabilité économique est fondatrice (I), mais insuffisante face aux externalités et aux attentes légitimes (II), ce qui appelle une refondation de l'entreprise comme acteur social responsable (III).
    \end{itemize}

    \textbf{DÉVELOPPEMENT}

    \textbf{I. La responsabilité économique : condition nécessaire et fondatrice de l'entreprise.}
    \begin{enumerate}
        \item \textbf{Contrainte de survie et efficience :} dans l'économie de marché camerounaise (zone CEMAC, concurrence des importations), l'entreprise qui ne dégage pas de profit disparaît. \emph{Ex. :} difficultés des entreprises textiles camerounaises face aux importations de friperie et aux produits asiatiques. Le profit rémunère le risque du capital, finance l'innovation et assure l'emploi.
        \item \textbf{Thèse de Friedman :} le dirigeant est le mandataire des actionnaires ; sa responsabilité est de maximiser la valeur de l'entreprise dans le respect du cadre légal. Toute dépense « sociale » non rentable se fait au détriment des actionnaires ou des salariés.
        \item \textbf{La création de valeur est déjà sociale :} créer des emplois décents, payer des impôts (budget national), former des compétences, développer la sous-traitance locale. \emph{Ex. :} les grandes entreprises agro-industrielles camerounaises emploient des milliers de salariés et structurent des filières locales.
    \end{enumerate}

    \textbf{II. L'insuffisance de la seule responsabilité économique : externalités et parties prenantes.}
    \begin{enumerate}
        \item \textbf{Défaillances de marché / externalités négatives :} le prix de marché n'intègre pas le coût de la pollution (air, eau, déforestation) ni celui des atteintes à la santé (poussières dans les menuiseries, produits chimiques en agriculture). \emph{Ex. :} les rejets industriels dans le Wouri à Douala : le coût est supporté par la collectivité et les riverains, non par le pollueur. D'où le principe \cle{pollueur-payeur} et le principe de précaution.
        \item \textbf{Théorie des parties prenantes (Freeman, 1984) :} l'entreprise a des devoirs envers \emph{tous} ceux qu'elle affecte : salariés (sécurité, salaire décent, dialogue social), clients (sécurité des produits), fournisseurs (délais de paiement), communauté locale (emploi des jeunes, infrastructures), générations futures (climat, biodiversité). \emph{Ex. :} les conflits sociaux récurrents dans les transports au Cameroun traduisent une responsabilité sociale négligée (conditions de travail, partage de la valeur).
        \item \textbf{Risque réputationnel et légal croissant :} les exigences de RSE progressent dans la réglementation et dans les appels d'offres ; ignorer la RSE expose à des contentieux, à la perte de marchés et à des difficultés de financement (les banques intègrent des critères environnementaux, sociaux et de gouvernance).
    \end{enumerate}

    \textbf{III. Vers l'entreprise responsable : intégrer la finalité économique dans une raison d'être sociétale.}
    \begin{enumerate}
        \item \textbf{Dépassement du dilemme : la performance globale.} De nombreuses études montrent une corrélation positive entre RSE et performance financière à long terme. L'\cle{entreprise à mission} inscrit une \cle{raison d'être} dans ses statuts (ex. : « produire du meuble durable en valorisant la forêt camerounaise et la jeunesse locale »), avec un comité de mission et une vérification indépendante.
        \item \textbf{Gouvernance partagée et reporting :} association des parties prenantes aux instances de décision ; publication d'indicateurs extra-financiers (émissions, égalité hommes-femmes, formation, éthique des achats). \emph{Ex. :} les plans de relance des grandes entreprises publiques camerounaises intègrent reboisement, écoles et santé pour les villages riverains, condition de leur acceptabilité sociale.
        \item \textbf{Responsabilité du dirigeant : de « gestionnaire » à « leader éthique ».} Le devoir de loyauté s'élargit à l'intérêt social de l'entreprise et de ses parties prenantes : formation à l'éthique, codes de conduite, dispositifs d'alerte.
    \end{enumerate}

    \textbf{CONCLUSION}
    \begin{itemize}
        \item \textbf{Bilan :} la responsabilité économique est le \textbf{socle} (sans profit, pas d'entreprise, pas d'emploi, pas d'impôt), mais elle n'en est pas le \textbf{plafond}. L'entreprise n'est pas une simple « machine à profit » : c'est une \cle{institution sociale} dotée d'un pouvoir (technique, économique, normatif).
        \item \textbf{Réponse :} non, sa responsabilité est \textbf{plurielle} : économique, sociale, sociétale et environnementale. Elle est \cle{globale} et \cle{intégrée}.
        \item \textbf{Ouverture :} l'avenir appartient peut-être aux entreprises « régénératives », qui restituent plus qu'elles ne prélèvent (carbone, biodiversité, capital social). Comment le droit OHADA et le Code du travail camerounais feront-ils évoluer la définition légale de l'« intérêt social » pour y intégrer la \cle{raison d'être} ?
    \end{itemize}
\end{exoresolu}

\begin{methode}[Réaliser un diagnostic global simplifié (TP intégré) — Approche « Bois \& Design »]
    \textbf{Objectif :} Synthétiser les acquis de la leçon en une fiche de diagnostic standardisée (outil d'aide à la décision).
    \begin{enumerate}
        \item \textbf{Identité et modèle économique :}
              \begin{itemize}
                  \item \cle{Proposition de valeur} : quoi ? pour qui ? quelle différenciation ?
                  \item \cle{Segments de clientèle} : B2B (hôtels, bureaux), B2C (particuliers haut de gamme), export.
                  \item \cle{Flux de revenus} : vente de meubles, prestations d'agencement, maintenance.
                  \item \cle{Structure de coûts} : charges fixes (loyer, salaires, amortissements, assurances) vs charges variables (bois, quincaillerie, sous-traitance, énergie, transport).
                  \item \cle{Seuil de rentabilité (point mort)} : $CA_{PM} = \dfrac{\text{Charges Fixes}}{\text{Taux de marge sur coût variable}}$, avec \cle{Taux de marge} $= \dfrac{\text{Prix de vente} - \text{Coût variable unitaire}}{\text{Prix de vente}}$.
              \end{itemize}
        \item \textbf{Diagnostic financier :}
              \begin{itemize}
                  \item \cle{Rentabilité} : marge brute, excédent brut d'exploitation (EBE), résultat net, rentabilité des capitaux propres (ROE = Résultat net / Capitaux propres).
                  \item \cle{Solvabilité / structure} : endettement (Dettes / Capitaux propres), autonomie financière (Capitaux propres / Total bilan), liquidité.
                  \item \cle{Trésorerie} : plan de trésorerie mensuel (encaissements vs décaissements), Besoin en Fonds de Roulement (BFR = Stocks + Créances clients - Dettes fournisseurs).
              \end{itemize}
        \item \textbf{Diagnostic RH et organisation :} effectifs, compétences clés (risque de perte), climat social (turn-over, absentéisme, accidents), plan de formation.
        \item \textbf{Diagnostic commercial :} parts de marché, fidélisation (taux de réachat), satisfaction client, digitalisation.
        \item \textbf{Synthèse SWOT $\rightarrow$ Plan d'actions (feuille de route) :} objectifs SMART (Spécifiques, Mesurables, Atteignables, Réalistes, Temporellement définis) et indicateurs de suivi.
    \end{enumerate}
    \textbf{Outils :} tableur pour les ratios, graphiques d'évolution du CA, des charges et de la trésorerie.
\end{methode}

\begin{exoresolu}[TP intégré : diagnostic complet « Bois \& Design » — Année N]
    \textbf{Énoncé :} Données de l'année N (en FCFA) :
    \begin{itemize}
        \item CA : 30 000 000 (500 meubles $\times$ 60 000).
        \item Achats de matières (bois, quincaillerie) : 12 000 000 (variable).
        \item Sous-traitance (vernis, tapisserie) : 3 000 000 (variable).
        \item Salaires + charges (5 personnes, gérant non rémunéré) : 9 000 000 (fixe).
        \item Loyer / énergie / assurance / comptable : 3 600 000 (fixe).
        \item Amortissements : 1 200 000 (fixe).
        \item Intérêts de l'emprunt : 600 000.
        \item Impôt sur les sociétés (taux retenu pour l'exercice : 30\,\%) : à calculer.
        \item \textbf{Bilan simplifié fin N :} Actif immobilisé net 15 M / Stocks 4 M / Créances clients 3 M / Banque 2 M. // Capitaux propres 10 M / Emprunt bancaire 8 M / Dettes fournisseurs 4 M / Dettes fiscales et sociales 2 M.
    \end{itemize}
    \begin{enumerate}
        \item Calculer : marge brute, taux de marge, EBE, résultat net, capacité d'autofinancement (CAF).
        \item Calculer le \cle{seuil de rentabilité} (en CA et en quantité).
        \item Calculer : BFR, FRNG, trésorerie nette. Interpréter.
        \item Calculer : autonomie financière, endettement, ROE.
        \item Formuler 3 recommandations stratégiques chiffrées pour N+1.
    \end{enumerate}
    \tcblower
    \textbf{Solution détaillée :}
    \begin{enumerate}
        \item \textbf{Compte de résultat et agrégats :}
              \begin{align*}
                  \text{CA} &= 30\,000\,000 \\
                  \text{Achats de matières premières} &= 12\,000\,000 \\
                  \text{Sous-traitance} &= 3\,000\,000 \\
                  \textbf{Marge brute} &= 30\,000\,000 - 15\,000\,000 = \textbf{15\,000\,000} \\
                  \textbf{Taux de marge} &= \frac{15\,000\,000}{30\,000\,000} = \textbf{50\,\%} \\
                  \text{Charges fixes (hors amortissements et intérêts)} &= 9\,000\,000 + 3\,600\,000 = 12\,600\,000 \\
                  \textbf{EBE} &= 15\,000\,000 - 12\,600\,000 = \textbf{2\,400\,000} \\
                  \text{Amortissements} &= 1\,200\,000 \\
                  \text{Résultat d'exploitation} &= 2\,400\,000 - 1\,200\,000 = 1\,200\,000 \\
                  \text{Intérêts} &= 600\,000 \\
                  \text{Résultat courant avant impôt} &= 1\,200\,000 - 600\,000 = 600\,000 \\
                  \text{Impôt (30\,\%)} &= 180\,000 \\
                  \textbf{Résultat net} &= 600\,000 - 180\,000 = \textbf{420\,000} \\
                  \textbf{CAF} &= \text{Résultat net} + \text{Amortissements} = 420\,000 + 1\,200\,000 = \textbf{1\,620\,000}
              \end{align*}
        \item \textbf{Seuil de rentabilité (point mort) :}
              \begin{itemize}
                  \item Charges fixes totales = salaires et charges + loyer/énergie/assurances + amortissements = 9 M + 3,6 M + 1,2 M = \textbf{13 800 000 FCFA}.
                  \item Taux de marge sur coût variable = 50\,\%.
                  \item $CA_{PM} = \dfrac{13\,800\,000}{0,50} = \textbf{27 600 000 FCFA}$.
                  \item $Q_{PM} = \dfrac{27\,600\,000}{60\,000} = \textbf{460 meubles}$.
                  \item \textbf{Marge de sécurité} = CA réel - $CA_{PM}$ = 30 M - 27,6 M = 2,4 M, soit 8\,\% du CA. \emph{Très faible : risque élevé en cas de baisse des ventes.}
              \end{itemize}
        \item \textbf{Équilibre financier et trésorerie :}
              \begin{itemize}
                  \item \textbf{BFR} = Stocks (4 M) + Créances clients (3 M) - Dettes fournisseurs (4 M) = \textbf{3 000 000 FCFA} (besoin de financement du cycle d'exploitation).
                  \item \textbf{FRNG} = Capitaux propres (10 M) + Emprunt à long terme (8 M) - Actif immobilisé net (15 M) = \textbf{3 000 000 FCFA} (ressources stables restant après financement des investissements).
                  \item \textbf{Trésorerie nette (TN)} = FRNG - BFR = 3 M - 3 M = \textbf{0 FCFA} (cohérent avec : Banque 2 M - Dettes fiscales et sociales 2 M = 0).
                  \item \textbf{Interprétation :} équilibre précaire (TN = 0). Tout le FRNG est absorbé par le BFR : aucune marge de manœuvre pour absorber un retard de paiement client ou une hausse des stocks. \textbf{Urgence :} négocier les délais fournisseurs, relancer les clients, réduire les stocks.
              \end{itemize}
        \item \textbf{Ratios de structure et de rentabilité :}
              \begin{itemize}
                  \item \textbf{Autonomie financière} = Capitaux propres / Total bilan = 10 M / 24 M = \textbf{41,7\,\%} ($> 30\,\%$ : structure saine).
                  \item \textbf{Endettement} = Dettes / Capitaux propres = (8 M + 4 M + 2 M) / 10 M = \textbf{140\,\%} (élevé, mais les dettes à long terme sont majoritaires).
                  \item \textbf{ROE} = Résultat net / Capitaux propres = 420 000 / 10 M = \textbf{4,2\,\%} (faible : la rémunération du capital est inférieure à celle d'un placement sans risque).
              \end{itemize}
        \item \textbf{Recommandations N+1 (chiffrées) :}
              \begin{enumerate}
                  \item \textbf{Augmenter les prix et la valeur :} objectif +10\,\% de prix moyen (gamme « Signature ») $\to$ CA de 33 M à volume constant ; nouveau taux de marge $\approx 55\,\%$ ; nouveau point mort $\approx$ 25 M ; marge de sécurité $\approx 25\,\%$ ; ROE en hausse.
                  \item \textbf{Réduire le BFR :} objectif stocks 3 M (gestion en flux tendus), créances clients 2 M (acompte de 30\,\% à la commande, relance à J+30), dettes fournisseurs 5 M (négociation à 60 jours). BFR cible $\approx$ 0 ; trésorerie dégagée $\approx$ +3 M.
                  \item \textbf{Investissement productif :} achat d'une machine à commande numérique d'occasion (5 M, financé par la CAF de 1,6 M + emprunt de 3,4 M sur 3 ans). Gain de productivité attendu : +30\,\% (650 meubles/an), sous condition d'un carnet de commandes signé au préalable.
              \end{enumerate}
    \end{enumerate}
    \textbf{Conclusion du diagnostic :} entreprise structurellement saine (autonomie financière de 41,7\,\%) mais \textbf{sous-performante} (ROE de 4,2\,\%, marge de sécurité de 8\,\%, trésorerie nulle). Leviers principaux : \textbf{montée en gamme (prix/valeur)} et \textbf{gestion rigoureuse du BFR}.
\end{exoresolu}

\begin{methode}[Appliquer les notions à l'économie informelle et au quotidien (transfert de concepts)]
    \textbf{Objectif :} Démontrer que les concepts (entreprise, facteurs, stratégie, RSE) s'appliquent au secteur informel, majoritaire dans l'emploi au Cameroun.
    \begin{enumerate}
        \item \textbf{Identifier l'unité économique :} le \cle{bayam-sellam} (commerçant ambulant), la \cle{call-box}, l'\cle{atelier de couture ou de soudure}, le \cle{taximan} (moto ou voiture) sont des \textbf{entreprises individuelles} (personnes physiques) au sens économique.
        \item \textbf{Analyser les facteurs de production :}
              \begin{itemize}
                  \item \cle{Travail} : auto-emploi + apprentis (souvent non déclarés, rémunérés en nature par le logement, la nourriture et la formation). Pénibilité, horaires étendus.
                  \item \cle{Capital} : stock (marchandises), outil de travail (moto, machine à coudre, téléphone), trésorerie (caisse journalière, tontines). \textbf{Accès au financement difficile} (exigences de garanties, taux élevés).
                  \item \cle{Nature} : espace public (trottoir, carrefour) utilisé comme lieu de vente, parfois contre redevances informelles.
              \end{itemize}
        \item \textbf{Évaluer la performance (comptabilité de survie) :}
              \begin{itemize}
                  \item \cle{Recettes journalières} - \cle{dépenses variables} (achats, carburant, redevances, repas) = \cle{marge journalière}.
                  \item \cle{Revenu net disponible} = marge - charges fixes (loyer du domicile, scolarité des enfants, santé, épargne tontine) - amortissement implicite (usure de la moto ou de la machine).
                  \item \textbf{Seuil de rentabilité vital :} le revenu net doit couvrir les besoins essentiels du ménage.
              \end{itemize}
        \item \textbf{Stratégie et RSE « par le bas » :}
              \begin{itemize}
                  \item \cle{Stratégie} : survie et résilience — diversification (ventes multiples), flexibilité horaire, réseaux de solidarité (tontines, groupements).
                  \item \cle{RSE implicite} : insertion des jeunes (apprentissage), service de proximité (quartiers non desservis), recyclage (récupération des métaux, plastiques, friperie). \textbf{Limites :} travail des enfants, insécurité, pollution, évasion fiscale.
              \end{itemize}
        \item \textbf{Problématique de la formalisation :} coûts (impôts, cotisations sociales, registre, comptable) vs avantages (accès au crédit formel, marchés publics, protection sociale, légitimité).
    \end{enumerate}
    \textbf{Exemple type Bac :} « Analysez l'activité de "Maman Ngo", vendeuse de beignets-haricots au carrefour Warda, comme une entreprise. »
\end{methode}

\begin{exoresolu}[Analyse d'une unité informelle : « Maman Ngo » — beignets au carrefour Warda]
    \textbf{Énoncé :} Maman Ngo vend des beignets et des haricots de 6 h à 10 h et de 16 h à 20 h. Données journalières moyennes : vente de 150 beignets $\times$ 100 FCFA + 50 assiettes de haricots $\times$ 500 FCFA. Achats : farine 5 kg (2 500), huile 3 L (3 000), haricots 5 kg (2 500), condiments et bois (2 000). Redevance d'occupation : 500. Transport : 1 000. Elle emploie 2 apprenties (logées et nourries, sans salaire). Livraison de farine en moto-taxi : 2 000/j. Épargne tontine : 2 000/j.
    \begin{enumerate}
        \item Construire le compte de résultat journalier (recettes / charges variables / marge / charges fixes / revenu net).
        \item Calculer le seuil de rentabilité journalier (en FCFA et en nombre d'unités).
        \item Identifier 2 forces, 2 faiblesses, 1 opportunité, 1 menace (SWOT express).
        \item Proposer 1 action de formalisation progressive (coût/bénéfice).
    \end{enumerate}
    \tcblower
    \textbf{Solution détaillée :}
    \begin{enumerate}
        \item \textbf{Compte de résultat journalier :}
              \begin{align*}
                  \textbf{Recettes} &= (150 \times 100) + (50 \times 500) = 15\,000 + 25\,000 = \textbf{40\,000 FCFA} \\
                  \textbf{Charges variables (achats directs)} &= 2\,500 + 3\,000 + 2\,500 + 2\,000 = \textbf{10\,000 FCFA} \\
                  \textbf{Marge brute} &= 40\,000 - 10\,000 = \textbf{30\,000 FCFA} \quad (\text{taux de marge } 75\,\%) \\
                  \textbf{Charges fixes journalières :} \\
                  \quad \text{Redevance d'occupation} &= 500 \\
                  \quad \text{Transport} &= 1\,000 + 2\,000 = 3\,000 \\
                  \quad \text{Amortissement implicite (matériel, réchaud)} &\approx 1\,000 \text{ (estimation)} \\
                  \quad \text{Coût des apprenties (logement, nourriture)} &\approx 2 \times 1\,500 = 3\,000 \\
                  \textbf{Total des charges fixes} &= \textbf{7\,500 FCFA} \\
                  \textbf{Revenu net avant épargne} &= 30\,000 - 7\,500 = \textbf{22\,500 FCFA} \\
                  \textbf{Épargne tontine} &= 2\,000 \\
                  \textbf{Revenu disponible du ménage} &= \textbf{20\,500 FCFA / jour}
              \end{align*}
              \emph{Interprétation :} environ 600 000 FCFA par mois, très supérieur au seuil de pauvreté : l'activité fait vivre correctement le ménage.
        \item \textbf{Seuil de rentabilité journalier :}
              \begin{itemize}
                  \item Charges fixes journalières = 7 500 FCFA ; taux de marge sur coût variable = 75\,\%.
                  \item $Recettes_{PM} = \dfrac{7\,500}{0,75} = \textbf{10\,000 FCFA}$.
                  \item Prix moyen pondéré par unité = $\dfrac{40\,000}{200} = 200$ FCFA, d'où $Q_{PM} = \dfrac{10\,000}{200} = \textbf{50 unités}$.
                  \item \textbf{Marge de sécurité} = 40 000 - 10 000 = 30 000 FCFA, soit 75\,\% des recettes. \textbf{Très confortable :} l'activité est robuste tant que la demande se maintient.
              \end{itemize}
        \item \textbf{SWOT express :}
              \begin{itemize}
                  \item \textbf{Forces :} marge élevée (75\,\%), clientèle fidèle (goût, horaires réguliers), faible investissement initial, maîtrise de la recette.
                  \item \textbf{Faiblesses :} dépendance totale à la personne de Maman Ngo (maladie = arrêt de l'activité), pénibilité (station debout, chaleur, fumée), apprenties non déclarées, trésorerie volatile (pas de fonds de roulement).
                  \item \textbf{Opportunité :} livraison aux bureaux environnants (commandes groupées), vente de beignets prêts à frire, adhésion à un groupement de vendeuses pour mutualiser les achats de farine et d'huile.
                  \item \textbf{Menace :} expulsion par la mairie, hausse brutale des prix de l'huile et de la farine, concurrence de nouveaux vendeurs, maladie ou accident (pas de couverture sociale).
              \end{itemize}
        \item \textbf{Action de formalisation progressive :}
              \textbf{Immatriculation auprès du Centre de Formalités de Création d'Entreprises (CFCE) et adhésion au régime fiscal simplifié des petites entreprises.}
              \begin{itemize}
                  \item \textbf{Coût :} impôt forfaitaire proportionnel au chiffre d'affaires + cotisations sociales si les apprenties sont déclarées. Charge mensuelle estimée : quelques dizaines de milliers de FCFA, supportable au vu de la marge.
                  \item \textbf{Bénéfice :} immatriculation au RCCM (existence légale), numéro d'identifiant unique (NIU), compte bancaire professionnel (accès au microcrédit à taux maîtrisé), couverture sociale (maladie, maternité, vieillesse, accidents du travail) pour elle et ses apprenties, accès aux marchés publics (cantines scolaires, casernes).
              \end{itemize}
    \end{enumerate}
    \textbf{Conclusion :} « Maman Ngo » est une \cle{véritable entreprise} performante (marge de 75\,\%, marge de sécurité de 75\,\%). Son passage au formel est \textbf{rentable} si l'accompagnement administratif (guichet unique, comptabilité simplifiée) réduit les coûts de la démarche.
\end{exoresolu}

\begin{attention}[Les 5 erreurs qui coûtent des points au Bac]
    \begin{enumerate}
        \item \textbf{Confondre « entreprise » (unité économique) et « société » (personne morale / forme juridique).} \emph{Piège :} dire « l'entreprise SARL... ». \textbf{Correct :} « la société SARL \emph{X} est une entreprise... » ou « l'entreprise \emph{X}, constituée sous forme de SARL... ». L'entreprise individuelle n'est PAS une société.
        \item \textbf{Oublier le critère du RISQUE dans la définition.} \emph{Piège :} définir l'entreprise seulement par « la production de biens et services ». \textbf{Correct :} intégrer systématiquement la \cle{prise de risque} sur le capital engagé (perte possible de l'apport). Sans risque $\to$ administration, salarié ou association non lucrative.
        \item \textbf{Mélanger charges fixes et charges variables dans le calcul du seuil de rentabilité.} \emph{Piège :} classer les salaires en charges variables ou l'amortissement en charges variables. \textbf{Correct :} \cle{charges fixes} = loyers, salaires fixes et charges, amortissements, assurances, intérêts ; \cle{charges variables} = matières premières, sous-traitance, commissions sur ventes, énergie du processus. \textbf{Formule :} $CA_{PM} = \dfrac{CF}{\text{Taux de marge sur coût variable}}$.
        \item \textbf{Rédiger une dissertation sans exemples camerounais concrets.} \emph{Piège :} ne citer que Ford, Taylor, Friedman ou Apple. \textbf{Correct :} illustrer \textbf{obligatoirement} par des références locales : grandes entreprises publiques et privées camerounaises, PME locales (« Bois \& Design »), bayam-sellam, GIC, tontines, dispositifs d'appui aux PME, OHADA, CNPS. Cela prouve l'ancrage territorial exigé par le MINESEC.
        \item \textbf{Confondre RSE (Responsabilité Sociétale) et philanthropie.} \emph{Piège :} « l'entreprise fait de la RSE car elle donne de l'argent à un orphelinat ». \textbf{Correct :} la RSE (ISO 26000) est l'\cle{intégration} des préoccupations sociales et environnementales dans la \textbf{stratégie et le cœur de métier} (gouvernance, achats responsables, éco-conception, conditions de travail, loyauté des pratiques). Le don (mécénat) n'en est qu'une petite partie (question « communautés et développement local »), pas la définition.
    \end{enumerate}
\end{attention}

\newpage

\coursec{Exercices}

\courssub{Je vérifie mes connaissances}

\begin{exercice}[QCM : Nature et définition de l'entreprise]
Parmi les propositions suivantes, une seule correspond à la définition juridique et économique de l'entreprise au sens de l'Acte uniforme OHADA relatif au droit des sociétés. Laquelle ?
\begin{enumerate}
    \item Une personne physique qui exerce une activité commerciale en son nom propre.
    \item Une organisation dotée de personnalité morale, réunissant des facteurs de production (capital, travail) pour produire des biens ou services en vue de les mettre sur le marché.
    \item Un service public administratif géré par l'État pour satisfaire un besoin d'intérêt général.
    \item Un groupement d'intérêt économique (GIE) qui n'a pas pour but la réalisation et le partage de bénéfices.
\end{enumerate}
\end{exercice}

\begin{exercice}[Vrai / Faux Justifié : Classification des entreprises]
Indiquez si les affirmations suivantes sont Vraies ou Fausses en justifiant brièvement votre réponse en vous référant au Décret n°2010/088 du 15 avril 2010 fixant les critères de classification des entreprises au Cameroun.
\begin{enumerate}
    \item Une entreprise comptant 15 salariés et réalisant un chiffre d'affaires annuel de 80 millions de FCFA est classée dans la catégorie des Petites Entreprises (PE).
    \item Le critère unique de classification est le nombre d'employés permanents.
    \item Une Grande Entreprise (GE) est nécessairement une Société Anonyme (SA).
    \item Le secteur informel ne fait l'objet d'aucune classification officielle dans cette nomenclature.
\end{enumerate}
\end{exercice}

\begin{exercice}[Définitions précises : Vocabulaire technique]
Donnez une définition rigoureuse (2 à 3 lignes maximum chacune) des notions suivantes en distinguant soigneusement les concepts voisins :
\begin{enumerate}
    \item L'\cle{Établissement} (au sens statistique et juridique) vs l'\cle{Entreprise}.
    \item La \cle{Valeur Ajoutée (VA)} vs le \cle{Chiffre d'Affaires (CA)}.
    \item L'\cle{Excédent Brut d'Exploitation (EBE)} vs le \cle{Résultat Net}.
    \item La \cle{Gouvernance d'entreprise} vs la simple \cle{Gestion courante}.
\end{enumerate}
\end{exercice}

\begin{exercice}[Calculs directs : Indicateurs de performance]
L'entreprise \textbf{« Bois \& Design »} (menuiserie industrielle à Douala) présente les données comptables simplifiées suivantes pour l'exercice 2023 (en millions de FCFA) :
\begin{itemize}
    \item Chiffre d'Affaires HT : 500
    \item Achats consommés (matières premières, marchandises) : 220
    \item Services extérieurs (loyer, électricité, maintenance, transport) : 80
    \item Impôts et taxes (hors IS) : 15
    \item Charges de personnel (brutes) : 120
    \item Dotations aux amortissements : 25
    \item Résultat financier : -10
    \item Résultat exceptionnel : +5
    \item Impôt sur les Sociétés (IS) : 18
\end{itemize}
Calculez :
\begin{enumerate}
    \item La Valeur Ajoutée (VA).
    \item L'Excédent Brut d'Exploitation (EBE).
    \item Le Résultat Net (RN).
    \item Le Taux de Marge Brute (TMB = VA / CA) et le Taux de Rentabilité Nette (RN / CA).
\end{enumerate}
\end{exercice}

\courssub{Je m'entraîne}

\begin{exercice}[Analyse de statuts : SARL vs SA sous droit OHADA]
M. \textsc{Kamga} et trois associés souhaitent créer une structure pour exploiter une plantation de palmiers à huile dans la région du Sud-Ouest. Ils hésitent entre la SARL et la SA.
\begin{enumerate}
    \item Complétez le tableau comparatif suivant en vous basant sur l'Acte uniforme OHADA (nombre d'associés min/max, capital minimum, responsabilité, organes de direction obligatoires, cessibilité des parts/actions).
    \item En tant que conseiller juridique, quelle forme leur recommanderiez-vous si : (a) ils veulent limiter la publicité sur leur patrimoine ; (b) ils envisagent une introduction en bourse dans 5 ans ; (c) l'un des associés est mineur émancipé ? Justifiez.
\end{enumerate}
\end{exercice}

\begin{exercice}[Organisation et structure : Le cas « Bois \& Design »]
L'organigramme de « Bois \& Design » (50 salariés) présente une structure fonctionnelle classique : Direction Générale \textrightarrow{} 4 Directions Fonctionnelles (Production, Commercial, Admin/Finances, RH) \textrightarrow{} Chefs d'atelier / Chefs de service \textrightarrow{} Ouvriers / Employés.
\begin{enumerate}
    \item Identifiez le type de structure (fonctionnelle, divisionnelle, matricielle, en réseau) et définissez le \cle{chaîne hiérarchique} et l'\cle{étendue de contrôle} (span of control) pour le Directeur de Production s'il encadre 3 chefs d'atelier qui gèrent chacun 10 ouvriers.
    \item L'entreprise reçoit une commande urgente nécessitant une collaboration étroite entre Production, Commercial et Achats. Quel dysfonctionnement structurel (\"effet silo\") risque d'apparaître ? Proposez deux modes de coordination transversale (comités de projet, intégrateur, etc.) pour y remédier.
    \item Calculez le \cle{taux d'encadrement} (cadres / total effectif) si l'entreprise compte 5 cadres dirigeants, 8 cadres moyens (chefs) et 37 agents d'exécution. Interprétez ce ratio.
\end{enumerate}
\end{exercice}

\begin{exercice}[RSE et Conflit social : Mise en situation]
Dans l'usine « Bois \& Design », un conflit oppose la direction aux délégués du personnel (SYNTRA-PME). Les revendications : augmentation des salaires de base de 15\% (inflation importée), amélioration des EPI (Équipements de Protection Individuelle) en atelier sciage, respect des repos hebdomadaires (travail dominical récurrent non majoré).
\begin{enumerate}
    \item Qualifiez ce conflit (intérêt, droit, reconnaissance) et citez les étapes de la procédure de conciliation obligatoire au Cameroun (Code du Travail, Loi n°92/007).
    \item En mobilisant la norme ISO 26000 (7 questions centrales), identifiez les 3 piliers de la RSE directement concernés par ce conflit.
    \item Rédigez une proposition de \cle{Protocole d'Accord} (3 articles clés) intégrant une clause de \cle{paix sociale} et un engagement RSE mesurable (ex: audit sécurité annuel, comité paritaire HSE).
\end{enumerate}
\end{exercice}

\begin{exercice}[Diagnostic financier simplifié : Évolution sur 2 ans]
Données comparées « Bois \& Design » (en millions FCFA) :
\begin{center}
\begin{tabular}{|l|c|c|}\hline
\textbf{Indicateur} & \textbf{2022} & \textbf{2023} \\ \hline
Chiffre d'Affaires & 450 & 500 \\
Valeur Ajoutée & 180 & 195 \\
EBE & 45 & 50 \\
Résultat Net & 12 & 15 \\
Effectif moyen & 45 & 50 \\
Masse salariale brute & 100 & 120 \\ \hline
\end{tabular}
\end{center}
\begin{enumerate}
    \item Calculez les taux de variation (en \%) du CA, de la VA, de l'EBE et du RN entre 2022 et 2023.
    \item Calculez la \cle{Productivité Apparente du Travail} (VA / Effectif) et la \cle{Rémunération Moyenne Annuelle} (Masse salariale / Effectif) pour les deux années. Analysez l'évolution du partage de la VA (Taux de redistribution = Masse salariale / VA).
    \item L'entreprise est-elle en \cle{croissance extensive} ou \cle{intensive} ? Justifiez par les chiffres.
\end{enumerate}
\end{exercice}

\courssub{Je me prépare au Bac}

\begin{exercice}[Sujet Type Dissertation — Épreuve de Philosophie (Coeff. 4, 4h)]
\textbf{Sujet :} \og L'entreprise a-t-elle pour seule finalité le profit ? \fg
\vspace{0.2cm}
\textbf{Consignes :} Vous traiterez ce sujet en mobilisant les connaissances du cours (notions de finalité lucrative, responsabilité sociale, parties prenantes / \emph{stakeholders}, droit OHADA, exemples camerounais concrets : SAFACAM, CDC, PME locales, secteur informel). Votre devoir comportera une introduction (accroche, définition des termes, problématique, annonce du plan), un développement dialectique en trois parties (Thèse : la contrainte économique et juridique du profit ; Antithèse : la responsabilité sociétale et les finalités non lucratives ; Synthèse : l'entreprise comme communauté de travail et acteur du développement durable), et une conclusion (réponse nuancée, ouverture).
\end{exercice}

\begin{exercice}[Sujet Type Explication de Texte — Épreuve de Philosophie (Coeff. 4, 4h)]
\textbf{Texte :} Extrait du rapport \emph{« Le « Missing Middle » au Cameroun : le chaînon manquant de l'émergence »} (Inspiré CNUCED / Banque Mondiale, 2023).
\og L'économie camerounaise se caractérise par une structure bimodale : d'un côté, quelques grandes entreprises formelles, souvent filiales de groupes étrangers ou publiques, capitalistiques, exportatrices (pétrole, bois, agro-industrie, banques) ; de l'autre, une myriade d'unités de production informelles, de survie, à très faible productivité. Entre les deux, le « Missing Middle » — les PME formelles de 20 à 200 employés — est atrophié. Elles souffrent d'un coût du crédit prohibitif (taux > 12\%), d'une fiscalité perçue comme prédatrice, d'une électricité instable et d'un cadre réglementaire complexe. Pourtant, ce sont ces PME qui, dans les pays émergents, portent l'innovation incrémentale, l'emploi décent des jeunes et la transformation locale des matières premières. Redynamiser ce tissu, c'est passer de la rente à la création de valeur partagée. \fg
\vspace{0.2cm}
\textbf{Questions :}
\begin{enumerate}
    \item Expliquez les concepts suivants dans le contexte du texte : \cle{structure bimodale}, \cle{Missing Middle}, \cle{productivité}, \cle{création de valeur partagée}.
    \item Dégagez la thèse principale de l'auteur et montrez comment il l'argumente (causes du blocage, conséquences, rôle attendu des PME).
    \item Discutez la pertinence de cette analyse pour le Cameroun d'aujourd'hui : la loi de finances 2024, le programme PADI (Programme d'Appui au Développement de l'Entrepreneuriat) et la ZLECAf répondent-ils aux obstacles identifiés ? Apportez des contre-arguments ou des nuances (ex: résilience de l'informel, rôle des GE dans les infrastructures).
\end{enumerate}
\end{exercice}

\begin{exercice}[Cas Pratique Intégré — Épreuve Écrite de Spécialité / Contrôle Continu (2h)]
\textbf{Dossier : « Bois \& Design » — Diagnostic Stratégique et Social}
\textbf{Document 1 : Extrait Statuts.} Société à Responsabilité Limitée (SARL) au capital de 50 000 000 FCFA. Gérant unique : M. \textsc{Fouda}. Siège : Douala-Bonabéri. Objet : Transformation du bois, fabrication meubles.
\textbf{Document 2 : Organigramme simplifié.} (Voir exercice « Je m'entraîne » n°2).
\textbf{Document 3 : Comptes de Résultat simplifiés 2022-2023.} (Voir exercice « Je m'entraîne » n°4).
\textbf{Document 4 : Article de presse « La Voix du Koa », mars 2024.} \og Grève à Bois \& Design : les ouvriers dénoncent l'absence de médecine du travail depuis 18 mois et le non-paiement des heures supplémentaires du dimanche. La direction invoque la trésorerie tendue suite à l'augmentation du prix du gaz industriel et du carburant pour les groupes électrogènes. \fg

\textbf{Travail demandé :}
\begin{enumerate}
    \item \textbf{Qualification juridique :} Précisez la forme sociale, la responsabilité des associés, l'organe de direction et le régime fiscal de l'IS applicable. Citez l'article de l'Acte Uniforme OHADA définissant la SARL.
    \item \textbf{Analyse Financière :} À partir du Document 3, calculez pour 2023 : la VA, l'EBE, le Taux de Marge Brute, la Productivité du Travail (VA/Effectif). Comparez à 2022. Diagnostiquez la santé financière (rentabilité, structure des charges, autofinancement).
    \item \textbf{Analyse Organisationnelle :} À partir du Document 2, identifiez le mode de structure, l'étendue de contrôle du Directeur Production, le taux d'encadrement. Quels sont les risques de dysfonctionnement (communication, réactivité) pour une PME de cette taille ?
    \item \textbf{Résolution du conflit \& RSE :} En vous appuyant sur le Document 4, le Code du Travail camerounais et les principes de la RSE (ISO 26000), rédigez une \textbf{Note de synthèse à l'attention du Gérant} (250 mots max) proposant : (a) la qualification juridique du conflit ; (b) 3 mesures correctives immédiates (droit du travail) ; (c) 2 engagements RSE structurels (santé/sécurité, dialogue social) pour prévenir la récidive.
\end{enumerate}
\end{exercice}

\newpage

\coursec{Corrigés des exercices}

\begin{corrige}[Exercice 1 — QCM : Nature et définition de l'entreprise]
\textbf{Bonne réponse : Proposition 2.}

\textbf{Justification :}
\begin{itemize}
    \item \textbf{Proposition 1 (Fausse) :} Décrit une \emph{entreprise individuelle} (personne physique), qui n'a pas de personnalité morale distincte de son exploitant (patrimoine unique, responsabilité illimitée), sauf option pour l'EIRL (peu courante en droit OHADA classique).
    \item \textbf{Proposition 2 (Vraie) :} Reprend la définition consensuelle (économique et juridique OHADA, Acte Uniforme relatif au droit des sociétés commerciales et du GIE) : \textbf{personnalité morale} (distincte des associés), \textbf{réunion de facteurs de production} (capital, travail), \textbf{but lucratif} (production de biens et services pour le marché), \emph{affectio societatis} (volonté de collaborer à une entreprise commune).
    \item \textbf{Proposition 3 (Fausse) :} Décrit un \emph{service public administratif} (SPA) : pas de personnalité morale propre (sauf EPA), pas de but lucratif, financement public, régime de droit public.
    \item \textbf{Proposition 4 (Fausse) :} Le \emph{GIE} (Groupement d'Intérêt Économique) a pour objet de \emph{faciliter} ou \emph{développer} l'activité économique de ses membres, non de réaliser et de partager des bénéfices pour lui-même.
\end{itemize}
\end{corrige}

\begin{corrige}[Exercice 2 — Vrai / Faux Justifié : Classification des entreprises]
\textbf{Rappel du Décret n°2010/088 (Art. 2) :} la classification repose sur \textbf{deux critères combinés} : l'\textbf{effectif salarié permanent} ET le \textbf{chiffre d'affaires annuel} (ou le total du bilan). L'entreprise est classée dans la catégorie la plus élevée atteinte par l'un des deux critères.

\begin{enumerate}
    \item \textbf{VRAI.} Effectif = 15 (tranche 10--50 $\rightarrow$ PE) ; CA = 80 M FCFA (tranche 50 M -- 1 Md $\rightarrow$ PE). Les deux critères concordent vers la catégorie \textbf{Petite Entreprise (PE)}.
    \item \textbf{FAUX.} Le décret retient \textbf{deux critères} (effectif \textbf{et} chiffre d'affaires / total du bilan), pas un critère unique. C'est une classification \emph{bi-critère}.
    \item \textbf{FAUX.} La taille (GE, ME, PE, micro-entreprise) est une classification \textbf{économique et statistique}. La forme juridique (SA, SARL, SNC, etc.) est un choix \textbf{juridique}. Une SARL peut être une GE si elle dépasse les seuils (effectif $>$ 200 et/ou CA $>$ 5 Mds FCFA).
    \item \textbf{VRAI (avec nuance).} Le décret s'applique aux entreprises au sens formel (immatriculées, tenues à la comptabilité). Le \textbf{secteur informel} (unités non immatriculées, non assujetties à l'IS/TVA) n'entre pas dans cette nomenclature officielle ; il fait l'objet d'enquêtes spécifiques (ex. : EESI de l'INS) mais pas de cette classification réglementaire.
\end{enumerate}
\end{corrige}

\begin{corrige}[Exercice 3 — Définitions précises : Vocabulaire technique]
\begin{enumerate}
    \item \textbf{Établissement vs Entreprise :}
    \begin{itemize}
        \item \cle{Établissement} : unité de production \textbf{géographiquement individualisée} (une adresse, un lieu), dépendante juridiquement de l'entreprise, où s'exercent une ou plusieurs activités (ex. : l'usine de Bonabéri, le magasin d'Akwa).
        \item \cle{Entreprise} : unité \textbf{juridique} (personnalité morale, identifiant unique), \textbf{économique} (autonomie de décision, gestion unifiée) et \textbf{sociale} (chef d'entreprise unique), pouvant regrouper plusieurs établissements.
    \end{itemize}
    \item \textbf{Valeur Ajoutée (VA) vs Chiffre d'Affaires (CA) :}
    \begin{itemize}
        \item \cle{CA} : somme des ventes de biens et services facturées aux clients (Prix $\times$ Quantités). Mesure l'\emph{activité commerciale} (le volume des ventes).
        \item \cle{VA} : richesse \textbf{créée} par l'entreprise = Production $-$ Consommations intermédiaires (achats de matières premières + services extérieurs). Mesure la \emph{contribution économique} à la richesse nationale (PIB).
        \[ \text{VA} = \text{CA} - \text{Achats consommés} - \text{Services extérieurs} \]
    \end{itemize}
    \item \textbf{Excédent Brut d'Exploitation (EBE) vs Résultat Net (RN) :}
    \begin{itemize}
        \item \cle{EBE} : excédent dégagé par le \textbf{cycle d'exploitation seul}, \textbf{avant} amortissements, provisions, charges financières et éléments exceptionnels. Indicateur de \emph{performance industrielle} pure (capacité à générer de la trésorerie par l'activité courante).
        \[ \text{EBE} = \text{VA} - \text{Charges de personnel} - \text{Impôts et taxes (hors IS)} \]
        \item \cle{RN} : résultat \textbf{final} de l'exercice = Résultat d'exploitation + Résultat financier + Résultat exceptionnel $-$ Impôt sur les Sociétés (IS). Mesure la \emph{performance globale} pour les propriétaires (associés ou actionnaires).
    \end{itemize}
    \item \textbf{Gouvernance d'entreprise vs Gestion courante :}
    \begin{itemize}
        \item \cle{Gestion courante} : exécution opérationnelle quotidienne (achats, production, ventes, paie, trésorerie) par la \textbf{Direction Générale} et les cadres. Horizon de court terme, recherche d'efficacité.
        \item \cle{Gouvernance} : ensemble des \textbf{règles, organes et processus} (Assemblée Générale, Conseil d'Administration, Commissaires aux comptes) qui \textbf{orientent, contrôlent et rendent compte} aux parties prenantes (actionnaires, État, salariés). Horizon stratégique, légitimité, éthique, conformité (OHADA, ISO 26000).
    \end{itemize}
\end{enumerate}
\end{corrige}

\begin{corrige}[Exercice 4 — Calculs directs : Indicateurs de performance « Bois \& Design » 2023]
\textbf{Données (en millions de FCFA) :} CA = 500 ; Achats = 220 ; Services extérieurs = 80 ; Impôts et taxes = 15 ; Charges de personnel = 120 ; Amortissements = 25 ; Résultat financier = $-10$ ; Résultat exceptionnel = $+5$ ; IS = 18.

\begin{enumerate}
    \item \textbf{Valeur Ajoutée (VA)}
    \[ \text{VA} = \text{CA} - \text{Achats} - \text{Services extérieurs} = 500 - 220 - 80 = \textbf{200 M FCFA} \]

    \item \textbf{Excédent Brut d'Exploitation (EBE)}
    \[ \text{EBE} = \text{VA} - \text{Charges de personnel} - \text{Impôts et taxes} = 200 - 120 - 15 = \textbf{65 M FCFA} \]

    \item \textbf{Résultat Net (RN)}
    \begin{align*}
    \text{Résultat d'Exploitation (RE)} &= \text{EBE} - \text{Amortissements} = 65 - 25 = 40 \\
    \text{Résultat Courant Avant Impôt (RCAI)} &= \text{RE} + \text{Résultat financier} = 40 + (-10) = 30 \\
    \text{Résultat Avant Impôt (RAI)} &= \text{RCAI} + \text{Résultat exceptionnel} = 30 + 5 = 35 \\
    \text{RN} &= \text{RAI} - \text{IS} = 35 - 18 = \textbf{17 M FCFA}
    \end{align*}

    \item \textbf{Taux de Marge sur Valeur Ajoutée et Taux de Rentabilité Nette}
    \[ \text{Taux de marge} = \frac{\text{VA}}{\text{CA}} = \frac{200}{500} = 0,40 = \textbf{40\,\%} \]
    \[ \text{Rentabilité nette} = \frac{\text{RN}}{\text{CA}} = \frac{17}{500} = 0,034 = \textbf{3,4\,\%} \]
\end{enumerate}
\textbf{Interprétation :} l'entreprise crée 40\,\% de valeur sur son chiffre d'affaires (bon niveau pour l'industrie du bois). L'EBE de 65 M assure l'autofinancement (amortissements de 25 M + remboursement des dettes). La rentabilité nette est faible (3,4\,\%), en raison des charges financières (10 M) et de l'IS (18 M).
\end{corrige}

\begin{corrige}[Exercice 5 — Analyse de statuts : SARL vs SA sous droit OHADA]
\begin{enumerate}
    \item \textbf{Tableau comparatif (Acte Uniforme OHADA révisé)}
    \begin{center}
    \begin{tabularx}{\linewidth}{|l|X|X|}\hline
    \rowcolor{popblueL}
    \textbf{Critère} & \textbf{SARL} & \textbf{SA} \\ \hline
    \textbf{Nombre d'associés / actionnaires} & Min 1 (SARL unipersonnelle) -- Max 100 & Min 1 (SA unipersonnelle possible) -- Max illimité \\ \hline
    \textbf{Capital minimum} & \textbf{100 000 FCFA} (libération d'au moins 1/4 à la constitution) & \textbf{10 000 000 FCFA} (libération d'au moins 1/4 des apports en numéraire) \\ \hline
    \textbf{Responsabilité} & Limitée aux apports & Limitée aux apports \\ \hline
    \textbf{Organes de direction} & \textbf{Gérant(s)} (personne physique) + Assemblée des associés. Pas de Conseil d'Administration obligatoire. & \textbf{Conseil d'Administration} (3 à 12 membres) + \textbf{Directeur Général}, ou Administrateur Général (SA à 2 actionnaires ou moins). \\ \hline
    \textbf{Cession des parts / actions} & \textbf{Agrément obligatoire} (sauf cession entre conjoints, ascendants, descendants). & \textbf{Libre} en principe (négociabilité des actions). Clauses d'agrément possibles mais encadrées. \\ \hline
    \textbf{Commissaire aux Comptes (CAC)} & Obligatoire si dépassement de 2 des 3 seuils légaux (total bilan, CA, effectif) ou sur demande d'un associé. & \textbf{Obligatoire} (au moins 1 titulaire + 1 suppléant). \\ \hline
    \textbf{Publicité / Transparence} & Moins contraignante (dépôt des comptes au greffe). & Plus contraignante (publicité des comptes annuels, registre du commerce). \\ \hline
    \end{tabularx}
    \end{center}

    \item \textbf{Recommandation du conseiller juridique :}
    \begin{enumerate}
        \item[(a)] \textbf{Limiter la publicité sur le patrimoine :} \textbf{SARL.} Publicité comptable allégée (dépôt au greffe), contrairement à la SA. Adaptée à la discrétion d'une entreprise familiale.
        \item[(b)] \textbf{Introduction en bourse dans 5 ans :} \textbf{SA.} \textbf{Seule la SA peut faire appel public à l'épargne} et être cotée (Bourse des Valeurs Mobilières de l'Afrique Centrale, Douala). La SARL ne peut émettre ni actions ni obligations auprès du public. La transformation SARL $\rightarrow$ SA est possible mais coûteuse (augmentation de capital, CAC, assemblée extraordinaire).
        \item[(c)] \textbf{Associé mineur émancipé :} \textbf{SARL (ou SA).} Le mineur émancipé dispose d'une capacité étendue et peut être associé dans les deux formes. \emph{Nuance :} il ne peut être ni Administrateur d'une SA ni Gérant d'une SARL (fonctions réservées aux personnes physiques majeures). S'il ne gère pas, les deux formes conviennent ; la SARL reste plus simple.
    \end{enumerate}
\end{enumerate}
\end{corrige}

\begin{corrige}[Exercice 6 — Organisation et structure : Le cas « Bois \& Design »]
\begin{enumerate}
    \item \textbf{Type de structure :} \textbf{structure fonctionnelle (hiérarchique pyramidale)} : regroupement par spécialités ou métiers (Production, Commercial, Administratif et Finances, Ressources Humaines).

    \textbf{Chaîne hiérarchique (Directeur de Production) :} DG $\rightarrow$ Directeur de Production $\rightarrow$ 3 Chefs d'atelier $\rightarrow$ Ouvriers (4 niveaux).

    \textbf{Étendue de contrôle (\emph{span of control}) du Directeur de Production :} \textbf{3} (il encadre directement 3 Chefs d'atelier). \emph{Note :} la définition managériale standard ne compte que les \emph{subordonnés directs}, et non les 30 ouvriers encadrés indirectement.

    \item \textbf{Dysfonctionnement « effet silo » :} chaque direction optimise son objectif local (la Production veut des séries longues, le Commercial veut du sur-mesure urgent, les Achats veulent de gros volumes fournisseurs) $\rightarrow$ \textbf{manque de communication transversale, lenteur, conflits d'objectifs, non-respect des délais clients.}

    \textbf{Deux modes de coordination transversale :}
    \begin{enumerate}
        \item \textbf{Comité de projet / task force « Commande urgente » :} réunion quotidienne pilotée par un \emph{chef de projet} (intégrateur) réunissant les directions Production, Commercial et Achats, avec pouvoir de décision rapide sur l'ordonnancement et les approvisionnements.
        \item \textbf{Intégrateur permanent (responsable méthodes / ordonnancement) :} poste transversal rattaché à la DG, chargé de la \emph{planification fine} et de l'interface Production / Achats / Commercial (suivi d'indicateurs : taux de livraison à l'heure, respect des délais).
    \end{enumerate}

    \item \textbf{Taux d'encadrement :}
    \[ \text{Cadres} = 5 \text{ (dirigeants)} + 8 \text{ (cadres moyens)} = 13 \]
    \[ \text{Taux d'encadrement} = \frac{13}{50} = 0,26 = \textbf{26\,\%} \]

    \textbf{Interprétation :} ratio élevé (la norme industrielle se situe souvent entre 10\,\% et 15\,\%). Signe d'une \textbf{structure lourde en encadrement} pour 50 personnes. Avantages : proximité, contrôle qualité, formation. Risques : \textbf{coût de structure élevé} (masse salariale des cadres), lenteur décisionnelle (trop de niveaux), risque de bureaucratie. À surveiller si le chiffre d'affaires stagne.
\end{enumerate}
\end{corrige}

\begin{corrige}[Exercice 7 — RSE et Conflit social : Mise en situation]
\begin{enumerate}
    \item \textbf{Qualification du conflit :} \textbf{conflit mixte.}
    \begin{itemize}
        \item \textbf{Conflit d'intérêt (économique) :} revendication d'augmentation des salaires (+15\,\%), majoration des heures supplémentaires dominicales (création de droits nouveaux).
        \item \textbf{Conflit de droit (juridique) :} non-respect du repos hebdomadaire, non-paiement des majorations d'heures supplémentaires, absence de médecine du travail, équipements de protection individuelle (EPI) défaillants : violations de la loi et du contrat de travail.
    \end{itemize}
    \textbf{Procédure de règlement (Loi n°92/007 portant Code du Travail) :}
    \begin{enumerate}
        \item Dépôt d'un \textbf{préavis de grève} ;
        \item Saisine de l'\textbf{Inspecteur du Travail} (conciliateur) ;
        \item \textbf{Réunions de conciliation} ;
        \item Procès-verbal de conciliation (accord total ou partiel) ou PV de non-conciliation ;
        \item En cas d'échec : \textbf{médiation} (médiateur choisi par les parties ou désigné) ;
        \item En cas d'échec : \textbf{arbitrage volontaire} (la sentence s'impose aux parties) ;
        \item Exercice licite du \textbf{droit de grève} (décision prise à la majorité absolue des salariés, au scrutin secret).
    \end{enumerate}

    \item \textbf{Piliers RSE (ISO 26000, questions centrales) concernés :}
    \begin{enumerate}
        \item \textbf{Conditions de travail et relations de travail :} santé et sécurité (EPI, médecine du travail), temps de travail, rémunération, dialogue social, formation.
        \item \textbf{Droits de l'homme :} droit à des conditions de travail justes et favorables, non-discrimination, liberté syndicale (délégués du personnel).
        \item \textbf{Pratiques loyales :} relations équitables avec les parties prenantes (salariés), lutte contre la corruption, responsabilité vis-à-vis des sous-traitants.
    \end{enumerate}

    \item \textbf{Projet de Protocole d'Accord (3 articles clés) :}

    \textbf{Article 1 -- Mesures immédiates de mise en conformité (Droit du Travail) :}
    \og La Direction s'engage à : (a) régulariser le paiement des heures supplémentaires dominicales avec majoration légale sous 30 jours ; (b) garantir le repos hebdomadaire de 24 heures consécutives, sauf dérogation de l'Inspecteur du Travail ; (c) fournir des EPI conformes aux ateliers de sciage sous 15 jours ; (d) mandater le médecin du travail (visite annuelle obligatoire) sous 60 jours. \fg

    \textbf{Article 2 -- Clause de paix sociale et dialogue structuré :}
    \og Les parties s'interdisent toute action de grève ou de lock-out pendant 24 mois. Création d'un \textbf{Comité Paritaire de Suivi (CPS)} (3 délégués du personnel, 3 représentants de la Direction, 1 Inspecteur du Travail observateur) réuni mensuellement. Ordre du jour : conditions de travail, sécurité, grille salariale, formation. \fg

    \textbf{Article 3 -- Engagements RSE structurels (mesurables) :}
    \og (a) \textbf{Audit hygiène-sécurité-environnement annuel externe}, avec publication du rapport aux salariés ; (b) \textbf{Accord d'entreprise « Santé-Sécurité-Qualité de Vie au Travail »} négocié avant fin 2024 (objectifs : baisse du taux de fréquence des accidents, 100\,\% d'EPI conformes, 1 formation sécurité par an et par salarié) ; (c) \textbf{Révision de la grille salariale} indexée sur l'inflation (INS) + partage des gains de productivité (intéressement). \fg
\end{enumerate}
\end{corrige}

\begin{corrige}[Exercice 8 — Diagnostic financier simplifié : Évolution sur 2 ans]
\begin{enumerate}
    \item \textbf{Taux de variation 2022 $\rightarrow$ 2023 (en \%)}
    \begin{align*}
    \Delta\text{CA} &= \frac{500-450}{450} = \textbf{+11,1\,\%} \\
    \Delta\text{VA} &= \frac{195-180}{180} = \textbf{+8,3\,\%} \\
    \Delta\text{EBE} &= \frac{50-45}{45} = \textbf{+11,1\,\%} \\
    \Delta\text{RN} &= \frac{15-12}{12} = \textbf{+25,0\,\%}
    \end{align*}

    \item \textbf{Productivité apparente du travail (VA / Effectif) et rémunération moyenne (Masse salariale / Effectif)}
    \begin{center}
    \begin{tabular}{|l|c|c|c|}\hline
    \rowcolor{popblueL}
    \textbf{Indicateur} & \textbf{2022} & \textbf{2023} & \textbf{Évolution} \\ \hline
    Productivité (M FCFA/pers.) & $180/45 = 4,00$ & $195/50 = 3,90$ & \textbf{$-2,5$\,\%} (baisse) \\ \hline
    Rémunération moyenne (M FCFA/pers.) & $100/45 = 2,22$ & $120/50 = 2,40$ & \textbf{+8,1\,\%} (hausse) \\ \hline
    Part des salaires (Masse/VA) & $100/180 = 55,6\,\%$ & $120/195 = 61,5\,\%$ & \textbf{+5,9 points} \\ \hline
    \end{tabular}
    \end{center}
    \textbf{Analyse du partage de la VA :} la part salariale dans la VA augmente fortement (55,6\,\% $\rightarrow$ 61,5\,\%). L'effort salarial (rémunération moyenne +8,1\,\% et effectif +11,1\,\%) absorbe l'essentiel de la croissance de la VA (+8,3\,\%). La marge dégagée pour l'autofinancement (EBE/VA) reste stable mais faible : $45/180 = 25,0\,\%$ puis $50/195 = 25,6\,\%$. \textbf{Risque :} fragilisation de la capacité d'autofinancement (CAF) pour investir.

    \item \textbf{Type de croissance :} \textbf{croissance EXTENSIVE.} \\
    \textbf{Preuves :} l'effectif progresse (+11,1\,\%) pendant que la productivité recule ($-2,5$\,\%). La croissance du CA (+11,1\,\%) et de la VA (+8,3\,\%) est tirée quasi exclusivement par l'\textbf{augmentation du volume de travail (embauches)}, et non par le progrès technique ou l'efficacité. L'entreprise « grossit » sans gagner en productivité. Nécessité d'investir (machines, formation, organisation) pour basculer vers une croissance intensive.
\end{enumerate}
\end{corrige}

\begin{corrige}[Exercice 9 — Sujet Type Dissertation : « L'entreprise a-t-elle pour seule finalité le profit ? »]
\textbf{BARÈME INDICATIF (20 pts) :} Introduction (3) + Thèse (4) + Antithèse (4) + Synthèse (4) + Conclusion/Ouverture (3) + Qualité de la langue et de la structure (2).

\textbf{Points de vigilance (attendus au Probatoire / Bac technique) :}
\begin{itemize}
    \item \textbf{Définitions :} profit (résultat net, rémunération du capital), finalité (but assigné), entreprise (organisation durable réunissant des parties prenantes).
    \item \textbf{Auteurs et concepts clés :} \textbf{A. Smith} (main invisible, intérêt propre $\rightarrow$ intérêt général), \textbf{M. Friedman} (1970 : « la responsabilité sociale de l'entreprise est d'augmenter ses profits »), \textbf{R. E. Freeman} (théorie des parties prenantes, \emph{stakeholders}), \textbf{Porter \& Kramer} (création de valeur partagée, \emph{shared value}), \textbf{droit OHADA} (intérêt social), \textbf{ISO 26000 / RSE}, \textbf{entreprise à mission}.
    \item \textbf{Exemples camerounais obligatoires :} SAFACAM / HEVECAM (certification, écoles, dispensaires), la CDC (poids social d'une entreprise publique), les PME locales (Bois \& Design, conflictualité), le secteur informel (survie vs profit formel).
\end{itemize}

\textbf{Plan détaillé (corrigé type) :}

\textbf{INTRODUCTION}
\begin{itemize}
    \item \textbf{Accroche :} « Une entreprise qui ne fait que du profit est une mauvaise entreprise » (H. Ford).
    \item \textbf{Définitions :} l'entreprise est une organisation pérenne combinant capital et travail pour produire une valeur marchande. Le profit est l'excédent des revenus sur les charges, condition de survie et de réinvestissement. La finalité est la raison d'être ultime.
    \item \textbf{Problématique :} le profit est-il une \textbf{condition nécessaire} (contrainte vitale) ou une \textbf{finalité suffisante} (but unique) ? La tension entre \emph{rentabilité actionnariale} et \emph{responsabilité sociétale} redéfinit-elle l'objet social de l'entreprise ?
    \item \textbf{Annonce du plan :} 1) la contrainte juridique et économique du profit (thèse) ; 2) l'émergence de finalités non lucratives : RSE et parties prenantes (antithèse) ; 3) l'entreprise comme communauté de travail et acteur du développement durable (synthèse).
\end{itemize}

\textbf{DÉVELOPPEMENT}

\textbf{I. THÈSE : Le profit, contrainte vitale et obligation juridique.}
\begin{itemize}
    \item \textbf{A. Impératif économique :} survie dans la concurrence (Schumpeter : l'innovation procure un profit temporaire). L'autofinancement (CAF) permet d'investir, d'innover, de rembourser les dettes. Sans profit $\rightarrow$ faillite $\rightarrow$ destruction d'emplois et de valeur. Ex. : les PME camerounaises (Bois \& Design, RN de 3,4\,\%) restent fragiles.
    \item \textbf{B. Impératif juridique (droit OHADA) :} l'\cle{intérêt social} est l'intérêt de la société en tant que personne morale. Les dirigeants ont un \textbf{devoir de loyauté} : gérer pour la prospérité de l'entité (donc pour le profit). \textbf{Friedman (1970)} : la responsabilité sociale de l'entreprise est d'augmenter ses profits, dans le respect des règles du jeu. Les actionnaires sont les propriétaires résiduels : le profit rémunère leur prise de risque.
    \item \textbf{C. Limite :} le profit \emph{maximum} de court terme n'est pas le profit \emph{optimal} de long terme. La myopie managériale (coupes dans la R\&D, la maintenance, la sécurité) détruit de la valeur (accidents chez Bois \& Design).
\end{itemize}

\textbf{II. ANTITHÈSE : L'entreprise, personne morale responsable : des finalités plurielles.}
\begin{itemize}
    \item \textbf{A. Théorie des parties prenantes (Freeman, 1984) :} l'entreprise doit créer de la valeur pour \textbf{tous} : actionnaires (profit), salariés (salaire, sécurité, sens), clients (qualité), fournisseurs (équité), État (impôts, respect de la loi), collectivités et environnement (RSE). Ex. : SAFACAM (certification, écoles, dispensaires) $\rightarrow$ « licence sociale d'opérer » (acceptabilité locale).
    \item \textbf{B. Responsabilité Sociétale des Entreprises (RSE / ISO 26000) :} 7 questions centrales (gouvernance, droits humains, conditions de travail, environnement, pratiques loyales, consommateurs, communautés). Au Cameroun, les programmes d'appui (PADI) et la ZLECAf exigent des standards RSE pour l'exportation.
    \item \textbf{C. Nouvelles formes juridiques :} \textbf{entreprise à mission} (raison d'être statutaire, objectifs sociaux et environnementaux contraignants) ; \textbf{coopératives} (gouvernance démocratique, profit partagé ou réinvesti) ; \textbf{économie sociale et solidaire} au Cameroun (GIC, coopératives cacao-café) : finalité de \emph{service aux membres} avant le profit du capital.
\end{itemize}

\textbf{III. SYNTHÈSE : L'entreprise comme « communauté de travail » et acteur du développement durable.}
\begin{itemize}
    \item \textbf{A. Le profit comme \textbf{moyen} et \textbf{indicateur}, non comme fin unique :} « le profit est à l'entreprise ce que la respiration est à l'homme : vital, mais on ne vit pas pour respirer ». La \textbf{création de valeur partagée} (Porter \& Kramer) articule compétitivité et progrès social. Ex. : investir dans les EPI et la formation chez Bois \& Design $\rightarrow$ baisse des accidents $\rightarrow$ productivité en hausse $\rightarrow$ profit durable.
    \item \textbf{B. Gouvernance élargie :} passage de la \emph{valeur actionnariale} (\emph{shareholder value}) à la \emph{valeur partagée entre parties prenantes} (\emph{stakeholder value}). Le Conseil d'Administration intègre des compétences RSE et la représentation des salariés.
    \item \textbf{C. Ancrage camerounais :} les PME (le « chaînon manquant ») sont un levier d'emploi décent et de transformation locale (bois, palme, cacao). Nécessité d'une \textbf{profitabilité responsable} : crédits adaptés, marchés publics (notation RSE), standards ZLECAf. Pour le secteur informel, le passage au formel donne accès aux droits et à la protection, en contrepartie de devoirs (impôts, normes).
\end{itemize}

\textbf{CONCLUSION}
\begin{itemize}
    \item \textbf{Réponse nuancée :} non, le profit n'est pas la \emph{seule} finalité, mais il est la \emph{condition nécessaire} de toute finalité. L'entreprise moderne est une \textbf{coalition de parties prenantes} dont la pérennité exige l'alignement du profit (économique), du sens (social) et de la responsabilité (environnementale).
    \item \textbf{Ouverture :} vers une mesure élargie de la performance : la \textbf{triple performance} (financière, sociale, environnementale -- \emph{Triple Bottom Line}, Elkington). Défis camerounais : éducation financière des dirigeants, cadre fiscal incitatif (bonus RSE), sécurité juridique.
\end{itemize}
\end{corrige}

\begin{corrige}[Exercice 10 — Sujet Type Explication de Texte : « Missing Middle » au Cameroun]
\textbf{BARÈME INDICATIF (20 pts) :} Q1 Concepts (5) + Q2 Thèse et argumentation (7) + Q3 Discussion (8).

\textbf{1. Explication des concepts :}
\begin{itemize}
    \item \textbf{Structure bimodale :} dualisme structurel de l'économie : un \textbf{pôle « moderne / formel »} (peu de grandes entreprises capitalistiques, exportatrices, productives, souvent étrangères ou publiques : pétrole, bois, banques, agro-industrie) face à un \textbf{pôle « informel / traditionnel »} (myriade de micro-unités de survie, faible productivité, hors du cadre légal et fiscal). \textbf{Absence d'une classe moyenne entrepreneuriale.}
    \item \textbf{Missing Middle (chaînon manquant) :} déficit structurel des \textbf{PME formelles (20 à 200 employés)}. Dans les pays émergents (Asie, Chili), ce segment porte la croissance inclusive, l'innovation incrémentale, l'emploi décent des jeunes et la transformation locale. Au Cameroun, il est « atrophié ».
    \item \textbf{Productivité :} rapport entre la \textbf{production (VA)} et un \textbf{facteur de production (travail ou capital)}. Faible dans l'informel (travail précaire, peu de capital, peu de technologie). Clé de la croissance intensive (cf. Exo 8 : baisse de productivité = croissance extensive). Le texte relie hausse de la productivité $\rightarrow$ emplois décents $\rightarrow$ transformation locale.
    \item \textbf{Création de valeur partagée (\emph{shared value} -- Porter \& Kramer) :} dépasser la RSE « philanthropique » ou de simple « conformité » : intégrer le progrès social \textbf{dans} le modèle d'affaires. Ex. : une PME du bois qui forme et emploie des jeunes, certifie sa production durable et vend des meubles finis (valeur ajoutée locale) au lieu d'exporter des grumes (rente). « Passer de la rente à la création de valeur partagée ».
\end{itemize}

\textbf{2. Thèse principale et argumentation :}
\begin{itemize}
    \item \textbf{Thèse :} \og le blocage du « chaînon manquant » (PME formelles) empêche l'émergence du Cameroun ; le redynamiser est la condition du passage d'une économie de rente à une économie de valeur partagée. \fg
    \item \textbf{Argumentation (causes $\rightarrow$ conséquences $\rightarrow$ rôle attendu) :}
    \begin{enumerate}
        \item \textbf{Causes du blocage (diagnostic) :} \textbf{financement :} coût du crédit élevé (souvent supérieur à 12\,\%), exclusion bancaire (garanties foncières exigées, asymétrie d'information). \textbf{Fiscalité :} perçue comme « prédatrice » (pression des contrôles, TVA/IS lourdes, absence d'incitations aux PME). \textbf{Infrastructures :} électricité instable (coût des groupes électrogènes -- cf. Bois \& Design). \textbf{Réglementation :} complexe, coûts de conformité disproportionnés (procédures, délais).
        \item \textbf{Conséquences :} maintien de la bimodalité. Les GE (rente, peu d'emplois qualifiés) coexistent avec l'informel (survie, sans protection). Pas de relais de croissance endogène ; jeunes diplômés au chômage.
        \item \textbf{Rôle attendu des PME (levier) :} \textbf{innovation incrémentale} (adaptation locale), \textbf{emploi décent des jeunes} (formation sur le tas, insertion), \textbf{transformation locale des matières premières} (bois $\rightarrow$ meubles, cacao $\rightarrow$ chocolat, coton $\rightarrow$ textile) = capture de la valeur ajoutée et amélioration de la balance commerciale.
    \end{enumerate}
\end{itemize}

\textbf{3. Discussion : pertinence pour le Cameroun (nuances et contre-arguments) :}
\begin{itemize}
    \item \textbf{Réponses publiques identifiées :}
    \begin{itemize}
        \item \textbf{Lois de finances récentes :} mesures en faveur des PME (exonérations fiscales pour les jeunes entreprises, guichet unique, plafonnement des contrôles). \emph{Nuance :} application effective ? Portée limitée si l'électricité et l'accès au foncier ne sont pas résolus.
        \item \textbf{PADI (Programme d'Appui au Développement de l'Entrepreneuriat) :} guichets de financement, formation, incubateurs. \emph{Nuance :} effet d'aubaine ? Montants insuffisants au regard du déficit de financement des PME ; cible souvent les « startups » technologiques plutôt que les PME industrielles du « chaînon manquant ».
        \item \textbf{ZLECAf (Zone de Libre-Échange Continentale Africaine) :} marché de plus d'un milliard de consommateurs ; opportunité pour les PME exportatrices (règles d'origine, normes). \emph{Nuance :} risque de concurrence de produits moins chers (Nigeria, Égypte, Afrique du Sud) si les PME ne sont pas compétitives (productivité, énergie, logistique). Nécessite une \textbf{politique industrielle} active (clusters, zones économiques spéciales type Kribi).
    \end{itemize}
    \item \textbf{Contre-arguments et nuances majeures :}
    \begin{enumerate}
        \item \textbf{Résilience de l'informel :} pas seulement une « survie » : économie de la débrouille, flexibilité, inclusion massive (environ 90\,\% de l'emploi), filet social. La transition vers le formel implique des coûts fixes (impôts, charges, locaux) souvent insupportables pour des micro-marges. Une \textbf{approche de transition progressive} (formalisation simplifiée, régime fiscal synthétique amélioré, protection sociale) est plus réaliste que la répression.
        \item \textbf{Rôle irremplaçable des GE :} les infrastructures lourdes (port de Kribi, barrages, télécoms, banques) exigent des capitaux massifs et des économies d'échelle : \textbf{seules les GE, les groupes étrangers ou l'État peuvent les porter}. Elles créent l'écosystème (sous-traitance aux PME) si le \textbf{contenu local} est imposé et suivi.
        \item \textbf{Gouvernance et institutions :} le texte omet le facteur \textbf{institutionnel} : corruption, insécurité juridique, instabilité. Sans État de droit effectif, aucune réforme technique (PADI, ZLECAf) ne débloquera durablement le chaînon manquant.
    \end{enumerate}
\end{itemize}
\end{corrige}

\begin{corrige}[Exercice 11 — Cas Pratique Intégré : « Bois \& Design » — Diagnostic Stratégique et Social]
\textbf{BARÈME INDICATIF (20 pts) :} Q1 Juridique (4) + Q2 Financière (6) + Q3 Organisationnelle (4) + Q4 Note de synthèse (6).

\textbf{1. Qualification juridique (Doc 1) :}
\begin{itemize}
    \item \textbf{Forme sociale :} \textbf{Société à Responsabilité Limitée (SARL)} (Acte Uniforme OHADA relatif au droit des sociétés commerciales et du GIE).
    \item \textbf{Responsabilité des associés :} \textbf{limitée au montant de leurs apports} : les associés ne sont tenus des dettes sociales qu'à concurrence de leurs apports. Le patrimoine personnel est protégé (sauf faute de gestion).
    \item \textbf{Organe de direction :} \textbf{gérant unique : M. FOUDA} (personne physique, nommé dans les statuts ou par acte séparé). Il est investi des pouvoirs les plus étendus pour agir au nom de la société. Pas de Conseil d'Administration (structure légère).
    \item \textbf{Régime fiscal :} \textbf{Impôt sur les Sociétés (IS) de droit commun}, au taux de 30\,\% majoré des centimes additionnels (soit 33\,\% en pratique) sur le bénéfice imposable, avec acompte provisionnel calculé sur le chiffre d'affaires. Une option pour l'IR est possible sous conditions strictes (non précisées ici) : l'IS s'applique donc par défaut.
\end{itemize}

\textbf{2. Analyse financière (Doc 3 = tableau de l'Exo 8)} \\
\textit{Note méthodologique : le tableau du Doc 3 donne VA = 195, EBE = 50, RN = 15 pour 2023, tandis que le détail comptable de l'Exo 4 donne VA = 200, EBE = 65, RN = 17. Pour la cohérence interne du dossier, on utilise les \textbf{données du Doc 3}, en signalant l'écart avec le détail comptable.}

\begin{itemize}
    \item \textbf{VA 2023 (Doc 3) :} \textbf{195 M FCFA} (donnée).
    \item \textbf{EBE 2023 (Doc 3) :} \textbf{50 M FCFA} (donnée).
    \item \textbf{Taux de marge sur VA 2023 :} $\frac{195}{500} = \textbf{39\,\%}$ (contre $180/450 = 40\,\%$ en 2022). \emph{Quasi stable.}
    \item \textbf{Productivité du travail (VA/Effectif) 2023 :} $\frac{195}{50} = \textbf{3,9 M FCFA/pers.}$ (contre 4,0 M en 2022). \emph{Baisse confirmée.}
    \item \textbf{Comparaison 2022 / 2023 (synthèse) :}
    \begin{center}
    \begin{tabular}{|l|c|c|c|}\hline
    \rowcolor{popblueL}
    \textbf{Indicateur} & \textbf{2022} & \textbf{2023} & \textbf{Variation} \\ \hline
    CA (M FCFA) & 450 & 500 & +11,1\,\% \\ \hline
    VA (M FCFA) & 180 & 195 & +8,3\,\% \\ \hline
    EBE (M FCFA) & 45 & 50 & +11,1\,\% \\ \hline
    RN (M FCFA) & 12 & 15 & +25,0\,\% \\ \hline
    Taux de marge (VA/CA) & 40,0\,\% & 39,0\,\% & $-1,0$ point \\ \hline
    Productivité (M FCFA/pers.) & 4,00 & 3,90 & $-2,5$\,\% \\ \hline
    Part des salaires (Masse/VA) & 55,6\,\% & 61,5\,\% & +5,9 points \\ \hline
    \end{tabular}
    \end{center}
    \textbf{Diagnostic de santé financière :}
    \begin{itemize}
        \item \textbf{Rentabilité :} en progression (RN +25\,\%), mais à un niveau faible (marge nette de 3\,\%). Taux de marge stable autour de 39\,\% (correct pour la transformation du bois).
        \item \textbf{Structure des charges :} \textbf{masse salariale dynamique (+20\,\%)}, bien au-dessus de la croissance de la VA (+8,3\,\%). Poids élevé des charges fixes (personnel + amortissements) : risque sur le seuil de rentabilité.
        \item \textbf{Autofinancement (CAF) :} CAF $\approx$ RN + Amortissements = 15 + 25 = 40 M (estimation). Les investissements nécessaires (maintenance, EPI, énergie) dépassent probablement la CAF $\rightarrow$ \textbf{trésorerie tendue} (confirmé par le Doc 4 : carburant, groupes électrogènes). Endettement fournisseurs ou bancaire probable.
        \item \textbf{Verdict :} \textbf{croissance extensive fragile.} Rentabilité en trompe-l'œil (tirée par le volume, pas par l'efficacité). Besoins urgents : \textbf{investissement productif} (énergie solaire, machines), \textbf{maîtrise de la masse salariale} (productivité), \textbf{négociation d'un crédit d'investissement} (guichets PME).
    \end{itemize}
\end{itemize}

\textbf{3. Analyse organisationnelle (Doc 2 = organigramme de l'Exo 6) :}
\begin{itemize}
    \item \textbf{Mode de structure :} \textbf{fonctionnelle (hiérarchique, par spécialités).}
    \item \textbf{Étendue de contrôle (Directeur de Production) :} \textbf{3} (3 Chefs d'atelier directs). Adaptée (norme usuelle de 5 à 8), elle permet la proximité.
    \item \textbf{Taux d'encadrement :} $\frac{5+8}{50} = \textbf{26\,\%}$ (élevé).
    \item \textbf{Risques de dysfonctionnement pour une PME de 50 personnes :}
    \begin{enumerate}
        \item \textbf{Effet silo / cloisonnement :} communication uniquement verticale. Le Commercial vend, la Production planifie, les Achats approvisionnent, sans coordination : délais, erreurs, stocks.
        \item \textbf{Lenteur décisionnelle :} 4 niveaux (DG $\rightarrow$ Directeurs $\rightarrow$ Chefs d'atelier $\rightarrow$ Ouvriers). Information filtrée et déformée ; réactivité faible face aux urgences (Doc 4 : grève, panne d'énergie).
        \item \textbf{Coût de structure :} 13 cadres pour 50 personnes = charge fixe lourde ; rigidité en cas de baisse d'activité.
        \item \textbf{Dépendance au gérant unique (M. Fouda) :} gouvernance centralisée, risque d'« homme-clé », absence de collégialité (pas de Conseil d'Administration), succession non préparée.
    \end{enumerate}
\end{itemize}

\textbf{4. Note de synthèse au Gérant (250 mots maximum)}

\textbf{Note de synthèse confidentielle -- À l'attention de M. FOUDA, Gérant -- Objet : résolution du conflit social et prévention RSE (mars 2024).}

\textbf{Monsieur le Gérant,}

\textbf{1. Qualification juridique :} le mouvement observé relève d'un \textbf{conflit collectif mixte} : \textbf{conflit de droit} (violations du Code du Travail : repos hebdomadaire, majorations d'heures supplémentaires, médecine du travail, EPI) et \textbf{conflit d'intérêt} (revendication d'augmentation de 15\,\%). La grève est \textbf{licite} (préavis déposé, revendications professionnelles) : l'employeur ne peut la sanctionner.

\textbf{2. Trois mesures correctives immédiates (droit du travail) :}
\begin{enumerate}
    \item \textbf{Régularisation de la paie :} paiement intégral des heures supplémentaires dominicales avec majoration légale sur la prochaine paie.
    \item \textbf{Santé et sécurité d'urgence :} mise à disposition immédiate d'EPI conformes (lunettes, casques, protections auditives) à l'atelier de sciage ; convocation du médecin du travail sous 15 jours ; affichage des consignes de sécurité.
    \item \textbf{Respect du temps de travail :} cessation du travail dominical non autorisé ; demande de dérogation auprès de l'Inspecteur du Travail si impératif de production ; planning équitable (repos compensateur + majoration).
\end{enumerate}

\textbf{3. Deux engagements RSE structurels (ISO 26000) pour prévenir la récidive :}
\begin{enumerate}
    \item \textbf{Comité paritaire Hygiène-Sécurité-Environnement opérationnel :} composition paritaire (3 délégués du personnel, 3 représentants de la Direction), réunions mensuelles obligatoires ; missions : enquêtes sur les accidents, inspections des ateliers, programme de formation sécurité (objectif : 100\,\% des salariés formés par an), suivi d'indicateurs ; budget dédié.
    \item \textbf{Accord d'entreprise « Dialogue social et partage de la valeur » (négociation 2024) :} institutionnalisation de \textbf{négociations annuelles obligatoires} (salaires, classifications, conditions de travail) ; clause d'\textbf{intéressement} indexée sur l'EBE et la productivité (alignement des intérêts du capital et du travail) ; audit RSE annuel externe publié aux parties prenantes.
\end{enumerate}

\textbf{Le coût de l'inaction (grève, contentieux, image, accès aux marchés) dépasse largement le coût de l'investissement humain et sécuritaire. La paix sociale est un actif productif.}

\textit{Fait à Douala, le 15 mars 2024. Le Conseiller Juridique \& RSE.}
\end{corrige}

\newpage

\coursec{Fiche bilan}

\begin{popfigure}
\begin{center}
\begin{tikzpicture}[mindmap, grow cyclic, every node/.style={concept, text=white, font=\small\bfseries},
root concept/.append style={concept color=popink, font=\large\bfseries, minimum size=3.2cm},
level 1/.append style={level distance=3.6cm, sibling angle=60, minimum size=2.4cm},
level 1 concept/.append style={font=\footnotesize\bfseries}]
\node[root concept, align=center]{LES\\ENTREPRISES}
child[concept color=popblue]{ node {Définition et nature} }
child[concept color=poporange]{ node {Typologie au Cameroun} }
child[concept color=popgreen]{ node {Facteurs de production} }
child[concept color=poppurple]{ node {Organisation interne} }
child[concept color=poppink]{ node {Gouvernance et stratégie} }
child[concept color=popgold]{ node {RSE} };
\end{tikzpicture}
\end{center}
\legende{Carte mentale de la leçon : les six idées maîtresses autour de la notion d'entreprise.}
\end{popfigure}

\begin{bilan}
\textbf{Définitions clés}
\begin{itemize}
\item \cle{Entreprise} : unité économique, juridiquement autonome, qui combine des facteurs de production (travail et capital) pour produire des biens ou des services destinés à être vendus sur un marché, en vue d'un profit ou d'un service d'intérêt général.
\item \cle{Facteurs de production} : le \cle{travail} (activité humaine, physique et intellectuelle) et le \cle{capital} (machines, bâtiments, outils, capitaux financiers), auxquels on ajoute souvent le savoir-faire de l'entrepreneur.
\item \cle{Valeur ajoutée} : richesse réellement créée par l'entreprise ; elle se calcule ainsi : $\text{Valeur ajoutée} = \text{Chiffre d'affaires} - \text{Consommations intermédiaires}$.
\item \cle{Gouvernance} : ensemble des règles et organes (direction, conseil, assemblée) qui décident, contrôlent et répartissent le pouvoir dans l'entreprise.
\item \cle{RSE} (Responsabilité Sociétale des Entreprises) : prise en compte volontaire, par l'entreprise, des impacts sociaux et environnementaux de ses activités (conditions de travail, pollution, vie des communautés locales).
\end{itemize}

\textbf{Typologie des entreprises (à connaître par c\oe ur)}
\begin{center}
\begin{tabularx}{\linewidth}{|l|X|X|}
\hline
\rowcolor{popblueL} \textbf{Critère} & \textbf{Catégories} & \textbf{Exemples camerounais} \\
\hline
Taille & TPE, PME, grandes entreprises & Atelier de quartier ; PME agroalimentaire ; SABC, ENEO \\
\hline
Secteur & Primaire, secondaire, tertiaire & Exploitation de cacao ; scierie ; banque, transport \\
\hline
Propriété du capital & Privée, publique, mixte & PME familiale ; CAMTEL ; société à capitaux mixtes \\
\hline
Forme juridique & Entreprise individuelle, SARL, SA & Commerçant de Mokolo ; SARL artisanale ; SA cotée ou non \\
\hline
\end{tabularx}
\end{center}

\textbf{Méthodes express}
\begin{enumerate}
\item Pour définir : donner la nature (unité économique autonome), la fonction (produire), le moyen (combiner travail et capital) et la finalité (marché, profit ou service public).
\item Pour analyser un cas concret : identifier la taille, le secteur, la forme juridique, puis les facteurs de production mobilisés, et conclure sur la gouvernance et la RSE.
\item Pour la dissertation : problématiser (l'entreprise ne poursuit-elle que le profit ?), opposer les thèses (finalité lucrative / responsabilité sociale), illustrer par des exemples camerounais, conclure par une réponse nuancée et justifiée.
\end{enumerate}
\end{bilan}

\begin{aretenir}[Checklist avant le Bac]
\begin{itemize}
\item[$\square$] Je sais définir une entreprise en une phrase complète (nature, fonction, moyens, finalité).
\item[$\square$] Je distingue les deux facteurs de production : travail et capital.
\item[$\square$] Je connais la formule de la valeur ajoutée et je sais l'appliquer à un cas simple.
\item[$\square$] Je cite les quatre critères de classification : taille, secteur, propriété du capital, forme juridique.
\item[$\square$] Je donne au moins deux exemples d'entreprises camerounaises pour chaque grande catégorie.
\item[$\square$] Je sais décrire l'organisation interne : direction, services (production, commercial, administratif), employés.
\item[$\square$] Je définis la gouvernance et j'identifie qui décide dans une PME et dans une grande entreprise.
\item[$\square$] Je définis la RSE et je cite deux exemples concrets (environnement, conditions de travail).
\item[$\square$] Je sais construire une problématique : l'entreprise a-t-elle d'autres devoirs que le profit ?
\item[$\square$] Je rédige une conclusion argumentée et justifiée, sans simple opinion personnelle.
\item[$\square$] Je sais mener le diagnostic simplifié d'une PME comme dans le cas « Bois \& Design ».
\end{itemize}
\end{aretenir}

\findecours

\end{document}
